\chapter{Fields}
\label{chap:fields}

The minuscule amount of algebra we have developed so far allows us to
scratch the surface of the unfathomable subject of \emph{field theory}.
Fields are of basic importance in many subjects---number theory and
algebraic geometry come immediately to mind---and it is not surprising that
their study has been developed to a particularly high level of
sophistication. While my profoundly limited knowledge will prevent me from
even hinting at this sophistication, even a cursory overview of the basic
notions will allow us to deal with remarkable applications, such as the
famous problems of constructibility of geometric figures. We will also get a
glimpse of the beautiful subject of \emph{Galois theory}, an amazing
interplay of the theory of field extensions, the solvability of polynomials,
and group theory.

\section{Field extensions, I}
\label{sec:fields:field-extensions-i}

We first deal with several basic notions in field theory, mostly inspired
by easy linear algebra considerations. The keywords here are \emph{finite},
\emph{simple}, \emph{finitely generated}, \emph{algebraic}.

\subsection{Basic definitions}
\label{subsec:fields:basic-definitions}
The study of fields is (of course) the study of the category \textsf{Fld}
of fields (Definition~III.1.14), with ring homomorphisms as morphisms. The
task is to understand what fields are and above all what they are \emph{in
relation to one another}.

The place to start is a reminder from elementary ring theory\footnote{The
last time I have made this point was back in \S{}V.5.2.}: \emph{every ring
homomorphism from a field to a nonzero ring is injective}. Indeed, the
kernel of a ring homomorphism to a nonzero ring is a proper ideal (because
a homomorphism maps $1$ to $1$ by definition and $1 \neq 0$ in nonzero
rings), and the only proper ideal in a field is $(0)$. In particular, every
ring homomorphism of fields is injective (fields are nonzero rings by
definition!); every morphism in \textsf{Fld} is a monomorphism (cf.\
Proposition~III.2.4).
Thus, every morphism $k \to K$ between two fields identifies the first with
a subfield of the second. In other words, one can view $K$ as a particular
way to enlarge $k$, that is, as an \emph{extension} of $k$. Field theory
is first and foremost the study of \emph{field extensions}. I will denote
a field extension by $k \subseteq K$ (which is less than optimal, because
there may be many ways to embed a field into another); other popular choices
are $K/k$ (which I do not like, as it suggests a quotienting operation) and
$\frac{K}{k}$ (which we will avoid most of the time, as it is hard to
typeset).
The coarsest invariant of a field $k$ is its \emph{characteristic}; cf.\
Exercise~V.4.17. We have a unique ring homomorphism $i : \mathbb{Z} \to k$
($\mathbb{Z}$ is initial in \textsf{Ring}: $1$ must go to $1$ by definition
of homomorphism, and this fixes the value of $i(n)$ for all $n \in
\mathbb{Z}$); the \emph{characteristic} of $k$, $\operatorname{char} k$, is
defined to be the nonnegative generator of the ideal $\ker i$; that is,
$\operatorname{char} k = 0$ if $i$ is injective, and $\operatorname{char} k
= p > 0$ if $\ker i = (p) \neq (0)$.

Put otherwise, either $n \cdot 1$ is only $0$ in $k$ for $n = 0$, in which
case $\operatorname{char} k = 0$, or $n \cdot 1 = 0$ for some $n \neq 0$;
$\operatorname{char} k = p > 0$ is then the smallest positive integer for
which $p \cdot 1 = 0$ in $k$. Since fields are integral domains, the image
$i(\mathbb{Z})$ must be an integral domain; hence $\ker i$ must be a prime
ideal. Therefore, the characteristic of a field is either $0$ or a prime
number.
If $k \subseteq K$ is an extension, then $\operatorname{char} k =
\operatorname{char} K$ (Exercise~1.1). Thus, we could define and study
categories $\textsf{Fld}_0$, $\textsf{Fld}_p$ of fields of a given
characteristic, without losing any information: these categories live within
\textsf{Fld}, without interacting with each other.\footnote{They do interact
in other ways, however---for example, because $\mathbb{Z}$ is initial in
\textsf{Ring}, so $\mathbb{Z}$ maps to all fields.} Further, each of these
categories has an initial object: $\mathbb{Q}$ is initial in
$\textsf{Fld}_0$, and\footnote{I am going to denote the field
$\mathbb{Z}/p\mathbb{Z}$ by $\mathbb{F}_p$ in this chapter, to stress its
role as a field (rather than `just' as a group).} $\mathbb{F}_p :=
\mathbb{Z}/p\mathbb{Z}$ in $\textsf{Fld}_p$; the reader has checked this
easy fact in Exercise~V.4.17, and now is an excellent time to contemplate it
again. In each case, the initial object is called the \emph{prime subfield}
of the given field. Thus, every field is in a canonical way an extension of
its prime subfield: studying fields really means studying field extensions.
To a large extent, the `small' field $k$ will be fixed in what follows, and
our main object of study will be the category $\textsf{Fld}_k$ of extensions
of $k$,
with the evident (and by now familiar to the reader) notions of morphisms.

The first general remark concerning field extensions is that the larger
field is an algebra, and hence a \emph{vector space} over the smaller one
(by the very definition of algebra; cf.\ Example~III.5.6).

\begin{defn}{}{fields:1.1}
A field extension $k \subseteq F$ is \emph{finite}, of \emph{degree} $n$,
if $F$ has (finite) dimension $\dim F = n$ as a vector space over $k$. The
extension is \emph{infinite} otherwise.

The degree of a finite extension $k \subseteq F$ is denoted by $[F : k]$
(and we write $[F : k] = \infty$ if the extension is infinite).
\end{defn}
In \S{}V.5.2 we have encountered a prototypical example of finite field
extension: a procedure starting from an irreducible polynomial $f(x) \in
k[x]$ with coefficients in a field and producing an extension $K$ of $k$ in
which $f(t)$ has a root. Explicitly,
\[K = \frac{k[t]}{(f(t))}\]
is such an extension (Proposition~V.5.7). The quotient is a field because
$k[t]$ is a PID; hence irreducible elements generate maximal ideals---prime
because of the UFD property (cf.\ Theorem~V.2.5) and then maximal because
nonzero prime ideals are maximal in a PID (cf.\ Proposition~III.4.13). The
coset of $t$ is a root of $f(x)$ when this is viewed as an element of
$K[x]$. The degree of the extension $k \subseteq K$ equals the degree of
the polynomial $f(x)$ (this is hopefully clear at this point and was checked
carefully in Proposition~III.4.6).
The reader may wonder whether perhaps \emph{all} finite extensions are of
this type. This is unfortunately not the case, but as it happens, it is the
case in a large class of examples. As I often do, I promise that the
situation will clarify itself considerably once we have accumulated more
general knowledge; in this case, the relevant fact will be
Proposition~5.19.

Also recall that we proved that these extensions are almost universal with
respect to the problem of extending $k$ so that the polynomial $f(t)$
acquires a root. I pointed out, however, that the \emph{uni} in \emph{uni}versal
is missing; that is, if $k \subseteq F$ is an extension of $k$ in which
$f(t)$ has a root, there may be many different ways to put $K$ in between:
\[k \subseteq K \subseteq F.\]
Such questions are central to the study of extensions, so we begin by
giving a second look at this situation.

\subsection{Simple extensions}
\label{subsec:fields:simple-extensions}
Let $k \subseteq F$ be a field extension, and let $\alpha \in F$. The
smallest subfield of $F$ containing both $k$ and $\alpha$ is denoted
$k(\alpha)$; that is, $k(\alpha)$ is the intersection of all subfields of
$F$ containing $k$ and $\alpha$.

\begin{defn}{}{fields:1.2}
A field extension $k \subseteq F$ is \emph{simple} if there exists an
element $\alpha \in F$ such that $F = k(\alpha)$.
\end{defn}
The extensions recalled above are of this kind: if $K = k[t]/(f(t))$ and
$\alpha$ denotes the coset of $t$, then $K = k(\alpha)$: indeed, if a
subfield of $K$ contains the coset of $t$, then it must contain (the coset
of) every polynomial expression in $t$, and hence it must be the whole of
$K$.

I have always found the notation $k(\alpha)$ somewhat unfortunate, since it
suggests that all such extensions are in some way isomorphic and possibly
all isomorphic to the field $k(t)$ of rational functions in one
indeterminate $t$ (cf.\ Definition~V.4.13). This is not true, although it
is clear that every element of $k(\alpha)$ may be written as
a rational function in $\alpha$ with coefficients in $k$ (Exercise~1.3).
In any case, it is easy to classify simple extensions: they are either
isomorphic to $k(t)$ or they are of the prototypical kind recalled above.
Here is the precise statement.

\begin{prop}{}{fields:1.3}
Let $k \subseteq k(\alpha)$ be a simple extension. Consider the evaluation
map $\epsilon : k[t] \to k(\alpha)$, defined by $f(t) \mapsto f(\alpha)$.
Then we have the following:
\begin{itemize}
\item $\epsilon$ is injective if and only if $k \subseteq k(\alpha)$ is an
  infinite extension. In this case, $k(\alpha)$ is isomorphic to the field
  of rational functions $k(t)$.
\item $\epsilon$ is not injective if and only if $k \subseteq k(\alpha)$ is
  finite. In this case there exists a unique monic irreducible nonconstant
  polynomial $p(t) \in k[t]$ of degree $n = [k(\alpha) : k]$ such that
\[k(\alpha) \cong \frac{k[t]}{(p(t))}.\]
  Via this isomorphism, $\alpha$ corresponds to the coset of $t$. The
  polynomial $p(t)$ is the monic polynomial of smallest degree in $k[t]$
  such that $p(\alpha) = 0$ in $k(\alpha)$.
\end{itemize}
\end{prop}
The polynomial $p(t)$ appearing in this statement is called the
\emph{minimal polynomial} of $\alpha$ over $k$.

Of course the minimal polynomial of an element $\alpha$ of a (`large')
field depends on the base (`small') field $k$. For example, $\sqrt{2} \in
\mathbb{C}$ has minimal polynomial $t^2 - 2$ over $\mathbb{Q}$, but $t -
\sqrt{2}$ over $\mathbb{R}$.

\begin{proof}
Let $F = k(\alpha)$. By the `first isomorphism theorem', the image of
$\epsilon : k[t] \to F$ is isomorphic to $k[t]/\ker(\epsilon)$. Since $F$
is an integral domain, so is $k[t]/\ker(\epsilon)$; hence $\ker(\epsilon)$
is a prime ideal in $k[t]$.
---Assume $\ker(\epsilon) = 0$; that is, $\epsilon$ is an injective map
from the integral domain $k[t]$ to the field $F$. By the universal property
of fields of fractions (cf.\ \S{}V.4.2), $\epsilon$ extends to a unique
homomorphism
\[k(t) \to F.\]
The (isomorphic) image of $k(t)$ in $F$ is a field containing $k$ and
$\alpha$; hence it equals $F$ by definition of simple extension.

Since $\epsilon$ is injective, the powers $\alpha^0 = 1, \alpha, \alpha^2,
\alpha^3, \dots$ (that is, the images $\epsilon(t^i)$) are all distinct and
linearly independent over $k$ (because the powers $1, t, t^2, \dots$ are
linearly independent over $k$); therefore the extension $k \subseteq F$ is
infinite in this case.
---If $\ker(\epsilon) \neq 0$, then $\ker(\epsilon) = (p(t))$ for a unique
monic irreducible nonconstant polynomial $p(t)$, which has smallest degree
(cf.\ Exercise~III.4.4!) among all nonzero polynomials in $\ker(\epsilon)$.
As $(p(t))$ is then maximal in $k[t]$, the image of $\epsilon$ is a
subfield of $F$ containing $\alpha = \epsilon(t)$. By definition of simple
extension, $F = $ the image of $\epsilon$; that is, the induced
homomorphism
\[\frac{k[t]}{(p(t))} \to F\]
is an isomorphism. In this case $[F : k] = \deg p(t)$, as recalled in
\S{}\ref{subsec:fields:basic-definitions}, and in particular the extension
is finite, as claimed.
\end{proof}

The alert reader will have noticed that the proof of \cref{prop:fields:1.3}
is essentially a rehash of the argument proving the `versality' part of
Proposition~V.5.7.
\begin{exam}{}{fields:exam:1.4}
Consider the extension $\mathbb{Q} \subseteq \mathbb{R}$.

The polynomial $x^2 - 2 \in \mathbb{Q}[x]$ has roots in $\mathbb{R}$:
therefore, by Proposition~V.5.7 there exists a homomorphism (hence a field
extension)
\[\epsilon : \frac{\mathbb{Q}[t]}{(t^2 - 2)} \hookrightarrow \mathbb{R},\]
such that the image of (the coset of) $t$ is a root $\alpha$ of $x^2 - 2$.
\cref{prop:fields:1.3} simply identifies the image of this homomorphism
with $\mathbb{Q}(\alpha) \subseteq \mathbb{R}$.
This is hopefully crystal clear; however, note that even this simple
example shows that the induced morphism $\epsilon$ is not unique (hence the
`lack of \emph{uni}'): because there are more than one root of $x^2 - 2$
in $\mathbb{R}$. Concretely, there are two possible choices for $\alpha$:
$\alpha = +\sqrt{2}$ and $\alpha = -\sqrt{2}$. The choice of $\alpha$
determines the evaluation map $\epsilon$ and therefore the specific
realization of $\mathbb{Q}[t]/(t^2 - 2)$ as a subfield of $\mathbb{R}$.

The reader may be misled by one feature of this example: clearly
$\mathbb{Q}(\sqrt{2}) = \mathbb{Q}(-\sqrt{2})$, and this may seem to
compensate for the lack of uniqueness lamented above. The morphism may not
be unique, but the image of the morphism surely is? No. The reader will
check that there are three distinct subfields of $\mathbb{C}$ isomorphic to
$\mathbb{Q}[t]/(t^3 - 2)$ (Exercise~1.5).
\end{exam}

Ay, there's the rub. One of our main goals in this chapter will be to
single out a condition on an extension $k \subseteq F$ that guarantees that
no matter how we
embed $F$ in a larger extension, the images of these (possibly many
different) embeddings will all coincide. Up to further technicalities, this
is what makes $\mathbb{Q} \subseteq \mathbb{Q}(\sqrt{2})$ a Galois
extension, while $\mathbb{Q} \subseteq \mathbb{Q}(\sqrt[3]{2})$ will not be
a Galois extension. But Galois extensions will have to wait until \S{}6,
and the reader can put this issue aside for now.

One way to recover uniqueness is to incorporate the choice of the root in
the data. This leads to the following refinement of the (uni)versality
statement, adapted for future applications.
\begin{prop}{}{fields:1.5}
Let $k_1 \subseteq F_1 = k_1(\alpha_1)$, $k_2 \subseteq F_2 = k_2(\alpha_2)$
be two finite simple extensions. Let $p_1(t) \in k_1[t]$, resp., $p_2(t)
\in k_2[t]$, be the minimal polynomials of $\alpha_1$, resp., $\alpha_2$.
Let $i : k_1 \to k_2$ be an isomorphism, such that\footnote{Of course any
homomorphism of rings $f : R \to S$ induces a unique ring homomorphism $f :
R[t] \to S[t]$ sending $t$ to $t$, named in the same way by a harmless
abuse of language.}
\[i(p_1(t)) = p_2(t).\]
Then there exists a unique isomorphism $j : F_1 \to F_2$ agreeing with $i$
on $k_1$ and such that $j(\alpha_1) = \alpha_2$.
\end{prop}
\begin{proof}
Since every element of $k_1(\alpha_1)$ is a linear combination of powers of
$\alpha_1$ with coefficients in $k_1$, $j$ is determined by its action on
$k_1$ (which agrees with $i$) and by $j(\alpha_1)$, which is prescribed to
be $\alpha_2$. Thus an isomorphism $j$ as in the statement is uniquely
determined.

To see that the isomorphism $j$ exists, note that as $i$ maps $p_1(t)$ to
$p_2(t)$, it induces an isomorphism
\[\frac{k_1[t]}{(p_1(t))} \xrightarrow{\;\sim\;} \frac{k_2[t]}{(p_2(t))}.\]
Composing with the isomorphisms found in \cref{prop:fields:1.3} gives $j$:
\[j : k_1(\alpha_1) \xrightarrow{\;\sim\;} \frac{k_1[t]}{(p_1(t))}
  \xrightarrow{\;\sim\;} \frac{k_2[t]}{(p_2(t))} \xrightarrow{\;\sim\;}
  k_2(\alpha_2),\]
as needed.
\end{proof}
Thus, isomorphisms lift uniquely to simple extensions, so long as they
preserve minimal polynomials. We may say that $j$ \emph{extends} $i$, and
draw the following diagram of extensions:
\[\begin{tikzcd}
  k_1(\alpha_1) \arrow[r, "j", "\sim"'] & k_2(\alpha_2) \\
  k_1 \arrow[u, hook] \arrow[r, "i", "\sim"'] & k_2 \arrow[u, hook]
\end{tikzcd}\]
We will be particularly interested in considering isomorphisms of a field
to itself, subject to the condition of fixing a specified subfield $k$
(that is, extending the identity on $k$), that is, automorphisms in
$\textsf{Fld}_k$. Explicitly, for $k \subseteq F$ an extension, we will
analyze isomorphisms $j : F \xrightarrow{\sim} F$ such that $\forall c \in
k$, $j(c) = c$. I will denote the group of such automorphisms by
$\operatorname{Aut}_k(F)$. This is probably the most important object of
study in this chapter, so I should make its definition
official:

\begin{defn}{}{fields:1.6}
Let $k \subseteq F$ be a field extension. The \emph{group of automorphisms
of the extension}, denoted $\operatorname{Aut}_k(F)$, is the group of field
automorphisms $j : F \to F$ such that $j|_k = \mathrm{id}_k$.
\end{defn}

\begin{coro}{}{fields:1.7}
Let $k \subseteq F = k(\alpha)$ be a simple finite extension, and let
$p(x)$ be the minimal polynomial of $\alpha$ over $k$. Then
$|\operatorname{Aut}_k(F)|$ equals the number of distinct roots of $p(x)$
in $F$; in particular,
\[|\operatorname{Aut}_k(F)| \leq [F : k],\]
with equality if and only if $p(x)$ factors over $F$ as a product of
distinct linear polynomials.
\end{coro}

\begin{proof}
Let $j \in \operatorname{Aut}_k(F)$. Since every element of $F$ is a
polynomial expression in $\alpha$ with coefficients in $k$, and $j$ extends
the identity on $k$, $j$ is determined by $j(\alpha)$. Now
\[p(j(\alpha)) = j(p(\alpha)) = j(0) = 0:\]
therefore, $j(\alpha)$ is necessarily a root of $p(x)$. This shows that
$|\operatorname{Aut}_k(F)|$ is no larger than the number of roots of $p(x)$.

On the other hand, by \cref{prop:fields:1.5} every choice of a root of
$p(t)$ determines a lift of the identity to an element $j \in
\operatorname{Aut}_k(F)$, establishing the other inequality.
\end{proof}
\cref{coro:fields:1.7} is our first hint of the powerful interaction
between group theory and the theory of field extensions; this theme will
haunt us throughout the chapter.

The proof of \cref{coro:fields:1.7} really sets up a bijection between the
roots of an irreducible polynomial $p(t) \in k[t]$ in a larger field $F$
and the group $\operatorname{Aut}_k(F)$. A particularly optimistic reader
may then hope that the roots of $p(t)$ come endowed with a group structure,
but this is hoping for too much: the bijection depends on the choice of one
root $\alpha$, and no one root of $p(t)$ is `more beautiful' than the
others.

However, we have established that if $F = k(\alpha)$ is a simple extension
of $k$, then the group $\operatorname{Aut}_k(F)$ acts faithfully and
transitively on the set of roots of $p(t)$ in $F$. In other words,
$\operatorname{Aut}_k(F)$ can be identified with a certain subgroup of the
symmetric group acting on the set of roots of $p(t)$.

More generally (Exercise~1.6) the automorphisms of any extension act on
roots of polynomials. One conclusion we can draw from these preliminary
considerations is that the analysis of $\operatorname{Aut}_k(F)$ is
substantially simplified if $k \subseteq F$ is a simple extension $k(\alpha)$
and further if the minimal polynomial $p(x)$ of $\alpha$ factors into
$\deg p(x)$ distinct linear terms in $F$. It will not come as a surprise
that we will focus our attention on such extensions in later sections: an
extension is \emph{Galois} precisely if it is of this type. (The reader is
invited to remember this fact, but its import will not be fully appreciated
until we develop substantially more material.)

\subsection{Finite and algebraic extensions}
\label{subsec:fields:finite-and-algebraic-extensions}
The dichotomy encountered in \cref{prop:fields:1.3} provides us with one
of the most important definitions in field theory:

\begin{defn}{}{fields:1.8}
Let $k \subseteq F$ be a field extension, and let $\alpha \in F$. Then
$\alpha$ is \emph{algebraic} over $k$, of \emph{degree} $n$ if $n =
[k(\alpha) : k]$ is finite; $\alpha$ is \emph{transcendental} over $k$
otherwise.

The extension $k \subseteq F$ is \emph{algebraic} if every $\alpha \in F$
is algebraic over $k$.
\end{defn}
By \cref{prop:fields:1.3}, $\alpha \in F$ is algebraic over $k$ if and only
if there exists a nonzero polynomial $f(x) \in k[x]$ such that $f(\alpha) =
0$. The minimal polynomial of $\alpha$ is the monic polynomial of smallest
degree satisfying this condition; as we have seen, it is necessarily
irreducible.
Also note that if $\alpha$ is algebraic over $k$, then every element of
$k(\alpha)$ may in fact be written as a \emph{polynomial} with coefficients
in $k$.

\emph{Finite extensions are necessarily algebraic:}
\begin{lem}{}{fields:1.9}
Let $k \subseteq F$ be a finite extension. Then every $\alpha \in F$ is
algebraic over $k$, of degree $\leq [F : k]$.
\end{lem}
\begin{proof}
Since $k \subseteq k(\alpha) \subseteq F$, the dimension of $k(\alpha)$ as a
$k$-vector space is bounded by $\dim_k F = [F : k]$.
\end{proof}
Concretely, if $k \subseteq F$ is finite and $\alpha \in F$, then the powers
$1, \alpha, \alpha^2 \dots$ are necessarily linearly dependent; and any
nontrivial linear dependence relation among them provides us with a nonzero
polynomial $f(x) \in k[x]$ such that $f(\alpha) = 0$.
The literal converse of the claim that `finite extensions are algebraic' is
not true; we will see an example in a moment. However, something along these
lines holds; understanding the situation requires a more careful look at
finiteness conditions.
First of all, compositions of finite extensions are finite extensions, and
the degree behaves nicely with respect to this operation:
\begin{prop}{}{fields:1.10}
Let $k \subseteq E \subseteq F$ be field extensions. Then $k \subseteq F$ is
finite if and only if both $k \subseteq E$ and $E \subseteq F$ are finite.
In this case,
\[
  [F : k] = [F : E][E : k].
\]
\end{prop}
\begin{proof}
If $F$ is finite-dimensional as a vector space over $k$, then so is its
subspace $E$; and any linear dependence relation of elements of $F$ over the
field $k$ gives one over the larger field $E$. It follows that if $k
\subseteq F$ is finite, then so are $k \subseteq E$ and $E \subseteq F$.
Conversely, assume $k \subseteq E$ and $E \subseteq F$ are both finite. Let
$(e_1, \dots, e_m)$ be a basis for $E$ over $k$, and let $(f_1, \dots, f_n)$
be a basis for $F$ over $E$. It will suffice to show that the $mn$ products
\[
  (e_1 f_1, e_1 f_2, \dots, e_m f_n)
\]
form a basis for $F$ over $k$.
Let $g \in F$. Then $\exists d_1, \dots, d_n \in E$ such that
\[
  g = \sum_{j=1}^{n} d_j f_j,
\]
since the elements $f_j$ span $F$ over $E$. Since $E$ is spanned by the
elements $e_i$ over $k$, $\exists c_{1j}, \dots, c_{mj} \in k$ such that
$\forall j$
\[
  d_j = \sum_{i=1}^{m} c_{ij} e_i.
\]
It follows that
\[
  g = \sum_{i=1}^{m} \sum_{j=1}^{n} c_{ij} e_i f_j:
\]
therefore the products $e_i f_j$ span $F$ over $k$. (This suffices to prove
that $k \subseteq F$ is finite.)
To verify that these elements are linearly independent, assume
\[
  \sum_{i,j} \lambda_{ij} e_i f_j = 0
\]
for $\lambda_{ij} \in k$. Then we have
\[
  \sum_{j} \Bigl(\sum_{i} \lambda_{ij} e_i\Bigr) f_j = 0,
\]
implying $\sum_i \lambda_{ij} e_i = 0$ for all $j$, since the elements $f_j$
are linearly independent over $E$. As the elements $e_i$ are linearly
independent over $k$, this shows that all $\lambda_{ij}$ equal $0$, as
needed.
\end{proof}
The formula given in the statement should look familiar to the reader: glance
at the end of \S{}II.8.5. Of course this is not accidental; the reader
should take it as another hint that the world of groups and the world of
fields have deep interaction. As in the case of groups, we draw the
immediate (but powerful) consequence, reminiscent of Lagrange's theorem:
\begin{coro}{}{fields:1.11}
Let $k \subseteq F$ be a finite extension, and let $E$ be an intermediate
field (that is, $k \subseteq E \subseteq F$). Then both $[E : k]$ and $[F :
E]$ divide $[F : k]$.
\end{coro}
As with Lagrange's theorem, judicious use of this result will make otherwise
mysterious statements melt into trivialities.
\begin{exam}{}{fields:exam:1.12}
Let $k \subseteq F$ be a field extension, and let $\alpha \in F$ be an
algebraic element over $k$, of odd degree. Then I claim that $\alpha$ may be
written as a polynomial in $\alpha^2$, with coefficients in $k$.
Indeed, $k(\alpha^2)$ is intermediate between $k$ and $k(\alpha)$:
\[
  k \subseteq k(\alpha^2) \subseteq k(\alpha);
\]
what can we say about the degree $d$ of $k(\alpha)$ over $k(\alpha^2)$?
Since $\alpha$ satisfies the polynomial $t^2 - \alpha^2 \in k(\alpha^2)[t]$,
$d \leq 2$. On the other hand, $d$ divides $[k(\alpha) : k]$ by
\cref{coro:fields:1.11}, and $[k(\alpha) : k]$ is odd, so $d \neq 2$.
Therefore $d = 1$, proving $k(\alpha) = k(\alpha^2)$, and in particular
$\alpha \in k(\alpha^2)$, which is the claim.
\end{exam}
Here is something else that should evoke fond memories for the reader: back
in \S{}III.6.5 we had encountered an important distinction between
\emph{finite} algebras and algebras \emph{of finite type}. An algebra is
finite over the base ring if it is finitely generated as a module, that is,
if it admits an onto homomorphism (of modules) from a finitely generated free
module; a commutative algebra is of finite type if it admits an onto
homomorphism (of algebras) from a polynomial ring in finitely many
variables.
Something along the same lines is going to occur here. An extension $k
\subseteq F$ is finite if and only if $\dim_k F$ is finite, that is, if and
only if $F$ is a finite $k$-algebra. The other finiteness condition takes the
following form.
\begin{defn}{}{fields:1.13}
A field extension $k \subseteq F$ is \emph{finitely generated} if there exist
$\alpha_1, \dots, \alpha_n \in F$ such that
\[
  F = k(\alpha_1)(\alpha_2)\dots(\alpha_n).
\]
\end{defn}
\begin{rem}{}{fields:1.14}
This is not the same as saying that $F$ is generated by the $\alpha_i$'s as
a (finite-type) $k$-algebra: even in the case of simple extensions, if
$\alpha$ is transcendental, then the natural evaluation map $\epsilon : k[t]
\to k(\alpha)$ of \cref{prop:fields:1.3} is not onto. However, this is an
epimorphism of rings in the categorical sense, as the reader will check
(Exercise~1.17); thus, the `categorical spirit' behind the finite-type
condition is preserved in this context.
This subtlety raises an interesting question: what if a field extension $k
\subseteq F$ really is a finite-type $k$-algebra? This turns out to be an
important issue, and we will come back to it in \S{}2.2.
\end{rem}
We will use the notation
\[
  k(\alpha_1, \alpha_2, \dots, \alpha_n)
\]
for the field $k(\alpha_1)(\alpha_2)\dots(\alpha_n)$. For $k \subseteq F$
and $\alpha_1, \dots, \alpha_n \in F$, the elements of the field $F =
k(\alpha_1, \dots, \alpha_n)$ are all the elements of $F$ which may be
written as rational functions in the $\alpha_i$'s, with coefficients in $k$
(cf.\ Exercise~1.3). Put otherwise, $k(\alpha_1, \dots, \alpha_n)$ is the
smallest subfield of $F$ containing $k(\alpha_1), \dots, k(\alpha_n)$ (it is
the \emph{composite} of these subfields). It is clear that the order of the
elements $\alpha_i$ is irrelevant.
Coming back to the issue of finite vs.\ algebraic, these two notions do
coincide for finitely generated extensions.
\begin{prop}{}{fields:1.15}
Let $k \subseteq F = k(\alpha_1, \dots, \alpha_n)$ be a finitely generated
field extension. Then the following are equivalent:
\begin{enumerate}[label=(\roman*)]
  \item $k \subseteq F$ is a finite extension.
  \item $k \subseteq F$ is an algebraic extension.
  \item Each $\alpha_i$ is algebraic over $k$.
\end{enumerate}
If these conditions are satisfied, then $[F : k] \leq$ the product of the
degrees of $\alpha_i$ over $k$.
\end{prop}
\begin{proof}
\cref{lem:fields:1.9} shows that (i)$\implies$(ii); (ii)$\implies$(iii)
trivially. Thus, we only need to prove that (iii)$\implies$(i), and to bound
the degree of $F$ over $k$ in the process.
Assume that each $\alpha_i$ is algebraic over $k$, and let $d_i$ be the
degree of $\alpha_i$ over $k$. By definition, $k \subseteq k(\alpha_i)$ is
finite, of degree $d_i$. By Exercise~1.16, each extension $k(\alpha_1, \dots,
\alpha_{i-1}) \subseteq k(\alpha_1, \dots, \alpha_i)$ is finite, of degree
$\leq d_i$. Applying \cref{prop:fields:1.10} to the composition of
extensions
\[
  k \subseteq k(\alpha_1) \subseteq k(\alpha_1, \alpha_2) \subseteq \dots
  \subseteq k(\alpha_1, \dots, \alpha_n) = F
\]
proves that $k \subseteq F$ is finite and $[F : k] \leq d_1 \dots d_n$, as
needed.
\end{proof}
While rather straightforward, \cref{prop:fields:1.15} has always seemed
remarkable to us: it says (in particular) that if $\alpha$ and $\beta$ are
algebraic over a field $k$, then so are $\alpha \pm \beta$, $\alpha\beta$,
$\alpha\beta^{-1}$, and any other rational function of $\alpha$ and $\beta$.
If all we knew were the simple-minded definition of `algebraic' (that is,
`root of a polynomial'), this would seem rather mysterious: given a
polynomial $f(x)$ of which $\alpha$ is a root and a polynomial $g(x)$ of
which $\beta$ is a root, how do we construct a polynomial $h(x)$ such that
(for example) $h(\alpha + \beta) = 0$? The answer is that we do not need to
perform any such construction, by virtue of \cref{prop:fields:1.15}.
One immediate consequence of this observation is that the set of algebraic
elements of any extension forms a field:

\begin{coro}{}{fields:1.16}
Let $k \subseteq F$ be a field extension. Let
\[
  E = \{\alpha \in F \mid \alpha \text{ is algebraic over } k\}.
\]
Then $E$ is a field.
\end{coro}
\begin{exam}{}{fields:exam:1.17}
Let $\overline{\Q} \subseteq \C$ be the set of complex numbers that are
algebraic over $\Q$; then $\overline{\Q}$ is a field, by
\cref{coro:fields:1.16}, and the extension $\Q \subseteq \overline{\Q}$ is
(tautologically) algebraic. Note that $\Q \subseteq \overline{\Q}$ is not a
finite extension, because in it there are elements of arbitrarily high degree
over $\Q$: indeed, there exist irreducible polynomials in $\Q[x]$ of
arbitrarily high degree, as we observed in Corollary~V.5.16.
Elements of $\overline{\Q}$ are called \emph{algebraic numbers}; the complex
numbers in the complement of $\overline{\Q}$ are transcendental over $\Q$ (by
definition); they are simply called \emph{transcendental numbers}. Elementary
cardinality considerations show that $\overline{\Q}$ is countable; thus, the
set of transcendental numbers is uncountable. Surprisingly, it may be
substantially difficult to prove that a given number is transcendental: for
example, the fact that $e$ is transcendental was only proved around 1870, by
\name{Charles Hermite}. The number $\pi$ is transcendental (\name{Lindemann},
ca.\ 1880); $e^\pi$ is transcendental (\name{Gelfond}--\name{Schneider},
1934). No one knows whether $\pi^e$, $\pi + e$, or $\pi e$ are
transcendental.
\end{exam}
Another consequence of \cref{prop:fields:1.15} is the important fact that
compositions of algebraic extensions are algebraic (whether finitely
generated or not):
\begin{coro}{}{fields:1.18}
Let $k \subseteq E \subseteq F$ be field extensions. Then $k \subseteq F$
is algebraic if and only if both $k \subseteq E$ and $E \subseteq F$ are
algebraic.
\end{coro}
\begin{proof}
If $k \subseteq F$ is algebraic, then every element of $F$ is algebraic
over $k$, hence over $E$, and every element of $E$ is algebraic over $k$;
thus $E \subseteq F$ and $k \subseteq E$ are algebraic.

Conversely, assume $k \subseteq E$ and $E \subseteq F$ are both algebraic,
and let $\alpha \in F$. Since $\alpha$ is algebraic over $E$, there exists
a polynomial
\[
  f(x) = x^n + e_{n-1}x^{n-1} + \dots + e_0 \in E[x]
\]
such that $f(\alpha) = 0$. This implies that in fact $\alpha$ is `already'
algebraic over the subfield $k(e_0, \dots, e_{n-1}) \subseteq E$;
therefore,
\[
  k(e_0, \dots, e_{n-1}) \subseteq k(e_0, \dots, e_{n-1}, \alpha)
\]
is a finite extension. On the other hand,
\[
  k \subseteq k(e_0, \dots, e_{n-1})
\]
is a finite extension by \cref{prop:fields:1.15}, since each $e_i$ is in
$E$ and therefore algebraic over $k$. By \cref{prop:fields:1.10},
\[
  k \subseteq k(e_0, \dots, e_{n-1}, \alpha)
\]
is finite. This implies that $\alpha$ is algebraic over $k$, as needed,
by \cref{lem:fields:1.9}.
\end{proof}
We will essentially deal only with finitely generated extensions, and
\cref{prop:fields:1.15} will simplify our work considerably. Also, since
finitely generated extensions are compositions of simple extensions, the
reader should expect that we will give a careful look at automorphisms of
such extensions.

It may in fact come as a surprise that, in many cases, finitely generated
extensions turn out to be simple to begin with; we will prove a precise
statement to this effect in due time (Proposition~5.19). The reader already
has enough tools to contemplate easy (but interesting) examples, such as
the following. This should serve as an encouragement to look at many more.
\begin{exam}{}{fields:exam:1.19}
Consider the extension $\Q \subseteq \Q(\sqrt{2}, \sqrt{3})$.

By \cref{prop:fields:1.15} we know that this is a finite (hence algebraic)
extension, of degree at most $4$.

Thus any five elements in $\Q(\sqrt{2}, \sqrt{3})$ must be linearly
dependent over $\Q$. We consider powers of $\sqrt{2} + \sqrt{3}$:
\[
  1,\quad (\sqrt{2}+\sqrt{3}),\quad (\sqrt{2}+\sqrt{3})^2,\quad
  (\sqrt{2}+\sqrt{3})^3,\quad (\sqrt{2}+\sqrt{3})^4
\]
must be linearly dependent. Thus, there must be rational numbers
$q_0, \dots, q_3$ such that
\[
  (\sqrt{2}+\sqrt{3})^4 + q_3(\sqrt{2}+\sqrt{3})^3
  + q_2(\sqrt{2}+\sqrt{3})^2 + q_1(\sqrt{2}+\sqrt{3}) + q_0 = 0.
\]

Elementary arithmetic shows that this relation is satisfied for $q_0 = 1$,
$q_1 = 0$, $q_2 = -10$, $q_3 = 0$: that is, $\sqrt{2}+\sqrt{3}$ is a
root of the polynomial
\[
  f(t) = t^4 - 10t^2 + 1.
\]
In fact, it is easily checked that $f(t)$ vanishes at all four combinations
$\pm\sqrt{2} \pm \sqrt{3}$.

It follows that
\[
  f(t) = (t-(-\sqrt{2}-\sqrt{3}))(t-(-\sqrt{2}+\sqrt{3}))
         (t-(\sqrt{2}-\sqrt{3}))(t-(\sqrt{2}+\sqrt{3}))
\]
and that $f(t)$ is irreducible over $\Q$ (no proper factor of $f(t)$ has
rational coefficients). Thus $\sqrt{2}+\sqrt{3}$ has degree $4$ over $\Q$.

Consider the composition of extensions
\[
  \Q \subseteq \Q(\sqrt{2}+\sqrt{3}) \subseteq \Q(\sqrt{2},\sqrt{3}).
\]
By \cref{coro:fields:1.11},
\[
  4 = [\Q(\sqrt{2}+\sqrt{3}) : \Q]
    \leq [\Q(\sqrt{2},\sqrt{3}) : \Q].
\]
On the other hand we know that $[\Q(\sqrt{2},\sqrt{3}) : \Q] \leq 4$;
thus we can conclude that this degree is exactly $4$, and it follows that
\[
  \Q(\sqrt{2},\sqrt{3}) = \Q(\sqrt{2}+\sqrt{3}):
\]
the extension was simple to begin with. (In due time we will prove that
this is a typical example, not a contrived one.)

Staring now at the composition
\[
  \Q \subseteq \Q(\sqrt{2}) \subseteq \Q(\sqrt{2},\sqrt{3}),
\]
\cref{prop:fields:1.10} tells us that
$[\Q(\sqrt{2},\sqrt{3}) : \Q(\sqrt{2})] = 2$.
That is, a side-effect of the computation carried out above is that the
polynomial $t^2 - 3$ must be irreducible over $\Q(\sqrt{2})$; this is
not surprising, but note that this method is very different from anything
we encountered back in \S{}V.5.

The fact that the other roots of the minimal polynomial of
$\sqrt{2}+\sqrt{3}$ `looked so much like' $\sqrt{2}+\sqrt{3}$ is not
surprising. Indeed, since $\Q(\sqrt{2}) \subseteq \Q(\sqrt{2},\sqrt{3})$
is a simple extension, we know (cf.\ \cref{prop:fields:1.5} and following)
that there is an automorphism of $\Q(\sqrt{2},\sqrt{3})$ fixing
$\Q(\sqrt{2})$ (and hence $\Q$) and swapping $\sqrt{3}$ and $-\sqrt{3}$.
Similarly, there is an automorphism fixing $\Q(\sqrt{3})$ and swapping
$\sqrt{2}$ and $-\sqrt{2}$. These automorphisms must act on the set of
roots of $f(t)$ (Exercise~1.6): applying them and their composition to
$\sqrt{2}+\sqrt{3}$ produces the other three roots of $f(t)$.

In fact, at this point we know that
$G = \operatorname{Aut}_{\Q}(\Q(\sqrt{2},\sqrt{3}))$ must consist of $4$
elements (by \cref{coro:fields:1.7}) and has at least two elements of
order $2$ (both automorphisms found above have order $2$). This is enough
to conclude that $G$ is not a cyclic group, and hence it must be isomorphic
to the group $\Z/2\Z \times \Z/2\Z$.
\end{exam}

\subsection*{Exercises}

\begin{exer}{}{fields:exer:1.1}
Prove that if $k \subseteq K$ is a field extension, then
$\operatorname{char} k = \operatorname{char} K$. Prove that the category
$\mathbf{Fld}$ has no initial object.
\end{exer}

\begin{exer}{}{fields:exer:1.2}
Define carefully the category $\mathbf{Fld}_k$ of extensions of $k$.
\end{exer}

\begin{exer}{}{fields:exer:1.3}
Let $k \subseteq F$ be a field extension, and let $\alpha \in F$. Prove
that the field $k(\alpha)$ consists of all the elements of $F$ which may
be written as rational functions in $\alpha$, with coefficients in $k$.
Why does this not give (in general) an onto homomorphism $k(t) \to
k(\alpha)$?
\end{exer}

\begin{exer}{}{fields:exer:1.4}
Let $k \subseteq k(\alpha)$ be a simple extension, with $\alpha$
transcendental over $k$. Let $E$ be a subfield of $k(\alpha)$ properly
containing $k$. Prove that $k(\alpha)$ is a finite extension of $E$.
\end{exer}

\begin{exer}{}{fields:exer:1.5}
(Cf.\ \cref{exam:fields:exam:1.4}.)
\begin{itemize}
\item Prove that there is exactly one subfield of $\R$ isomorphic to
  $\Q[t]/(t^2 - 2)$.
\item Prove that there are exactly three subfields of $\C$ isomorphic to
  $\Q[t]/(t^3 - 2)$.
\end{itemize}
From a `topological' point of view, one of these copies of $\Q[t]/(t^3 -
2)$ looks very different from the other two: it is not dense in $\C$,
but the others are.
\end{exer}

\begin{exer}{}{fields:exer:1.6}
Let $k \subseteq F$ be a field extension, and let $f(x) \in k[x]$ be a
polynomial. Prove that $\operatorname{Aut}_k(F)$ acts on the set of
roots of $f(x)$ contained in $F$. Provide examples showing that this
action need not be transitive or faithful.
\end{exer}

\begin{exer}{}{fields:exer:1.7}
Let $k \subseteq F$ be a field extension, and let $\alpha \in F$ be
algebraic over $k$.
\begin{itemize}
\item Suppose $p(x) \in k[x]$ is an irreducible monic polynomial such
  that $p(\alpha) = 0$; prove that $p(x)$ is the minimal polynomial of
  $\alpha$ over $k$, in the sense of \cref{prop:fields:1.3}.
\item Let $f(x) \in k[x]$. Prove that $f(\alpha) = 0$ if and only if
  $p(x) \mid f(x)$.
\item Show that the minimal polynomial of $\alpha$ is the minimal
  polynomial of a certain $k$-linear transformation of $F$, in the
  sense of Definition~VI.6.12.
\end{itemize}
\end{exer}

\begin{exer}{}{fields:exer:1.8}
Let $f(x) \in k[x]$ be a polynomial over a field $k$ of degree $d$, and
let $\alpha_1, \dots, \alpha_d$ be the roots of $f(x)$ in an extension
of $k$ where the polynomial factors completely. For a subset $I \subseteq
\{1, \dots, d\}$, denote by $\alpha_I$ the sum $\sum_{i \in I} \alpha_i$.
Assume that $\alpha_I \in k$ only for $I = \emptyset$ and $I =
\{1, \dots, d\}$. Prove that $f(x)$ is irreducible over $k$.
\end{exer}

\begin{exer}{}{fields:exer:1.9}
Let $k$ be a finite field. Prove that the order $|k|$ is a power of a
prime integer.
\end{exer}

\begin{exer}{}{fields:exer:1.10}
Let $k$ be a field. Prove that the ring of square $n \times n$ matrices
$M_n(k)$ contains an isomorphic copy of every extension of $k$ of degree
$\leq n$. (Hint: If $k \subseteq F$ is an extension of degree $n$ and
$\alpha \in F$, then `multiplication by $\alpha$' is a $k$-linear
transformation of $F$.)
\end{exer}
\begin{exer}{}{fields:exer:1.11}
Let $k \subseteq F$ be a finite field extension, and let $p(x)$ be the
characteristic polynomial of the $k$-linear transformation of $F$ given
by multiplication by $\alpha$. Prove that $p(\alpha) = 0$.

This gives an effective way to find a polynomial satisfied by an element
of an extension. Use it to find a polynomial satisfied by $\sqrt{2} +
\sqrt{3}$ over $\Q$, and compare this method with the one used in
\cref{exam:fields:exam:1.19}.
\end{exer}

\begin{exer}{}{fields:exer:1.12}
Let $k \subseteq F$ be a finite field extension, and let $\alpha \in F$.
The \emph{norm} of $\alpha$, $N_{k \subseteq F}(\alpha)$, is the
determinant of the linear transformation of $F$ given by multiplication
by $\alpha$ (cf.\ Exercise~1.11, Definition~VI.6.4). Prove that the norm
is multiplicative: for $\alpha, \beta \in F$,
\[
  N_{k \subseteq F}(\alpha\beta)
  = N_{k \subseteq F}(\alpha)\, N_{k \subseteq F}(\beta).
\]
Compute the norm of a complex number viewed as an element of the
extension $\R \subseteq \C$ (and marvel at the excellent choice of
terminology). Do the same for elements of an extension $\Q(\sqrt{d})$ of
$\Q$, where $d$ is an integer that is not a square, and compare the
result with Exercise~III.4.10.
\end{exer}

\begin{exer}{}{fields:exer:1.13}
Define the \emph{trace} $\operatorname{tr}_{k \subseteq F}(\alpha)$ of an
element $\alpha$ of a finite extension $F$ of a field $k$ by following
the lead of Exercise~1.12. Prove that the trace is additive:
\[
  \operatorname{tr}_{k \subseteq F}(\alpha + \beta)
  = \operatorname{tr}_{k \subseteq F}(\alpha)
  + \operatorname{tr}_{k \subseteq F}(\beta)
\]
for $\alpha, \beta \in F$. Compute the trace of an element of an
extension $\Q \subseteq \Q(\sqrt{d})$, for $d$ an integer that is not
a square.
\end{exer}

\begin{exer}{}{fields:exer:1.14}
Let $k \subseteq k(\alpha)$ be a simple algebraic extension, and let
$x^d + a_{d-1}x^{d-1} + \dots + a_0$ be the minimal polynomial of
$\alpha$ over $k$. Prove that
\[
  \operatorname{tr}_{k \subseteq k(\alpha)}(\alpha) = -a_{d-1}
  \quad\text{and}\quad
  N_{k \subseteq k(\alpha)}(\alpha) = (-1)^d a_0.
\]
(Cf.\ Exercises~1.12 and~1.13.)
\end{exer}

\begin{exer}{}{fields:exer:1.15}
Let $k \subseteq F$ be a finite extension, and let $\alpha \in F$. Assume
$[F : k(\alpha)] = r$. Prove that
\[
  \operatorname{tr}_{k \subseteq F}(\alpha)
  = r\,\operatorname{tr}_{k \subseteq k(\alpha)}(\alpha)
  \quad\text{and}\quad
  N_{k \subseteq F}(\alpha)
  = N_{k \subseteq k(\alpha)}(\alpha)^r.
\]
(Cf.\ Exercises~1.12 and~1.13.) (Hint: If $f_1, \dots, f_r$ is a basis
of $F$ over $k(\alpha)$ and $\alpha$ has degree $d$ over $k$, then
$(f_i \alpha^j)_{\substack{i = 1, \dots, r \\ j = 0, \dots, d-1}}$ is
a basis of $F$ over $k$. The matrix corresponding to multiplication by
$\alpha$ with respect to this basis consists of $r$ identical square
blocks.)
\end{exer}

\begin{exer}{}{fields:exer:1.16}
Let $k \subseteq L \subseteq F$ be fields, and let $\alpha \in F$. If
$k \subseteq k(\alpha)$ is a finite extension, then $L \subseteq
L(\alpha)$ is finite and $[L(\alpha) : L] \leq [k(\alpha) : k]$.
\end{exer}

\begin{exer}{}{fields:exer:1.17}
Let $k \subseteq F = k(\alpha_1, \dots, \alpha_n)$ be a finitely
generated extension. Prove that the evaluation map
\[
  k[t_1, \dots, t_n] \to F, \quad t_i \mapsto \alpha_i
\]
is an epimorphism of rings (although it need not be onto).
\end{exer}

\begin{exer}{}{fields:exer:1.18}
Let $R$ be a ring sandwiched between a field $k$ and an algebraic
extension $F$ of $k$. Prove that $R$ is a field.

Is it necessary to assume that the extension is algebraic?
\end{exer}

\begin{exer}{}{fields:exer:1.19}
Let $k \subseteq F$ be a field extension of degree $p$, a prime integer.
Prove that there are no subrings of $F$ properly containing $k$ and
properly contained in $F$. (Use Exercise~1.18.)
\end{exer}

\begin{exer}{}{fields:exer:1.20}
Let $p$ be a prime integer, and let $\alpha = \sqrt[p]{2} \in \R$. Let
$g(x) \in \Q[x]$ be any nonconstant polynomial of degree $< p$. Prove
that $\alpha$ may be expressed as a polynomial in $g(\alpha)$ with
rational coefficients.

Prove that an analogous statement for $\sqrt[4]{2}$ is false.
\end{exer}
\begin{exer}{}{fields:exer:1.21}
Let $k \subseteq F$ be a field extension, and let $E$ be the intermediate
field consisting of the elements of $F$ which are algebraic over $k$. For
$\alpha \in F$, prove that $\alpha$ is algebraic over $E$ if and only if
$\alpha \in E$. Deduce that $\overline{\Q}$ is algebraically closed.
\end{exer}

\begin{exer}{}{fields:exer:1.22}
Let $k \subseteq F$ be a field extension, and let $\alpha, \beta \in F$
be algebraic, of degree $d$, $e$, respectively. Assume $d$, $e$ are
relatively prime, and let $p(x)$ be the minimal polynomial of $\beta$
over $k$. Prove $p(x)$ is irreducible over $k(\alpha)$.
\end{exer}

\begin{exer}{}{fields:exer:1.23}
Express $\sqrt{2}$ explicitly as a polynomial function in $\sqrt{2} +
\sqrt{3}$ with rational coefficients.
\end{exer}

\begin{exer}{}{fields:exer:1.24}
Generalize the situation examined in \cref{exam:fields:exam:1.19}: let
$k$ be a field of characteristic $\neq 2$, and let $a, b \in k$ be
elements that are not squares in $k$; prove that $k(\sqrt{a}, \sqrt{b})
= k(\sqrt{a} + \sqrt{b})$.

Prove that $k(\sqrt{a}, \sqrt{b})$ has degree $2$, resp., $4$, over $k$
according to whether $ab$ is, resp., is not, a square in $k$.
\end{exer}

\begin{exer}{}{fields:exer:1.25}
Let $\xi := \sqrt{2 + \sqrt{2}}$.
\begin{itemize}
\item Find the minimal polynomial of $\xi$ over $\Q$, and show that
  $\Q(\xi)$ has degree $4$ over $\Q$.
\item Prove that $\sqrt{2 - \sqrt{2}}$ is another root of the minimal
  polynomial of $\xi$.
\item Prove that $\sqrt{2 - \sqrt{2}} \in \Q(\xi)$. (Hint:
  $(a+b)(a-b) = a^2 - b^2$.)
\item By \cref{prop:fields:1.5}, sending $\xi$ to $\sqrt{2 - \sqrt{2}}$
  defines an automorphism of $\Q(\xi)$ over $\Q$. Find the matrix of
  this automorphism w.r.t.\ the basis $1, \xi, \xi^2, \xi^3$.
\item Prove that $\operatorname{Aut}_{\Q}(\Q(\xi))$ is cyclic of order
  $4$.
\end{itemize}
\end{exer}

\begin{exer}{}{fields:exer:1.26}
Let $k \subseteq F$ be a field extension, let $I$ be an indexing set,
and let $\{\alpha_i\}_{i \in I}$ be a choice of elements of $F$. This
choice determines a homomorphism $\phi$ of $k$-algebras from the
polynomial ring $k[I]$ on the set $I$ to $F$ (the polynomial ring is a
free commutative $k$-algebra; cf.\ Proposition~III.6.4). We say that
$\{\alpha_i\}_{i \in I}$ is \emph{algebraically independent} over $k$
if $\phi$ is injective. For example, distinct elements
$\alpha_1, \dots, \alpha_n$ of $F$ are algebraically independent over
$k$ if there is no nonzero polynomial $f(x_1, \dots, x_n) \in
k[x_1, \dots, x_n]$ such that $f(\alpha_1, \dots, \alpha_n) = 0$.

Prove that $\alpha_1, \dots, \alpha_n$ are algebraically independent if
and only if the assignment $t_1 \mapsto \alpha_1, \dots, t_n \mapsto
\alpha_n$ defines a homomorphism of $k$-algebras (and hence an
isomorphism) from the field of rational functions $k(t_1, \dots, t_n)$
to $k(\alpha_1, \dots, \alpha_n)$.
\end{exer}

\begin{exer}{}{fields:exer:1.27}
With notation and terminology as in Exercise~1.26, the indexed set
$\{\alpha_i\}_{i \in I}$ is a \emph{transcendence basis} for $F$ over
$k$ if it is a maximal algebraically independent set in $F$.
\begin{itemize}
\item Prove that $\{\alpha_i\}_{i \in I}$ is a transcendence basis for
  $F$ over $k$ if and only if it is algebraically independent and $F$
  is algebraic over $k(\{\alpha_i\}_{i \in I})$.
\item Prove that transcendence bases exist. (Zorn)
\item Prove that any two transcendence bases for $F$ over $k$ have the
  same cardinality. (Mimic the proof of Proposition~VI.1.9. Don't feel
  too bad if you prefer to deal only with the case of finite
  transcendence bases.)
\end{itemize}
The cardinality of a transcendence basis is called the \emph{transcendence
degree} of $F$ over $k$, denoted $\operatorname{tr.deg}_{k \subseteq F}$.
\end{exer}

\begin{exer}{}{fields:exer:1.28}
Let $k \subseteq E \subseteq F$ be field extensions. Prove that
$\operatorname{tr.deg}_{k \subseteq F}$ is finite if and only if both
$\operatorname{tr.deg}_{k \subseteq E}$ and
$\operatorname{tr.deg}_{E \subseteq F}$ (see Exercise~1.27) are finite,
and in this case
\[
  \operatorname{tr.deg}_{k \subseteq F}
  = \operatorname{tr.deg}_{k \subseteq E}
  + \operatorname{tr.deg}_{E \subseteq F}.
\]
\end{exer}

\begin{exer}{}{fields:exer:1.29}
An extension $k \subseteq F$ is \emph{purely transcendental} if it
admits a transcendence basis $\{\alpha_i\}_{i \in I}$ (see
Exercise~1.27) such that $F = k(\{\alpha_i\}_{i \in I})$.

Prove that any field extension $k \subseteq F$ may be decomposed as a
purely transcendental extension followed by an algebraic extension. (Not
all field extensions may be decomposed as an algebraic extension followed
by a purely transcendental extension.)
\end{exer}

\begin{exer}{}{fields:exer:1.30}
Let $k \subseteq k(\alpha)$ be a simple extension, with $\alpha$
transcendental over $k$. Let $E$ be a subfield of $k(\alpha)$ properly
containing $k$. Prove that $\operatorname{tr.deg}_{k \subseteq E} = 1$.

L\"uroth's theorem asserts that in this situation $k \subseteq E$ is
itself a simple transcendental extension of $k$; that is, it is purely
transcendental (Exercise~1.29).
\end{exer}

\section{Algebraic closure, Nullstellensatz, and a little algebraic
geometry}
\label{sec:fields:algebraic-closure-nullstellensatz-and-a-little-algebraic}

One of the most important extensions of a field $k$ is its algebraic
closure $k \subseteq \overline{k}$; this was mentioned already in
\S{}V.5.2, and we can now prove that $\overline{k}$ exists and is unique
(up to isomorphism). Once this circle of ideas is approached, the
temptation to say a few words about the important result known as
Hilbert's Nullstellensatz, at the cost of a small digression into
algebraic geometry, will simply be irresistible.

\subsection{Algebraic closure}
\label{subsec:fields:algebraic-closure}
Recall (Definition~V.5.9) that a field $K$ is algebraically closed if all
irreducible polynomials in $K[x]$ have degree $1$, that is, if every
polynomial in $K[x]$ factors completely as a product of linear terms.
Equivalently (Exercise~III.4.21) every maximal ideal in $K[x]$ is of the
form $(x - c)$, for $c \in K$. Now that we have a little more vocabulary,
we can recast this definition yet again, as follows.
\begin{lem}{}{fields:2.1}
For a field $K$, the following are equivalent:
\begin{enumerate}[label=(\roman*)]
  \item $K$ is algebraically closed.
  \item $K$ has no nontrivial algebraic extensions.
  \item If $K \subseteq L$ is any extension and $\alpha \in L$ is
    algebraic over $K$, then $\alpha \in K$.
\end{enumerate}
\end{lem}
The proof is a straightforward application of the definitions and a good
exercise (Exercise~2.1).
\begin{defn}{}{fields:2.2}
An \emph{algebraic closure} of a field $k$ is an algebraic extension
$k \subseteq \overline{k}$ such that $\overline{k}$ is algebraically
closed.
\end{defn}
Of course, every polynomial $f(x) \in k[x]$ will split into linear factors
over $\overline{k}$: so does every polynomial in the much larger
$\overline{k}[x]$, as $\overline{k}$ is algebraically closed. The
requirement that $k \subseteq \overline{k}$ be algebraic ensures that no
other intermediate field $L$,
\[
  k \subseteq L \subsetneq \overline{k},
\]
may be algebraically closed: indeed, $L \subseteq \overline{k}$ would be a
nontrivial algebraic extension, contradicting \cref{lem:fields:2.1}. In
fact, the only subfield of $\overline{k}$ containing all roots of all
nonconstant polynomials in $k[x]$ is $\overline{k}$ itself
(Exercise~2.2); thus, the algebraic closure of a field is `as small as
possible' subject to this requirement. Not surprisingly, $\overline{k}$
turns out to be unique up to isomorphism, so that we can speak of
\emph{the} algebraic closure of $k$.
\begin{thm}{}{fields:2.3}
Every field $k$ admits an algebraic closure $k \subseteq \overline{k}$;
this extension is unique up to isomorphism.
\end{thm}
Concerning existence, the idea is to construct `by hand' a huge extension
$K$ of $k$ where every polynomial $f(x) \in k[x]$ factors completely. The
elements of $K$ which are algebraic over $k$ will form an algebraic
closure of $k$.

The construction is done in steps, each step including `one more root' of
all nonconstant polynomials in $k[x]$. Here is a formalization of this
step.
\begin{lem}{}{fields:2.4}
Let $k$ be a field. Then there exists an extension $k \subseteq K$ such
that every nonconstant polynomial $f(x) \in k[x]$ has at least one root
in $K$.
\end{lem}
\begin{proof}
(This construction is apparently due to \name{Emil Artin}.) Consider a
set $T = \{t_f\}$ in bijection with the set of nonconstant monic
polynomials $f(x) \in k[x]$, and let $k[T]$ be the corresponding
polynomial ring\footnote{Here is another rare occasion where we need to
consider polynomial rings with possibly infinitely many indeterminates. An
element of $k[T]$ is simply an ordinary polynomial with coefficients in
$k$, involving a finite number of indeterminates from $T$.} in all the
indeterminates $t_f$. Let $I \subseteq k[T]$ be the ideal generated by
all polynomials $f(t_f)$.

Then $I$ is a proper ideal. Indeed, otherwise we could write
\begin{equation}
  \label{eq:fields:2.4-star}
  1 = \sum_{i=1}^{n} a_i \cdot f_i(t_{f_i}),
\end{equation}
where $a_i \in k[T]$. I claim that this cannot be done: indeed, we can
construct an extension $k \subseteq F$ where the polynomials
$f_1(x), \dots, f_n(x)$ have roots $\alpha_1, \dots, \alpha_n$,
respectively (apply Proposition~V.5.7 $n$ times); view
\eqref{eq:fields:2.4-star} as an identity in $F[T]$, and plug in
$t_{f_i} = \alpha_i$, obtaining
\[
  1 = \sum_{i=1}^{n} a_i \cdot f_i(\alpha_i)
    = \sum_{i=1}^{n} a_i \cdot 0 = 0,
\]
which is nonsense.

Since $I$ is proper, it is contained in a maximal ideal $\mathfrak{m}$
(Proposition~V.3.5). Thus, we obtain a field extension
\[
  k \subseteq K := \frac{k[T]}{\mathfrak{m}};
\]
by construction every nonconstant monic (and hence every nonconstant)
polynomial $f(x)$ has a root in $K$, namely the coset of $t_f$.
\end{proof}
The field constructed in \cref{lem:fields:2.4} contains at least one root
of each nonconstant polynomial $f(x) \in k[x]$, but this is not good
enough: we are seeking a field which contains all roots of such
polynomials. Equivalently, not only should $f(x)$ have (at least) one
linear factor $x - \alpha$ in $K[x]$, but we need the quotient polynomial
$f(x)/(x - \alpha) \in K[x]$ to have a linear factor and the quotient by
that new factor to have one, etc. This prompts us to consider a whole
chain of extensions
\[
  k \subseteq K_1 \subseteq K_2 \subseteq K_3 \subseteq \dots
\]
where $K_1$ is obtained from $k$ by applying \cref{lem:fields:2.4}, $K_2$
is similarly obtained from $K_1$, etc.

Now consider the union $L$ of this chain\footnote{This is an example of
direct limit, a notion which we will introduce more formally in
\S{}VIII.1.4.}. For every two $a, b \in L$, there is an $i$ such that
$a, b \in K_i$; we can define $a + b$, $a \cdot b$ in $L$ by adopting
their definition in $K_i$, and the result does not depend on the choice
of $i$. It follows easily that $L$ is a field.
\begin{fact}{}{fields:2.5}
The field $L$ is algebraically closed.
\end{fact}
\begin{proof}
If $f(x) \in L[x]$ is a nonconstant polynomial, then $f(x) \in K_i[x]$
for some $i$; hence $f(x)$ has a root in $K_{i+1} \subseteq L$. That is,
every nonconstant polynomial in $L[x]$ has a root in $L$, as needed.
\end{proof}
\paragraph{Proof of the existence of algebraic closures.}
The existence of algebraic closures is now completely transparent, because
of the following simple observation:
\begin{lem}{}{fields:2.6}
Let $k \subseteq L$ be a field extension, with $L$ algebraically closed.
Let
\[
  \overline{k} := \{\alpha \in L \mid \alpha \text{ is algebraic over }
  k\}.
\]
Then $\overline{k}$ is an algebraic closure of $k$.
\end{lem}
The construction reviewed above provides us with an algebraically closed
field $L$ containing any given field $k$, so the lemma is all we need to
prove.

By \cref{coro:fields:1.16}, $\overline{k}$ is a field, and the extension
$k \subseteq \overline{k}$ is tautologically algebraic. To verify that
$\overline{k}$ is algebraically closed, let $k \subseteq \overline{k}(\alpha)$
be a simple algebraic extension. The minimal polynomial of $\alpha$ has a
root in $L$ since $L$ is algebraically closed, so by versality
(Proposition~V.5.7) there exists an embedding $\overline{k}(\alpha)
\subseteq L$. We can then view $\alpha$ as an element of $L$;
$k \subseteq \overline{k} \subseteq \overline{k}(\alpha)$ is a
composition of algebraic extensions, so $k \subseteq \overline{k}(\alpha)$
is algebraic (\cref{coro:fields:1.18}), and in particular $\alpha$ is
algebraic over $k$. But then $\alpha \in \overline{k}$, by definition of
the latter. It follows that $\overline{k}$ is algebraically closed, by
\cref{lem:fields:2.1}. \qed

\begin{rem}{}{fields:2.7}
Zorn's lemma was sneakily used in this proof (we used the existence of
maximal ideals, which relies upon Zorn's lemma), and it is natural to
wonder whether the existence of algebraic closures may be another
statement equivalent to the axiom of choice. Apparently, this is not the
case.\footnote{Allegedly, the existence of algebraic closures is a
consequence of the compactness theorem for first-order logic, which is
known to be weaker than the axiom of choice.}
\end{rem}
Next, we deal with uniqueness. It may be tempting to try to set things up
in such a way that algebraic closures would end up being solutions to a
universal problem; this would guarantee uniqueness by abstract nonsense.
However, we run into the same obstacle encountered in Proposition~V.5.7:
morphisms between extensions depend on choices of roots of polynomials,
so that algebraic closures are not `universal' in the most straightforward
sense of the term. Therefore, a bit of independent work is needed.
\begin{lem}{}{fields:2.8}
Let $k \subseteq L$ be a field extension, with $L$ algebraically closed.
Let $k \subseteq F$ be any algebraic extension. Then there exists a
morphism of extensions $i : F \to L$.
\end{lem}
As pointed out above, $i$ is by no means unique!
\begin{proof}
This argument also relies on Zorn's lemma. Consider the set $\mathcal{Z}$
of homomorphisms
\[
  i_K : K \to L
\]
where $K$ is an intermediate field, $k \subseteq K \subseteq F$, and
$i_K$ restricts to the identity on $k$; $\mathcal{Z}$ is nonempty, since
the extension $i_k : k \hookrightarrow L$ defines an element of
$\mathcal{Z}$. We give a poset structure to $\mathcal{Z}$ by defining
$i_K \leq i_{K'}$ if $K \subseteq K' \subseteq F$ and $i_{K'}$ restricts
to $i_K$ on $K$. To verify that every chain $\mathcal{C}$ in $\mathcal{Z}$
has an upper bound in $\mathcal{Z}$, let $K_{\mathcal{C}}$ be the union
of the sources of all $i_K \in \mathcal{C}$ ($K_{\mathcal{C}}$ is clearly
a field); if $\alpha \in K_{\mathcal{C}}$, define $i_{K_{\mathcal{C}}}(\alpha)$
to be $i_K(\alpha)$, where $i_K$ is any element of $\mathcal{C}$ such
that $\alpha \in K$. This prescription is clearly independent of the
chosen $K$ and defines a homomorphism $K_{\mathcal{C}} \to L$ restricting
to the identity on $k$. This homomorphism is an upper bound for
$\mathcal{C}$.

By Zorn's lemma, $\mathcal{Z}$ admits a maximal element $i_G$,
corresponding to an intermediate field $k \subseteq G \subseteq F$. Let
$H = i_G(G)$ be the image of $G$ in $L$.

I claim that $G = F$: this will prove the statement, because it will imply
that there is a homomorphism $i_F : F \to L$ extending the identity on $k$.

Arguing by contradiction, assume that there exists an $\alpha \in F
\setminus G$, and consider the extension $G \subseteq G(\alpha)$. Since
$\alpha \in F$ is algebraic over $k$, it is algebraic over $G$; thus, it
is a root of an irreducible polynomial $g(x) \in G[x]$. Consider the
induced homomorphism
\[
  i_G : G[x] \to H[x],
\]
and let $h(x) = i_G(g(x))$. Then $h(x)$ is an irreducible polynomial
over $H$, and it has a root $\beta$ in $L$ (this is where we use the
hypothesis that $L$ is algebraically closed!). We are in the situation
considered in \cref{prop:fields:1.5}: we have an isomorphism of fields
$i_G : G \to H$; we have simple extensions $G(\alpha)$, $H(\beta)$; and
$\alpha$, $\beta$ are roots of irreducible polynomials $g(x)$,
$h(x) = i_G(g(x))$. By \cref{prop:fields:1.5}, $i_G$ lifts to an
isomorphism
\[
  i_{G(\alpha)} : G(\alpha) \to H(\beta) \subseteq L
\]
sending $\alpha$ to $\beta$. This contradicts the maximality of $i_G$;
hence $G = F$, concluding the argument.
\end{proof}
\paragraph{Proof of uniqueness of algebraic closures.}
Let $k \subseteq \overline{k}$, $k \subseteq \overline{k}_1$ be two
algebraic closures of $k$; we have to prove that there exists an
isomorphism $\overline{k}_1 \to \overline{k}$ extending the identity on
$k$.

Since $k \subseteq \overline{k}_1$ is algebraic and $\overline{k}$ is
algebraically closed, by \cref{lem:fields:2.8} there exists a
homomorphism $i : \overline{k}_1 \to \overline{k}$ extending the identity
on $k$; this homomorphism is trivially injective, since $\overline{k}_1$
is a field. It is also surjective, by \cref{lem:fields:2.1}: otherwise
$\overline{k}_1 \subseteq \overline{k}$ is a nontrivial algebraic
extension, contradicting the fact that $\overline{k}_1$ is algebraically
closed. Therefore $i$ is an isomorphism, as needed. \qed

\subsection{The Nullstellensatz}
\label{subsec:fields:the-nullstellensatz}
If $K$ is an algebraically closed field, then every maximal ideal in
$K[x]$ is of the form $(x - c)$, for $c \in K$ (Exercise~III.4.21).
This statement has a straightforward-looking generalization to polynomial
rings in more indeterminates: if $K$ is algebraically closed, then every
maximal ideal in $K[x_1, \dots, x_n]$ is of the form
$(x_1 - c_1, \dots, x_n - c_n)$.

Proving this statement is more challenging than it may seem from the looks
of it; it is one facet of the famous theorem known as \name{Hilbert}'s
Nullstellensatz (``theorem on the position of zeros''). The Nullstellensatz
has made a few cameo appearances earlier (see Example~III.4.15 and
\S{}III.6.5 in particular), and we will spend a little time on it now.
We will not prove the Nullstellensatz in its natural generality, that is,
for all fields; this would take us too far. Actually, there are reasonably
short proofs of the theorem in this generality, but the short arguments I
have run into end up replacing a wider (and useful) context with ingenious
cleverness; I do not find these arguments particularly insightful or
memorable.

By contrast, the theorem has a very simple and memorable proof if we make
the further assumption that the field is uncountable. Therefore, the
reader will have a complete proof of the theorem (in the form given below,
which does not assume the field to be algebraically closed) for fields
such as $\R$ and $\C$ but will have to trust on faith regarding other
extremely important fields such as $\Q$.

The most vivid applications of the Nullstellensatz involve basic
definitions in algebraic geometry, and we will get a small taste of this
in the next section; but the theorem itself is best understood in the
context of the considerations on `finiteness' mentioned in \S{}1.3; see
especially \cref{rem:fields:1.14}. I pointed out that if $k \subseteq F$
is finitely generated as a field extension, then $F$ need not be finitely
generated (that is, `finite-type') as a $k$-algebra, and I raised the
question of understanding extensions $F$ that are finite-type
$k$-algebras. The following result answers this question. It is my
favorite version of the Nullstellensatz; therefore I will name it so:
\begin{thm}{Nullstellensatz}{fields:2.9}
Let $k \subseteq F$ be a field extension, and assume that $F$ is a
finite-type $k$-algebra. Then $k \subseteq F$ is a finite (hence
algebraic) extension.
\end{thm}
The reader should pause and savor this statement for a moment. As pointed
out in \S{}III.6.5, an algebra $k \to S$ may very well be of finite type
(as an algebra), without being finite (i.e., finitely generated as a
module): $S = k[t]$ is an obvious example. The content of
\cref{thm:fields:2.9} is that the distinction between finite-type and
finite disappears if $S$ is a field.

As mentioned above, we will only prove this statement under an additional
(and somewhat unnatural) assumption, which reduces the argument to
elementary linear algebra.
\begin{proof}[Proof for uncountable fields]
Assume that $k$ is uncountable. Let $k \subseteq F$ be a field extension,
and assume that $F$ is finitely generated as an algebra over $k$; in
particular, it is finitely generated as a field extension. We have to
prove that $k \subseteq F$ is a finite extension, or equivalently (by
\cref{prop:fields:1.15}) that it is algebraic, that is, that every
$\alpha \in F \setminus k$ is the root of a nonzero $f(x) \in k[x]$.

Now, by hypothesis $F$ is the quotient of a polynomial ring over $k$:
\[
  k[x_1, \dots, x_n] \twoheadrightarrow F;
\]
this implies that there is a countable basis of $F$ as a vector space
over $k$, because this is the case for $k[x_1, \dots, x_n]$. Consider
then the set
\[
  \left\{ \frac{1}{\alpha - c} \right\}_{c \in k};
\]
this is an uncountable subset of $F$; therefore it is linearly dependent
over $k$ (Proposition~VI.1.9). That is, there exist distinct
$c_1, \dots, c_m \in k$ and nonzero coefficients
$\lambda_1, \dots, \lambda_m \in k$ such that
\[
  \frac{\lambda_1}{\alpha - c_1} + \cdots + \frac{\lambda_m}{\alpha - c_m}
  = 0.
\]
Expressing the left-hand side with a common denominator gives
\[
  \frac{f(\alpha)}{g(\alpha)} = 0,
\]
for $f(x) \neq 0$ in $k[x]$ and $g(\alpha) \neq 0$ (Exercise~2.4).
This implies $f(\alpha) = 0$ for a nonzero $f(x) \in k[x]$, and we are
done.
\end{proof}
Regardless of its proof, the form of the Nullstellensatz given in
\cref{thm:fields:2.9} does not look much like the other results I have
mentioned as associated with it. For one thing, it is not too clear why
\cref{thm:fields:2.9} should be called a theorem on the `position of
zeros'. The result deserves a little more attention.

The form stated at the beginning of this section is obtained as follows:
\begin{coro}{}{fields:2.10}
Let $K$ be an algebraically closed field, and let $I$ be an ideal of
$K[x_1, \dots, x_n]$. Then $I$ is maximal if and only if
\[
  I = (x_1 - c_1, \dots, x_n - c_n),
\]
for $c_1, \dots, c_n \in K$.
\end{coro}
\begin{proof}
For $c_1, \dots, c_n \in K$,
\[
  \frac{K[x_1, \dots, x_n]}{(x_1 - c_1, \dots, x_n - c_n)} \cong K
\]
(Exercise~III.4.12) is a field; therefore $(x_1 - c_1, \dots, x_n - c_n)$
is maximal. Conversely, let $\mathfrak{m}$ be a maximal ideal in
$K[x_1, \dots, x_n]$, so that the quotient is a field and the natural map
\[
  K \to L := \frac{K[x_1, \dots, x_n]}{\mathfrak{m}}
\]
is a field extension. The field $L$ is finitely generated as a
$K$-algebra; therefore $K \subseteq L$ is an algebraic extension by
\cref{thm:fields:2.9}. Since $K$ is algebraically closed, this implies
that $L = K$ (\cref{lem:fields:2.1}); therefore, there is a surjective
homomorphism
\[
  \phi : K[x_1, \dots, x_n] \to K
\]
such that $\mathfrak{m} = \ker\phi$. Let $c_i := \phi(x_i)$. Then
\[
  (x_1 - c_1, \dots, x_n - c_n) \subseteq \ker\phi = \mathfrak{m};
\]
but $(x_1 - c_1, \dots, x_n - c_n)$ is maximal, so this implies
$\mathfrak{m} = (x_1 - c_1, \dots, x_n - c_n)$, as needed.
\end{proof}
Note how very (trivially) false the statement of \cref{coro:fields:2.10}
is over fields which are not algebraically closed: $(x^2 + 1)$ is maximal
in $\R[x]$. No such silliness may occur over $\C$, for any number of
variables.

\subsection{A little affine algebraic geometry}
\label{subsec:fields:a-little-affine-algebraic-geometry}
\cref{coro:fields:2.10} is one of the main reasons why `classical'
algebraic geometry developed over $\C$ rather than fields such as $\R$ or
$\Q$: it may be interpreted as saying that if $K$ is algebraically closed,
then there is a natural bijection between the points of the product space
$K^n$ and the maximal ideals in the ring $K[x_1, \dots, x_n]$. This is
the beginning of a fruitful dictionary translating geometry into algebra
and conversely. It is worthwhile formalizing this statement and exploring
other `geometric' consequences of the Nullstellensatz.

For $K$ a field, $\mathbb{A}^n_K$ denotes the \emph{affine space} of
dimension $n$ over $K$, that is, the set of $n$-tuples of elements of
$K$:
\[
  \mathbb{A}^n_K = \{(c_1, \dots, c_n) \mid c_i \in K\}.
\]
Elements of $\mathbb{A}^n_K$ are called \emph{points}.

\emph{Objection}: Why don't we just use `$K^n$' for this object? Because
the latter already stands for the standard $n$-dimensional vector space
over $K$; so it carries with it a certain baggage: the tuple
$(0, \dots, 0)$ is special, so are the vector subspaces of $K^n$, etc.
By contrast, no point of $\mathbb{A}^n_K$ is `special'; we can carry any
point to any other point by a simple translation. Similarly, although it
is useful to consider affine subspaces, these (unlike vector subspaces)
are not required to go through the origin. Affine geometry and linear
algebra are different in many respects.

We have already established that, by the Nullstellensatz, if $K$ is
algebraically closed, then there is a natural bijection between the points
of $\mathbb{A}^n_K$ and the maximal ideals of the polynomial ring
$K[x_1, \dots, x_n]$. Concretely, the bijection works like this: the
point $p = (c_1, \dots, c_n)$ corresponds to the set of polynomials
\[
  \mathcal{I}(p)
  := \{f(x) \in K[x_1, \dots, x_n] \mid f(p) = 0\},
\]
where $f(p)$ is the result of applying the evaluation map:
$f(x_1, \dots, x_n) \mapsto f(c_1, \dots, c_n) \in K$. This of course
is nothing but the homomorphism $\phi : K[x_1, \dots, x_n] \to K$
defined by prescribing $x_i \mapsto c_i$; the set $\mathcal{I}(p)$ is
simply $\ker\phi$, which is the maximal ideal
$(x_1 - c_1, \dots, x_n - c_n)$.

This correspondence $p \mapsto \mathcal{I}(p)$ may be defined over every
field; the content of \cref{coro:fields:2.10} is that every maximal ideal
of $K[x_1, \dots, x_n]$ corresponds to a point of $\mathbb{A}^n_K$ in
this way, if $K$ is algebraically closed.

It is natural to upgrade the correspondence as follows (again, over
arbitrary fields): for every subset $S \subseteq \mathbb{A}^n_K$,
consider the ideal
\[
  \mathcal{I}(S)
  := \{f(x) \in K[x_1, \dots, x_n] \mid \forall p \in S,\, f(p) = 0\},
\]
that is, the set of polynomials `vanishing along $S$'; it is immediately
checked that this is indeed an ideal of $K[x_1, \dots, x_n]$. We can
also consider a correspondence in the reverse direction, from ideals of
$K[x_1, \dots, x_n]$ to subsets of $\mathbb{A}^n_K$, defined by setting
for every ideal $I$,
\[
  \mathcal{V}(I)
  := \{p = (c_1, \dots, c_n) \in \mathbb{A}^n_K
      \mid \forall f \in I,\, f(c_1, \dots, c_n) = 0\}.
\]
Thus, $\mathcal{V}(I)$ is the set of common solutions of all the
polynomial equations $f = 0$ as $f \in I$: the set of points `cut out in
$\mathbb{A}^n_K$' by the polynomials in $I$.

In its most elementary manifestation, the dictionary mentioned in the
beginning of the section consists precisely of the pair of correspondences
\[
  \left\{\text{subsets of }\mathbb{A}^n_K\right\}
  \begin{matrix}
    \xrightarrow{\;\mathcal{I}\;} \\[3pt]
    \xleftarrow{\;\mathcal{V}\;}
  \end{matrix}
  \left\{\text{ideals in }K[x_1,\dots,x_n]\right\}.
\]
The set $\mathcal{V}(I)$ is often (somewhat improperly) called the
\emph{variety} of $I$, while $\mathcal{I}(S)$ is (also improperly) the
\emph{ideal} of $S$.

The function $\mathcal{V}$ can be defined (using the same prescription)
for any set $A \subseteq K[x_1, \dots, x_n]$, whether $A$ is an ideal
or not; it is clear that $\mathcal{V}(A) = \mathcal{V}(I)$ where $I$ is
the ideal generated by $A$ (Exercise~2.5). \name{Hilbert}'s basis theorem
(Theorem~V.1.2) tells us something rather interesting: $K[x_1, \dots,
x_n]$ is Noetherian when $K$ is a field; hence every ideal $I$ is
generated by a finite number of elements. Thus, for every ideal $I$ there
exist polynomials $f_1, \dots, f_r$ such that (abusing notation a little)
\[
  \mathcal{V}(I) = \mathcal{V}(f_1, \dots, f_r).
\]
In other words, if a set may be defined by any set of polynomial
equations, it may in fact be defined by a finite set of polynomial
equations.\footnote{I distinctly remember being surprised by this fact
the first time I ran into it.}

I will pass in silence many simple-minded properties of the functions
$\mathcal{V}$, $\mathcal{I}$ (they reverse inclusions, they behave
reasonably well with respect to unions and intersections, etc.); the
reader should feel free to research the topic at will. But I want to
focus the reader's attention on the fact that (of course) $\mathcal{V}$,
$\mathcal{I}$ are not bijections; this is a problem, if we really want
to construct a dictionary between `geometry' and `algebra'.

For example, there are (in general) lots of subsets of $\mathbb{A}^n_K$
which are not cut out by a set of polynomial equations, and there are
lots of ideals in $K[x_1, \dots, x_n]$ which are not obtained as
$\mathcal{I}(S)$ for any subset $S \subseteq \mathbb{A}^n_K$.

I will leave to the reader the task of finding examples of the first
phenomenon (Exercise~2.6). The way around it is to restrict our attention
to subsets of $\mathbb{A}^n_K$ which are of the form $\mathcal{V}(I)$
for some ideal $I$.
\begin{defn}{}{fields:2.11}
An \emph{(affine) algebraic set} is a subset $S \subseteq \mathbb{A}^n_K$
such that there exists an ideal $I \subseteq K[x_1, \dots, x_n]$ for
which $S = \mathcal{V}(I)$.
\end{defn}
This definition may seem like a cheap patch: we do not really know what
the image of $\mathcal{V}$ is, so we label it in some way and proceed.
This is not entirely true --- the family of algebraic subsets of a given
$\mathbb{A}^n_K$ has a solid, recognizable structure: it is the family of
closed subsets in a topology on $\mathbb{A}^n_K$ (Exercise~2.7); this
topology is called the \emph{Zariski topology}.

Studying affine algebraic sets is the business of affine algebraic
geometry. The aim is to understand `geometric' properties of these sets
in terms of `algebraic' properties of corresponding ideals or other
algebraic entities associated with them.
\begin{exam}{}{fields:exam:2.12}
Here are pictures of two algebraic subsets of $\mathbb{A}^2_{\R}$: the
first is $\mathcal{V}((y - x^2))$; the second is $\mathcal{V}((y^2 -
x^3))$. What feature of the ideal $(y^2 - x^3)$ is responsible for the
`cusp' at $(0, 0)$ in the second picture? The reader will find out in any
course in elementary algebraic geometry.
\end{exam}
The fact that $\mathcal{I}$ is not surjective leads to interesting
considerations. The standard example of an ideal that is not the ideal of
any set is $(x^2)$ in $K[x]$: wherever $x^2$ vanishes, so does $x$;
hence if $x^2 \in \mathcal{I}(S)$ for a set $S$, then $x \in
\mathcal{I}(S)$ as well. This motivates the following definitions.
\begin{defn}{}{fields:2.13}
Let $I$ be an ideal in a commutative ring $R$. The \emph{radical} of $I$
is the ideal
\[
  \sqrt{I} := \{r \in R \mid \exists k \geq 0,\; r^k \in I\}.
\]
An ideal $I$ is a \emph{radical ideal} if $I = \sqrt{I}$.
\end{defn}
The reader should check that $\sqrt{I}$ is indeed an ideal; it may be
characterized as the intersection of all prime ideals containing $I$
(Exercise~2.8).
\begin{exam}{}{fields:exam:2.14}
An element of a commutative ring $R$ is nilpotent if and only if it is in
the radical of the ideal $(0)$ (cf.\ Exercise~V.4.19). The reader has
encountered this ideal, called the \emph{nilradical} of $R$, in the
exercises, beginning with Exercise~III.3.12.

Prime ideals $\mathfrak{p}$ are clearly radical: because $f^k \in
\mathfrak{p} \implies f \in \mathfrak{p}$, and therefore $\mathfrak{p}
\subseteq \sqrt{\mathfrak{p}}$ (the other inclusion holds trivially for
all ideals).
\end{exam}
\begin{lem}{}{fields:2.15}
Let $K$ be a field, and let $S$ be a subset of $\mathbb{A}^n_K$. Then
the ideal $\mathcal{I}(S)$ is a radical ideal of $K[x_1, \dots, x_n]$.
\end{lem}
\begin{proof}
The inclusion $\mathcal{I}(S) \subseteq \sqrt{\mathcal{I}(S)}$ holds for
every ideal, so it is trivially satisfied. To verify the inclusion
$\sqrt{\mathcal{I}(S)} \subseteq \mathcal{I}(S)$, let $f \in
\sqrt{\mathcal{I}(S)}$. Then there is an integer $k \geq 0$ such that
$f^k \in \mathcal{I}(S)$; that is, $(\forall p \in S),\; f(p)^k = 0$.
But then $(\forall p \in S),\; f(p) = 0$, proving $f \in \mathcal{I}(S)$,
as needed.
\end{proof}
In view of these considerations, for any field we can refine the
correspondence between subsets of $\mathbb{A}^n_K$ and ideals of
$K[x_1, \dots, x_n]$ as follows:
\[
  \left\{\text{algebraic subsets of }\mathbb{A}^n_K\right\}
  \begin{matrix}
    \xrightarrow{\;\mathcal{I}\;} \\[3pt]
    \xleftarrow{\;\mathcal{V}\;}
  \end{matrix}
  \left\{\text{radical ideals in }K[x_1,\dots,x_n]\right\};
\]
as I have argued, it is necessary to do so if we want to have a good
dictionary. In the rest of the section, $\mathcal{I}$ and $\mathcal{V}$
will be taken as acting between these two sets.

While we have tightened the correspondence substantially, the situation
is still less than idyllic. Over arbitrary fields, the function
$\mathcal{V}$ is now surjective by the very definition of an affine
algebraic subset (cf.\ Exercise~2.9); but it is not necessarily
injective. An immediate counterexample is offered over $K = \R$ by the
ideal $(x^2 + 1)$: this is maximal, hence prime, hence radical, and
\[
  \mathcal{V}(x^2 + 1) = \emptyset = \mathcal{V}(1).
\]
Here is where the Nullstellensatz comes to our aid, and here is where we
see what the Nullstellensatz has to do with the `zeros' of an ideal
($ = \mathcal{V}$ of that ideal).
\begin{prop}[Weak Nullstellensatz]{}{fields:2.16}
Let $K$ be an algebraically closed field, and let $I \subseteq
K[x_1, \dots, x_n]$ be an ideal. Then $\mathcal{V}(I) = \emptyset$ if
and only if $I = (1)$.
\end{prop}
\begin{proof}
If $I = (1)$, then $\mathcal{V}(I) = \emptyset$ by definition.
Conversely, assume that $I \neq (1)$. By Proposition~V.3.5, $I$ is then
contained in a maximal ideal $\mathfrak{m}$. Since $K$ is algebraically
closed, by \cref{coro:fields:2.10} we have
$\mathfrak{m} = (x_1 - c_1, \dots, x_n - c_n)$
for some $c_1, \dots, c_n \in K$; so for all $f(x_1, \dots, x_n) \in I$
there exist $g_1, \dots, g_n \in K[x_1, \dots, x_n]$ such that
\[
  f(x_1, \dots, x_n)
  = \sum_{i=1}^{n} g_i(x_1, \dots, x_n)(x_i - c_i).
\]
% In particular,
% ?
In particular,
\[
  f(c_1, \dots, c_n)
  = \sum_{i=1}^{n} g_i(c_1, \dots, c_n)(c_i - c_i) = 0:
\]
that is, $(c_1, \dots, c_n) \in \mathcal{V}(I)$. This proves
$\mathcal{V}(I) \neq \emptyset$ if $I \neq (1)$, and we are done.
\end{proof}
This is encouraging, but it does not seem to really prove that $\mathcal{V}$
is injective. As it happens, however, it does: we can derive from the
`weak' Nullstellensatz the following stronger result, which does imply
injectivity:
\begin{prop}[Strong Nullstellensatz]{}{fields:2.17}
Let $K$ be an algebraically closed field, and let $I \subseteq
K[x_1, \dots, x_n]$ be an ideal. Then
\[
  \mathcal{I}(\mathcal{V}(I)) = \sqrt{I}.
\]
This implies that the composition $\mathcal{I} \circ \mathcal{V}$ is the
identity on the set of radical ideals; that is, $\mathcal{V}$ has a
left-inverse, and therefore it is injective (cf.\ Proposition~I.2.1!).
\end{prop}
\begin{proof}
Note that $\mathcal{I}(\mathcal{V}(I))$ is a radical ideal (by
\cref{lem:fields:2.15}). It follows immediately from the definitions that
$I \subseteq \mathcal{I}(\mathcal{V}(I))$; therefore
\[
  \sqrt{I} \subseteq \mathcal{I}(\mathcal{V}(I))
  = \mathcal{I}(\mathcal{V}(I)).
\]
We have to verify the reverse inclusion: assume that $f \in
\mathcal{I}(\mathcal{V}(I))$, that is, $f(p) = 0$ for all $p \in
\mathbb{A}^n_K$ such that $\forall g \in I,\, g(p) = 0$; we have to show
that there exists an $m$ for which $f^m \in I$.
The argument is not difficult after the fact, but it is fiendishly
clever.\footnote{This is called the \name{Rabinowitsch} trick and dates
back to 1929. It was published in a one-page, eighteen-line article in the
\textit{Mathematische Annalen}.} Let $y$ be an extra variable, and
consider the ideal $J$ generated by $I$ and $1 - fy$ in
$K[x_1, \dots, x_n, y]$. More explicitly, assume
$I = (g_1, \dots, g_r)$
(since $K[x_1, \dots, x_n]$ is Noetherian, finitely many generators
suffice); we can view $g_i(x_1, \dots, x_n)$, resp.,
$f(x_1, \dots, x_n) \in K[x_1, \dots, x_n]$, as polynomials
$G_i(x_1, \dots, x_n, y)$, resp., $F(x_1, \dots, x_n, y) \in
K[x_1, \dots, x_n, y]$, and let
\[
  J = (G_1, \dots, G_r, 1 - Fy).
\]
As $K[x_1, \dots, x_n, y]$ has $n+1$ variables, the ideal $J$ defines a
subset $\mathcal{V}(J)$ in $\mathbb{A}^{n+1}_K$. I claim $\mathcal{V}(J)
= \emptyset$.

Indeed, assume on the contrary that $\mathcal{V}(J) \neq \emptyset$, and
let $p = (a_1, \dots, a_n, b) \in \mathbb{A}^{n+1}_K$ be a point of
$\mathcal{V}(J)$. Then for $i = 1, \dots, r$, $G_i(a_1, \dots, a_n, b)
= 0$; but this means $g_i(a_1, \dots, a_n) = 0$ for $i = 1, \dots, r$;
that is, $(a_1, \dots, a_n) \in \mathcal{V}(I)$. Since $f$ vanishes at
all points of $\mathcal{V}(I)$, this implies $f(a_1, \dots, a_n) = 0$,
and then
\[
  (1 - Fy)(a_1, \dots, a_n, b)
  = 1 - f(a_1, \dots, a_n)\,b = 1 - 0 = 1.
\]
But this contradicts the assumption that $p \in \mathcal{V}(J)$.
Therefore, $\mathcal{V}(J) = \emptyset$.
Now use the weak Nullstellensatz: since $K$ is algebraically closed, we
can conclude $J = (1)$. Therefore, there exist polynomials
$H_i(x_1, \dots, x_n, y)$, $i = 1, \dots, r$, and
$L(x_1, \dots, x_n, y)$, such that
\[
  \sum_{i=1}^{r} H_i(\mathbf{x}, y)\,G_i(\mathbf{x}, y)
  + L(\mathbf{x}, y)(1 - F(\mathbf{x}, y)\,y) = 1.
\]
We are not done being clever: the next step is to view this equality in
the field of rational functions over $K$ rather than in the polynomial
ring, which we can do, since the latter is a subring of the former. We
can then plug in $y = 1/F$ and still get an equality, with the advantage
of killing the last summand on the left-hand side:
\[
  \sum_{i=1}^{r} H_i\!\left(\mathbf{x}, \tfrac{1}{F}\right)
  G_i\!\left(\mathbf{x}, \tfrac{1}{F}\right) = 1.
\]
Further, $G_i(\mathbf{x}, 1/F) = g_i(x_1, \dots, x_n)$ ($G_i$ is just
the name of $g_i$ in the larger polynomial ring and only depends on the
first $n$ variables), while $H_i(\mathbf{x}, 1/F) =
h_i(x_1, \dots, x_n)/f(x_1, \dots, x_n)^m$, where $h_i \in
K[x_1, \dots, x_n]$ and $m$ is an integer large enough to work for all
$i = 1, \dots, r$. Expressing in terms of a common denominator, we can
rewrite the identity as
\[
  \frac{h_1 g_1 + \dots + h_r g_r}{f^m} = 1,
\]
or $f^m = h_1 g_1 + \dots + h_r g_r$. We have proved this identity in
the field of rational functions; as it only involves polynomials, it holds
in $K[x_1, \dots, x_n]$. The right-hand side is an element of $I$, so
this proves $f^m \in I$, and we are done.
\end{proof}
The addition of the variable $y$ in the proof looks more reasonable (and
less like a trick) if one views it geometrically. The variety
$\mathcal{V}(1 - xy)$ in $\mathbb{A}^2_K$ is an ordinary hyperbola. The
projection $(x, y) \mapsto x$ may be used to identify $\mathcal{V}(1 -
xy)$ with the complement of the origin (that is, $\mathcal{V}(x)$) in the
``$x$-axis'' (viewed as $\mathbb{A}^1_K$). Similarly, $\mathcal{V}(1 -
fy)$ is a way to realize the Zariski open set $\mathbb{A}^n_K \setminus
\mathcal{V}(f)$ as a Zariski closed subset in a space $\mathbb{A}^{n+1}_K$
of dimension one higher. The proof of \cref{prop:fields:2.17} hinges on
this fact, and this fact is also an important building block in the basic
set-up of algebraic geometry, since it can be used to show that the
Zariski topology has a basis consisting of (sets which are isomorphic in
a suitable sense to) affine algebraic sets.
I will officially record the conclusion we were able to establish
concerning the functions $\mathcal{V}$, $\mathcal{I}$:
\begin{coro}{}{fields:2.18}
Let $K$ be an algebraically closed field. Then for any $n \geq 0$ the
functions
\[
  \left\{\text{algebraic subsets of }\mathbb{A}^n_K\right\}
  \begin{matrix}
    \xrightarrow{\;\mathcal{I}\;} \\[3pt]
    \xleftarrow{\;\mathcal{V}\;}
  \end{matrix}
  \left\{\text{radical ideals in }K[x_1,\dots,x_n]\right\}
\]
are inverses of each other.
\end{coro}
\begin{proof}
\cref{prop:fields:2.17} shows that $\mathcal{I} \circ \mathcal{V}$ is the
identity on radical ideals, so $\mathcal{V}$ is injective, and
$\mathcal{V}$ is surjective by definition of affine algebraic set. It
follows that $\mathcal{V}$ is a bijection and $\mathcal{I}$ is its
inverse.
\end{proof}
Summarizing: If $K$ is algebraically closed, then studying sets defined by
polynomial equations in a space $K^n$ is `the same thing as' studying
radical ideals in a polynomial ring $K[x_1, \dots, x_n]$.
This correspondence is actually realized even more effectively at the
level of $K$-algebras. We say that a ring is \emph{reduced} if it has no
nonzero nilpotents; then $R/I$ is reduced if and only if $I$ is a radical
ideal (Exercise~2.8).
\begin{defn}{}{fields:2.19}
Let $S \subseteq \mathbb{A}^n_K$ be an algebraic set. The
\emph{coordinate ring} of $S$ is the quotient\footnote{The notation
`$K[S]$' is fairly standard, although it may lead to confusion with the
usual polynomial rings; hopefully the context takes care of this. It makes
sense to incorporate the name of the field in the notation, since one
could in principle change the base field while keeping the `same' equations
encoded in the ideal $\mathcal{I}(S)$.}
\[
  K[S] := \frac{K[x_1, \dots, x_n]}{\mathcal{I}(S)}.
\]
\end{defn}
Thus, the coordinate ring of an algebraic set is a reduced, commutative
$K$-algebra of finite type. The reader will establish a `concrete'
interpretation of this ring, as the ring of `polynomial functions' on $S$,
in Exercise~2.12. One way to think about the Nullstellensatz (in its
manifestation as \cref{coro:fields:2.18}) is that, if $K$ is algebraically
closed, then every reduced commutative $K$-algebra of finite type $R$ may
be realized as the coordinate ring of an affine algebraic set $S$; the
points of $S$ correspond in a natural way to the maximal ideals of $R$
(Exercise~2.14).
In fact, in basic algebraic geometry one defines the category of affine
algebraic sets over a field $K$ in terms of this correspondence: morphisms
of algebraic sets are defined so that they match homomorphisms of the
corresponding (reduced, finite-type) $K$-algebras. (The reader will
encounter a more precise definition in Example~VIII.1.9.)
This dictionary allows us to translate every `geometric' feature of an
algebraic set (such as dimension, smoothness, etc.) into corresponding
`algebraic' features of the corresponding coordinate rings. Once these key
algebraic features have been identified, then one can throw away the
restriction of working on reduced finite-type algebras over an algebraically
closed field and try to `do geometry' on any (say commutative, Noetherian)
ring. For example, it turns out that PIDs correspond to (certain) smooth
curves in the affine geometric world; and then one can try to think of $\Z$
as a `smooth curve' (although $\Z$ is not a reduced finite-type $K$-algebra
over any field $K$!) and try to use theorems inspired by the geometry of
curves to understand features of $\Z$---that is, use geometry to do number
theory.\footnote{This is of course a drastic oversimplification.}
Another direction in which these simple considerations may be generalized is
by `gluing' affine algebraic sets into manifold-like objects: sets which may
not be affine algebraic sets `globally' but may be covered by affine
algebraic sets. A concrete way to do this is by working in projective space
rather than affine space; the globalization process can then be carried out
in a straightforward way at the algebraic level, where it translates into
the study of `graded' rings and modules---more notions for which,
regretfully, I will find no room in this book (except for a glance, in
\S{}VIII.4.3). In the second half of the twentieth century a more abstract
viewpoint, championed by \name{Alexander Grothendieck} and others, has
proven to be extremely effective; it has led to the intense study of
schemes, the current language of choice in algebraic geometry. We have
encountered the simplest kind of schemes when we introduced the spectrum of
a ring $R$, $\operatorname{Spec} R$, back in \S{}III.4.3.

\subsection*{Exercises}

\begin{exer}{}{fields:exer:2.1}
  Prove \cref{lem:fields:2.1}.
\end{exer}
\begin{exer}{}{fields:exer:2.2}
  Let $k \subseteq \bar{k}$ be an algebraic closure, and let $L$ be an
  intermediate field. Assume that every polynomial $f(x) \in k[x] \subseteq
  L[x]$ factors as a product of linear terms in $L[x]$. Prove that
  $L = \bar{k}$.
\end{exer}
\begin{exer}{}{fields:exer:2.3}
  Prove that if $k$ is a countable field, then so is $\bar{k}$.
\end{exer}
\begin{exer}{}{fields:exer:2.4}
  Let $k$ be a field, let $c_1, \dots, c_m \in k$ be distinct elements,
  and let $\lambda_1, \dots, \lambda_m$ be nonzero elements of $k$. Prove
  that
  \[
    \frac{\lambda_1}{x - c_1} + \cdots + \frac{\lambda_m}{x - c_m} \neq 0
  \]
  in $k(x)$. (This fact is used in the proof of \cref{thm:fields:2.9}.)
\end{exer}
\begin{exer}{}{fields:exer:2.5}
  Let $K$ be a field, let $A$ be a subset of $K[x_1, \dots, x_n]$, and
  let $I$ be the ideal generated by $A$. Prove that $\mathcal{V}(A) =
  \mathcal{V}(I)$ in $\mathbb{A}^n_K$.
\end{exer}
\begin{exer}{}{fields:exer:2.6}
  Let $K$ be your favorite infinite field. Find examples of subsets $S
  \subseteq \mathbb{A}^n_K$ which cannot be realized as $\mathcal{V}(I)$
  for any ideal $I \subseteq K[x_1, \dots, x_n]$. Prove that if $K$ is a
  finite field, then every subset $S \subseteq \mathbb{A}^n_K$ equals
  $\mathcal{V}(I)$ for some ideal $I \subseteq K[x_1, \dots, x_n]$.
\end{exer}
\begin{exer}{}{fields:exer:2.7}
  Let $K$ be a field and $n$ a nonnegative integer. Prove that the set of
  algebraic subsets of $\mathbb{A}^n_K$ is the family of closed sets of a
  topology on $\mathbb{A}^n_K$.
\end{exer}
\begin{exer}{}{fields:exer:2.8}
  With notation as in \cref{defn:fields:2.13}:
  \begin{itemize}
    \item Prove that the set $\sqrt{I}$ is an ideal of $R$.
    \item Prove that $\sqrt{I}$ corresponds to the nilradical of $R/I$ via
      the correspondence between ideals of $R/I$ and ideals of $R$
      containing $I$.
    \item Prove that $\sqrt{I}$ is in fact the intersection of all prime
      ideals of $R$ containing $I$. (Cf.\ Exercise~V.3.13.)
    \item Prove that $I$ is radical if and only if $R/I$ is reduced.
      (Cf.\ Exercise~III.3.13.)
  \end{itemize}
\end{exer}
\begin{exer}{}{fields:exer:2.9}
  Prove that every affine algebraic set equals $\mathcal{V}(I)$ for a
  radical ideal $I$.
\end{exer}
\begin{exer}{}{fields:exer:2.10}
  Prove that every ideal in a Noetherian ring contains a power of its
  radical.
\end{exer}
\begin{exer}{}{fields:exer:2.11}
  Assume a field $K$ is not algebraically closed. Find a reduced
  finite-type $K$-algebra which is not the coordinate ring of any affine
  algebraic set.
\end{exer}
\begin{exer}{}{fields:exer:2.12}
  Let $K$ be an infinite field. A \emph{polynomial function} on an affine
  algebraic set $S \subseteq \mathbb{A}^n_K$ is the restriction to $S$ of
  (the evaluation function of) a polynomial $f(x_1, \dots, x_n) \in
  K[x_1, \dots, x_n]$. Polynomial functions on $S$ manifestly form a ring
  and in fact a $K$-algebra. Prove that this $K$-algebra is isomorphic to
  the coordinate ring of $S$.
\end{exer}
\begin{exer}{}{fields:exer:2.13}
  Let $K$ be an algebraically closed field. Prove that every reduced
  commutative $K$-algebra of finite type is the coordinate ring of an
  algebraic set $S$ in some affine space $\mathbb{A}^n_K$.
\end{exer}
\begin{exer}{}{fields:exer:2.14}
  Prove that, over an algebraically closed field $K$, the points of an
  algebraic set $S$ correspond to the maximal ideals of the coordinate ring
  $K[S]$ of $S$, in such a way that if $p$ corresponds to the maximal ideal
  $\mathfrak{m}_p$, then the value of the function $f \in K[S]$ at $p$
  equals the coset of $f$ in $K[S]/\mathfrak{m}_p \cong K$.
\end{exer}
\begin{exer}{}{fields:exer:2.15}
  Let $K$ be an algebraically closed field. An algebraic subset $S$ of
  $\mathbb{A}^n_K$ is \emph{irreducible} if it cannot be written as the
  union of two algebraic subsets properly contained in it. Prove that $S$
  is irreducible if and only if its ideal $\mathcal{I}(S)$ is prime, if and
  only if its coordinate ring $K[S]$ is an integral domain.

  An irreducible algebraic set is `all in one piece', like $\mathbb{A}^n_K$
  itself, and unlike (for example) $\mathcal{V}(xy)$ in the affine plane
  $\mathbb{A}^2_K$ with coordinates $x$, $y$. Irreducible affine algebraic
  sets are called \emph{(affine algebraic) varieties}.
\end{exer}
\begin{exer}{}{fields:exer:2.16}
  Let $K$ be an algebraically closed field. The field of rational functions
  $K(x_1, \dots, x_n)$ is the field of fractions of $K[\mathbb{A}^n_K] =
  K[x_1, \dots, x_n]$; every rational function $\alpha = F/G$ (with
  $G \neq 0$ and $F$, $G$ relatively prime) may be viewed as defining a
  function on the open set $\mathbb{A}^n_K \setminus \mathcal{V}(G)$; we
  say that $\alpha$ is `defined' for all points in the complement of
  $\mathcal{V}(G)$.

  Let $G \in K[x_1, \dots, x_n]$ be irreducible. The set of rational
  functions that are defined in the complement of $\mathcal{V}(G)$ is a
  subring of $K(x_1, \dots, x_n)$. Prove that this subring may be
  identified with the localization (Exercise~V.4.7) of $K[\mathbb{A}^n_K]$
  at the multiplicative set $\{1, G, G^2, G^3, \dots\}$. (Use the
  Nullstellensatz.)

  The same considerations may be carried out for any irreducible algebraic
  set $S$, adopting as field of `rational functions' $K(S)$ the field of
  fractions of the integral domain $K[S]$.
\end{exer}
\begin{exer}{}{fields:exer:2.17}
  Let $K$ be an algebraically closed field, and let $\mathfrak{m}$ be a
  maximal ideal of $K[x_1, \dots, x_n]$, corresponding to a point $p$ of
  $\mathbb{A}^n_K$. A \emph{germ of a function} at $p$ is determined by an
  open set containing $p$ and a function defined on that open set; in our
  context (dealing with rational functions and where the open set may be
  taken to be the complement of a function that does not vanish at $p$)
  this is the same information as a rational function defined at $p$, in
  the sense of Exercise~2.16. Show how to identify the ring of germs with
  the localization $K[\mathbb{A}^n_K]_{\mathfrak{m}}$ (defined in
  Exercise~V.4.11).

  As in Exercise~2.16, the same discussion can be carried out for any
  algebraic set. This is the origin of the name `localization': localizing
  the coordinate ring of a variety $V$ at the maximal ideal corresponding
  to a point $p$ amounts to considering only functions defined in a
  neighborhood of $p$, thus studying $V$ `locally', `near $p$'.
\end{exer}
\begin{exer}{}{fields:exer:2.18}
  Let $K$ be an algebraically closed field. Consider the two `curves'
  $C_1 : y = x^2$, $C_2 : y^2 = x^3$ in $\mathbb{A}^2_K$ (pictures of
  the real points of these algebraic sets are shown in
  \cref{exam:fields:exam:2.12}).
  \begin{itemize}
    \item Prove that $K[C_1] \cong K[t] = K[\mathbb{A}^1_K]$, while
      $K[C_2]$ may be identified with the subring $K[t^2, t^3]$ of $K[t]$
      consisting of polynomials $a_0 + a_2 t^2 + \dots + a_d t^d$ with
      zero $t$-coefficient. (Note that every polynomial in $K[x, y]$ may
      be written as $f(x) + g(x)y + h(x,y)(y^2 - x^3)$ for uniquely
      determined polynomials $f(x)$, $g(x)$, $h(x,y)$.)
    \item Show that $C_1$, $C_2$ are both irreducible
      (cf.\ Exercise~2.15).
    \item Prove that $K[C_1]$ is a UFD, while $K[C_2]$ is not.
    \item Show that the Krull dimension of both $K[C_1]$ and $K[C_2]$ is
      $1$. (This is why these sets would be called `curves'. You may use
      the fact that maximal chains of prime ideals in $K[x,y]$ have
      length~$2$.)
    \item The origin $(0,0)$ is in both $C_1$, $C_2$ and corresponds to
      the maximal ideals $\mathfrak{m}_1$, resp., $\mathfrak{m}_2$, in
      $K[C_1]$, resp., $K[C_2]$, generated by the classes of $x$ and $y$.
    \item Prove that the localization $K[C_1]_{\mathfrak{m}_1}$ is a DVR
      (Exercise~V.4.13). Prove that the localization
      $K[C_2]_{\mathfrak{m}_2}$ is not a DVR. (Note that the relation
      $y^2 = x^3$ still holds in this ring; prove that
      $K[C_2]_{\mathfrak{m}_2}$ is not a UFD.)
  \end{itemize}
  The fact that a DVR admits a local parameter, that is, a single
  generator for its maximal ideal (cf.\ Exercise~V.2.20), is a good
  algebraic translation of the fact that a curve such as $C_1$ has a
  single, smooth branch through $(0,0)$. The maximal ideal of
  $K[C_2]_{\mathfrak{m}_2}$ cannot be generated by just one element, as
  the reader may verify.
\end{exer}
\begin{exer}{}{fields:exer:2.19}
  Prove that the fields of rational functions (Exercise~2.16) of the
  curves $C_1$ and $C_2$ of Exercise~2.18 are isomorphic and both have
  transcendence degree $1$ over $k$ (cf.\ Exercise~1.27).

  This is another reason why we should think of $C_1$ and $C_2$ as
  `curves'. In fact, it can be proven that the Krull dimension of the
  coordinate ring of a variety equals the transcendence degree of its
  field of rational functions. This is a consequence of
  \name{Noether}'s normalization theorem, a cornerstone of commutative
  algebra.
\end{exer}
\begin{exer}{}{fields:exer:2.20}
  Recall from Exercise~VI.2.13 that $\mathbb{P}^n_K$ denotes the
  `projective space' parametrizing lines in the vector space $K^{n+1}$.
  Every such line consists of multiples of a nonzero vector $(c_0, \dots,
  c_n) \in K^{n+1}$, so that $\mathbb{P}^n_K$ may be identified with the
  quotient (in the set-theoretic sense of \S{}I.1.5) of
  $K^{n+1} \setminus \{(0, \dots, 0)\}$ by the equivalence relation $\sim$
  defined by
  \[
    (c_0, \dots, c_n) \sim (c_0', \dots, c_n')
    \iff (\exists\, \lambda \in K^*)\;
    (c_0', \dots, c_n') = (\lambda c_0, \dots, \lambda c_n).
  \]
  The `point' in $\mathbb{P}^n_K$ determined by the vector $(c_0, \dots,
  c_n)$ is denoted $(c_0 : \cdots : c_n)$; these are the `projective
  coordinates'\footnote{This is a convenient abuse of language. Keep in
  mind that the $c_i$'s are not determined by the point, so they are not
  `coordinates' in any strict sense. Their ratios are, however, determined
  by the point.} of the point. Note that there is no `point'
  $(0 : \cdots : 0)$.

  Prove that the function $\mathbb{A}^n_K \to \mathbb{P}^n_K$ defined by
  \[
    (c_1, \dots, c_n) \mapsto (1 : c_1 : \cdots : c_n)
  \]
  is a bijection. This function is used to realize $\mathbb{A}^n_K$ as a
  subset of $\mathbb{P}^n_K$. By using similar functions, prove that
  $\mathbb{P}^n_K$ can be covered with $n+1$ copies of $\mathbb{A}^n_K$,
  and relate this fact to the cell decomposition obtained in
  Exercise~VI.2.13. (Suggestion: Work out carefully the case $n = 2$.)
\end{exer}

\begin{exer}{}{fields:exer:2.21}
  Let $F(x_0, \dots, x_n) \in K[x_0, \dots, x_n]$ be a homogeneous
  polynomial. With notation as in Exercise~2.20, prove that the condition
  `$F(c_0, \dots, c_n) = 0$' for a point $(c_0 : \cdots : c_n) \in
  \mathbb{P}^n_K$ is well-defined: it does not depend on the
  representative $(c_0, \dots, c_n)$ chosen for the point $(c_0 : \cdots :
  c_n)$. We can then define the following subset of $\mathbb{P}^n_K$:
  \[
    \mathcal{V}(F) :=
    \{(c_0 : \cdots : c_n) \in \mathbb{P}^n_K \mid F(c_0,\dots,c_n)=0\}.
  \]
  Prove that this `projective algebraic set' can be covered with $n+1$
  affine algebraic sets.

  The basic definitions in `projective algebraic geometry' can be
  developed along essentially the same path taken in this section for
  affine algebraic geometry, using `homogeneous ideals' (that is, ideals
  generated by homogeneous polynomials; see \S{}VIII.4.3) rather than
  ordinary ideals. This problem shows one way to relate projective and
  affine algebraic sets, in one template example.
\end{exer}

\section{Geometric impossibilities}
\label{sec:fields:geometric-impossibilities}

% Very simple considerations in field theory dispose easily of a whole class of
% geometric
% problems which had seemed unapproachable for a very long time---problems which
% preoccupied the Greeks and were only solved in the relatively recent past.
% These problems have to do with the construction of certain geometric fig-
% ures with `straightedge and compass', that is, subject to certain very strict
% rules.
% The practical utility of these constructions would appear to be absolute zero,
% but
% the intellectual satisfaction of being able to thoroughly understand matters
% which
% stumped extremely intelligent people for a very long time is well worth the
% little
% side trip.

\subsection{Constructions by straightedge and compass}
\label{subsec:fields:constructions-by-straightedge-and-compass}
% We begin with two points

% O, P in the ordinary, real plane. You are allowed to mark (`construct') more
% points
% and other geometric figures in the plane, but only according to the following
% rules:
% \bullet  If you have constructed two points A, B, then you can draw the line joining
% them (using your straightedge).
% \bullet  If you have constructed two points A, B, then you can draw the circle with
% center at A and containing B (using your compass).
% \bullet  You can mark any number of points of intersection of any two distinct lines,
% line and circle, or circles that you have drawn already.
% Performing these actions leads to complex (and beautiful) collections of lines
% and
% circles; we say that we have `constructed' a geometric figure if that figure
% appears
% as a subset of the whole picture.
% For example, here is a recipe to construct an equilateral triangle with O, P as
% two of its vertices:
% (1) draw the circle with center O, through P ;
% (2) draw the circle with center P , through O;
% (3) let Q be any of the two points of intersection of the two circles;
% (4) draw the lines through O, P ; O, Q; and P , Q.
% Of course O, P , Q are vertices of an equilateral triangle.
% It is a pleasant exercise to try to come up with a specific sequence of basic
% moves
% constructing a given figure or geometric configuration. The enterprising reader
% could try to produce the vertices of a pentagon by a straightedge-and-compass
% construction, without looking it up and before we learn enough to thwart pure
% geometric intuition.
% The general question is to decide whether a figure can or cannot be constructed
% this way. Three problems of this kind became famous in antiquity:
% \bullet  trisecting angles;
% \bullet  squaring circles;
% \bullet  doubling cubes.
% For example, one assumes one has already constructed two lines forming an angle
% \theta
% and asks whether one can construct two lines forming \theta /3. This is
% trivially possible
% for some constructible \theta : we just constructed an angle of \pi /3, so
% evidently we were
% able to trisect the angle \pi ; but can this be done for all constructible
% angles?
% The answer is no. Similarly, it is not possible to construct a square whose area
% equals the area of a (constructible) circle, and it is not possible to construct
% the
% side of a cube whose volume is twice as large as a cube with given
% (constructible)
% side.
% Field theory will allow us to establish all this13 . Later we will return to
% these
% constructions and study the question of the constructibility of regular polygons
% 13 However, for the impossibility of squaring circles we will use the
% transcendence of \pi , which
% we are taking on faith.
% with a given side: it is very easy to construct equilateral triangles (we just
% did),
% squares, hexagons; it is not too hard to construct pentagons; but the question
% of which polygons are constructible and which are not connects beautifully with
% deeper questions in field theory.
% The three problems listed above have an illustrious history; they apparently
% date back to at least 414 B.C., if one is to take this fragment from
% Aristophanes'
% ``The birds'' as an indication:
% METON: By measuring with a straightedge there
% the circle becomes a square with a court within.14
% The problems depend on the precise restrictions imposed on the constructions.
% For example, one can trisect angles if one is allowed to mark points on the
% straight-
% edge (thus making it a ruler ); in fact, this was already known to Archimedes
% (Ex-
% ercise 3.13).
% Translating these geometric problems into algebra requires establishing the fea-
% sibility of a few basic constructions.
% First of all, one can construct a line containing a given point A and perpendic-
% ular to a given line (thus, contains at least another constructible point B).
% For
% this, draw the circle with center A and containing B;
% ---if this circle only intersects
%  at B, then the line through A and B is perpen-
% dicular to :
% A
% B
% ---otherwise, the circle intersects
%  at a second point C (whether A is on
% or not):
% 14 I thank Matilde Marcolli for the translation. For those who are less ignorant
% than I am:
% o ' \rho \theta \tilde{\omega }\iota  \mu \varepsilon \tau \rho \eta \sigma
% \omega  \kappa \alpha \nu o\nu \iota  \pi \rho o\sigma \tau \iota \theta
% \varepsilon \iota \varsigma  ` \iota \nu \alpha
% ,
% o ` \kappa \upsilon \kappa \lambda o\varsigma  \gamma \varepsilon \nu \eta \tau
% \alpha \iota  \sigma o\iota  \tau \varepsilon \tau \rho \alpha \gamma \omega \nu
% o\varsigma  \kappa \alpha \nu  \mu \varepsilon \sigma \omega \iota  \alpha
% \gamma o\rho \alpha
% and once C is determined, the line through A and perpendicular to may be found
% by joining the two points of intersection D, E of the two circles with centers
% at B,
% resp., C, and containing C, resp., B:
% Applying this construction twice gives the line through a point A and parallel
% to a given line:
% A
% which is often handy.
% Judicious application of these operations allows us to bring a Cartesian refer-
% ence system into the picture. Starting with the initial two points O, P , we can
% construct perpendicular Cartesian axes centered at O, with P marking (say) the
% point (1, 0) on the `x-axis':
% and constructing a point A = (x, y) is equivalent to constructing its
% projections
% X = (x, 0), Y = (0, y) onto the axes, and in fact it is equivalent to
% constructing the
% \setminus two points X = (x, 0) and Y = (y, 0):
% It follows that determining which figures are constructible by straightedge and
% compass is equivalent to determining which numbers may be realized as
% coordinates
% of a constructible point.
% Definition 3.1. A real number r is constructible if the point (r, 0) is
% constructible
% with straightedge and compass (assuming O = (0, 0) and P = (1, 0), as above). I
% will denote by CR \subseteq  R the set of constructible real numbers.
%  ?
% Also, we can identify the real plane with C, placing O at 0 and P at 1, and we
% say that z = x+iy is constructible if the point (x, y) is constructible by
% straightedge
% and compass. I will denote by CC \subseteq  C the set of constructible complex
% numbers.
% Summarizing the foregoing discussion, we have proved:
% Lemma 3.2. A point (x, y) is constructible by straightedge and compass if and
% only
% if x + iy \in  CC , if and only if x, y \in  CR .
% The next obvious question is what kind of structure these sets of constructible
% numbers carry, and this prompts us to look at a few more basic straightedge-and-
% compass constructions.
% Lemma 3.3. The subset CR \subseteq  R of constructible numbers is a subfield of
% R.
% Likewise, CC is a subfield of C, and in fact CC = CR (i).
% Proof. The set CR \subseteq  R is nonempty, so in order to show it is a field,
% we only
% need to show that it is closed with respect to subtraction and division by a
% nonzero
% constructible number (cf. Proposition II.6.2).
% The reader will check that CR is closed under subtraction (Exercise 3.2). To
% see that CR is closed under division, let a, b \in  CR, with b = 0. Since A =
% (0, a) and
% B = (b, 0) are constructible by hypothesis, we can construct C as the
% y-intersect
% of the line through P = (1, 0) and parallel to the line through A and B, as in
% the
% following picture15 :
% 15 In the picture I am assuming a > 0 and b > 1; the reader will check that
% other possibilities
% lead to the same conclusion.
% Since the triangles AOB and COP are similar, we see that C = (0, a/b); it
% follows
% that a/b is constructible.
% The fact that CC is a subfield of C is an immediate consequence of the fact that
% CR is a field, and the proof of this is left to the enjoyment of the reader
% (Exercise 3.7).
% The fact that CC = CR (i) is a restatement of the fact that x + iy \in  CC if
% and only
% if x and y are in CR.
%  ?
% We may view CR and CC as extensions of Q, drawing a bridge between con-
% structibility by straightedge and compass and field theory: we will be able to
% under-
% stand constructibility of geometric figures if we can understand the field
% extensions
% Q \subseteq  CR \subseteq  CC.
% Happily, we can understand these extensions!

\subsection{Constructible numbers and quadratic extensions}
\label{subsec:fields:constructible-numbers-and-quadratic-extensions}
% Our goal is to prove

% the following amazingly explicit description of CR (immediately implying one for
% CC).
% Theorem 3.4. Let \gamma  \in  R. Then \gamma  \in  CR if and only if there exist
% real numbers
% \delta 1 , . . . , \delta k such that \forall j = 1, . . . , k
% [Q(\delta 1, . . . , \delta j ) : Q(\delta 1 , . . . , \delta j-1 )] = 2
% and \gamma  \in  Q(\delta 1, . . . , \delta k ).
% In other words, \gamma  \in  R is constructible if and only if it can be placed
% in the
% top field of a sequence of (real) quadratic extensions over Q. Since CC = CR
% (i),
% the same statement holds for constructible complex numbers, including i in the
% list
% of \delta j (or simply allowing the \delta j to be complex numbers; cf. Exercise
% 3.9).
% The proof of this theorem is as explicit as one can possibly hope. From a given
% straightedge-and-compass construction of a point (x, y) one can obtain an
% explicit
% sequence \delta 1 , . . . , \delta k as in the statement, such that x and y
% belong to the extension
% Q(\delta 1 , . . . , \delta k ); conversely, from any element \gamma  of any
% such extension one can obtain
% an explicit construction of a point (\gamma , 0) by straightedge and compass.
% Proof. Let's first argue in the `geometry to algebra' direction. A configuration
% of points, lines, and circles obtained by a straightedge-and-compass
% construction
% may be described by the coordinates of the points and the equations of the lines
% and circles. Suppose that at one stage in a given construction all coordinates
% of all
% points and all coefficients in the equations of lines and circles belong to a
% field F ;
% I will say that the configuration is defined over F .
% Then I claim that for every object constructed at the next stage, there exists
% a number \delta  \in  R, of degree at most 2 over F , such that the new
% configuration is
% defined over F (\delta ). The `only if' part of the theorem follows by induction
% on the
% number of steps in the construction, since at the beginning the configuration
% (that
% is, the pair of points O = (0, 0), P = (1, 0)) is defined over Q.
% Verifying my claim amounts to verifying it for the basic operations defining
% straightedge-and-compass constructions. The reader will check (Exercise 3.5)
% that
% the point of intersection of two lines defined over F has coordinates in F and
% that
% lines and circles determined by points with coordinates in F are defined over F
% . So
% \delta  = 1 works in all these cases.
% For the intersection of a line and a circle C, assume that is not parallel to
% the y-axis (the argument is entirely analogous otherwise) and that it does meet
% C;
% let
% y = mx + r
% be the equation of and let
% x2 + y 2 + ax + by + c = 0
% be the equation of C. We are assuming that a, b, c, m, r \in  F . Then the x-
% coordinates of the points of intersection of and C are the solutions of the
% equation
% x2 + (mx + r)2 + ax + b(mx + r) + c = 0.
% \sqrt 
% The `quadratic formula' shows that these coordinates belong to the field F ( D),
% where D is the discriminant of this polynomial: explicitly,
% D = (2mr + bm + a)2 - 4(m2 + 1)(r2 + br + c),
% \sqrt 
% but this is unimportant. What is important is that D \in  F ; hence \delta  =
%  D satisfies
% our requirement.
% For the intersection of two (distinct) circles defined over F , nothing new is
% needed: if
%  ? 2
% x + y2 + a1x + b1 y + c1 = 0,
% x2 + y2 + a2x + b2 y + c2 = 0
% are two circles, subtracting the two equations shows that their points of
% intersection
% coincide with the points of intersection of a circle and a line:
% ?x2
%  + y2 + a1 x + b1 y + c1 = 0,
% (a1 - a2)x + (b1 - b2 )y + (c1 - c2) = 0,
% with the same conclusion as in the previous case.
% This completes the verification of the `only if' part of the theorem.
% To prove that every element of an extension as stated is constructible, again
% argue by induction: it suffices to show that if (i) \delta  \in  R, (ii) all
% elements of F are
% constructible, and (iii) r = \delta  2 \in  F , then \delta  is constructible
% (note that in order to
% construct an element of degree 2 over F , it suffices to construct the square
% root of
% the discriminant of its minimal polynomial). Therefore, all we have to show is
% that
% we can `take square roots' by a straightedge-and-compass construction. Here is
% the
% picture (if r > 1):
% If A = (r, 0) is constructible, so is B = (-r, 0); C is the midpoint of the seg-
% ment BP (midpoints are constructible, Exercise 3.1); the circle with center C
% and
% containing P intersects the positive y-axis at a point Q = (0, \delta ), and
% elementary
% geometry shows that \delta  2 = r. Therefore \delta  is constructible,
% concluding the proof of
% the theorem.
%  ?
% For example, one can construct an angle of 3\circ  : allegedly
% +
% 1 \sqrt 
%  \sqrt 
%  1 \sqrt 
%  \sqrt  \sqrt 
% cos 3\circ  = ( 3 + 1) 5 + 5 + ( 6 - 2)( 5 - 1),
% and this expression shows that
% +
% \sqrt  \sqrt  \sqrt 
%  \sqrt 
% cos 3\circ 
%  \in  Q( 2, 3, 5)( 5 + 5);
% that is (by Theorem 3.4), cos 3\circ  is constructible. Of course once A = (cos
% \theta , 0) is
% constructed, so is the angle \theta :
% Here is another example of a `constructive' application of Theorem 3.4:
% Example 3.5. Regular pentagons are constructible.
% Indeed, it suffices to construct the point A = (cos(2\pi /5), 0), and it so
% happens that
% \gamma  = cos(2\pi /5) satisfies
% (*)
%  4\gamma  2 + 2\gamma  - 1 = 0
% \sqrt 
% (in fact, \gamma  is half of the inverse of the golden ratio: \gamma  = 5-1
%  ). It follows that
% \sqrt 
%  4\gamma  \in  Q( 5), and hence
%  it is constructible by Theorem 3.4. In fact, the proof shows
% \sqrt how to construct 5, so the reader should now have no difficulty producing a
% straightedge-and-compass construction of a regular pentagon (Exercise 3.6).
%  ?
% We will come back to constructions of regular polygons once we have studied
% field theory a little more thoroughly (cf. \S{}7.2). By the way, how does one
% figure
% out that \gamma  = cos(2\pi /5) should satisfy the identity (*)? This is easy
% modulo some
% complex number arithmetic; see Exercise 3.11.

\subsection{Famous impossibilities}
\label{subsec:fields:famous-impossibilities}
% Theorem 3.4 easily settles the three problems men-

% tioned in \S{}3.1, via the following immediate consequence:
% Corollary 3.6. Let \gamma  \in  CC be a constructible number. Then [Q(\gamma ) :
% Q] is a power
% of 2.
% Proof. By Lemma 3.3 and Theorem 3.4, there exist \delta 1 . . . , \delta k \in
% R such that
% \gamma  \in  Q(\delta 1 , . . . , \delta k , i),
% and each \delta j has degree \leq  2 over Q(\delta 1, . . . , \delta j-1).
% Repeated application of Proposi-
% tion 1.10 shows that
% [Q(\delta 1 , . . . , \delta k , i) : Q]
% is a power of 2, and since
% Q \subseteq  Q(\gamma ) \subseteq  Q(\delta 1, . . . , \delta k , i),
% the statement follows from Corollary 1.11.
%  ?
% In particular, constructible real numbers must satisfy the same condition.
% Consider the problem of trisecting an angle. We know we can construct an
% angle of 60\circ  (this is a subproduct of the construction of an equilateral
% triangle).
% Constructing an angle of 20\circ  is equivalent (Exercise 3.10) to constructing
% the com-
% plex number \gamma  on the unit circle and with argument \pi /9:
% Now \gamma  9 = -1; that is, \gamma  satisfies the polynomial
% t9 + 1 = (t3 + 1)(t6 - t3 + 1).
% It does not satisfy t3 + 1; therefore its minimal polynomial over Q is a factor
% of
% t6 - t3 + 1.
% However, this polynomial is irreducible over Q (substitute t ?\to  t - 1, and
% apply
% Eisenstein's criterion), so
% [Q(\gamma ) : Q] = 6.
% By Corollary 3.6, \gamma  is not constructible. Therefore there exist
% constructible angles
% which cannot be trisected.
% Similarly,
%  \sqrt 
%  3 cubes cannot be doubled because that would imply the constructibil-
% ity of 2, which has degree 3 over Q, contradicting Corollary 3.6.
% Squaring circles amounts to constructing \pi , and \pi  is not even algebraic
% (al-
% though, as I have mentioned, the proof of this fact is not elementary), so this
% is
% also not possible.
% As another example, constructing a regular 7-gon would amount to constructing
% a 7-th (complex) root of 1, \zeta :
% By definition, \zeta  satisfies
% t7 - 1 = (t - 1)(t6 + t5 + \dots  + t + 1);
% as \zeta  = 1, \zeta  must satisfy the cyclotomic polynomial t6 + \dots  + 1.
% This is irreducible
% (Example V.5.19); hence again we find that \zeta  has degree 6 over Q, and
% Corollary 3.6
% implies that the regular 7-gon cannot be constructed with straightedge and
% compass.
% Of course `7' is not too special: if p is a positive prime integer, the
% cyclotomic
% polynomial of degree p - 1 is irreducible (Example V.5.19 again); hence
% [Q(\zeta p ) : Q] = p - 1,
% where \zeta p is the complex p-th root of 1 with argument 2\pi /p. Therefore,
% Corollary 3.6
% reveals that if p is prime, then the regular p-gon can be constructed only if p
% - 1 is
% a power of 2. This is even more restrictive than it looks at first, since if p =
% 2k + 1

\subsection*{Exercises}

% is prime, then necessarily k is itself a power of 2 (Exercise 3.15). Primes of
% the
% ?
% form 22 + 1 are called Fermat primes:
% 3 , 5 , 17 , 257 , 65537
% (the `next one', 232 + 1 = 4294967297 = 641 \cdot  6700417, is not prime; in
% fact, no one
% knows if there are any other Fermat primes, and yet some people conjecture that
% there are infinitely many).
% These considerations alone do not tell us that if p is a Fermat prime, then the
% regular p-gon can be constructed with straightedge and compass; but they do tell
% us that the next case to consider would be p = 24 + 1 = 17, and it just so
% happens
% that
% \sqrt 
%  \sqrt 
%  ?
%  \sqrt 
%  \sqrt  \sqrt 
%  \setminus 
%  \sqrt 
%  \sqrt  \setminus 
%  \sqrt 
%  \sqrt  \setminus 
%  \sqrt 
% 2\pi 
%  17 - 1 +
%  2 34 + 6 17 + 2( 17 - 1) 17 - 17 - 8 2 17 + 17 + 2 17 - 17
% cos
%  =
% as noted by Gauss at age 19. Therefore (by Theorem 3.4) the 17-gon is
% constructible
% by straightedge and compass. As I have announced already, the situation will be
% further clarified soon, allowing us to bypass heavy-duty trigonometry.
% 3.1. ? Prove that if A, B are constructible, then the midpoint of the segment AB
% is also constructible. Prove that if two lines 1 , 2 are constructible and not
% parallel,
% then the two lines bisecting the angles formed by 1 and 2 are also
% constructible.
% [\S{}3.2]
% 3.2. ? Prove that if a, b are constructible numbers, then so is a - b. [\S{}3.1]
% 3.3. Find an explicit straightedge-and-compass construction for the product of
% two
% real numbers.
% 3.4. Show how to square a triangle by straightedge and compass.
% 3.5. ? Let F be a subfield of R.
% \bullet  Let A = (xA , yA ), B = (xB , yB ) be two points in R2 , with xA , yA , xB , yB \in  F .
% Prove that the line through A, B is defined over F (that is, it admits an
% equation
% with coefficients in F ).
% \bullet  Prove that the circle with center at A and containing B is defined over F .
% \bullet  Let 1, 2 be two distinct, nonparallel lines in R2, defined over F , and let (x, y) =
% 1 \cap  2 . Prove that x, y \in  F .
% [\S{}3.2]
% 3.6. ? Devise a way to construct a regular pentagon with straightedge and
% compass.
% [\S{}3.2]
% 3.7. ? Identify the (real) plane with C, and place O, P at 0, 1 \in  C. Let CC
% be the
% set of all constructible points, viewed as a subset of C. Prove that CC is a
% subfield
% of C. [\S{}3.1]
% 3.8. For \delta  \in  C, \delta  = 0, let \theta \delta  be the argument of
% \delta  (that is, the angle formed by
% the line through 0 and \delta  with the real axis). Prove that \delta  \in  CC
% if and only if |\delta |,
% cos \theta \delta  , sin \theta \delta  are all constructible real numbers.
% 3.9. ? Let \gamma 1 , . . . , \gamma k \in  CC be constructible complex numbers,
% and let K be the
% field Q(\gamma 1 , . . . , \gamma k ) \subseteq  CC . Let \delta  be a complex
% number such that [K(\delta ) : K] = 2.
% Prove that \delta  \in  CC .
% Deduce that there are no irreducible polynomials of degree 2 over CC and that
% CC is the smallest subfield of C with this property. [\S{}3.2]
% 3.10. ? Prove that if two lines in any constructible configuration form an angle
% of \theta , then the line through O forming an angle of \theta  clockwise with
% respect to the
% line OP is constructible.
% Deduce that an angle \theta  is constructible anywhere in the plane if and only
% if
% cos \theta  is constructible. [\S{}3.3]
% 3.11. ? Verify that \gamma  = cos(2\pi /5) satisfies the relation (*) given in
% \S{}3.2. (If
% \sigma  = sin(2\pi /5), note that z = \gamma  + i\sigma  satisfies z 5 - 1.) [\S{}3.2]
% 3.12. Prove that the angles of 1\circ  and 2\circ  are not constructible. (Hint:
% Given what
% we know at this point, you only need to recall that there exist trigonometric
% formulas
% for the sum of two angles; the exact shape of these formulas is not important.)
% For
% what integers n is the angle n\circ  constructible?
% 3.13. ? Prove that \alpha  = \theta  in the following picture:
% 3This says that angles can be trisected if we allow the use of a `ruler', that
% is, a
% straightedge with markings (here, the construction works if we can mark a fixed
% distance of 1 on the ruler). Apparently, this construction was known to
% Archimedes.
% [\S{}3.1]
% 3.14. Prove that the regular 9-gon is not constructible.
% 3.15. ? Prove that if 2k + 1 is prime, then k is a power of 2. [\S{}3.3]

\section{Field extensions, II}
\label{sec:fields:field-extensions-ii}

% It is time to continue our survey of different flavors of field extensions. The
% keywords
% here are splitting fields, normal, separable.

\subsection{Splitting fields and normal extensions}
\label{subsec:fields:splitting-fields-and-normal-extensions}
% In \S{}2 we have constructed the

% algebraic closure k of any given field k: every polynomial in k[x] factors as a
% product of linear terms (that is, `splits' ) in k[x], and k is the `smallest'
% extension
% of k satisfying this property.
% Here is an analogous, but more modest, requirement: given a subset F \subseteq
% k[x]
% of polynomials, construct an extension k \subseteq  F such that every polynomial
% in F
% splits as a product of linear terms over F , and require F to be as small as
% pos-
% sible with this property. We then call F the splitting field for F . In practice
% we
% will only be interested in the case in which F is a finite collection of
% polynomi-
% als f1 (x), . . . , fr (x); then requiring that each fi (x) splits over F is
% equivalent to
% requiring that the product f 1 (x) \dots  fr (x) splits over F .
% That is, for our purposes it is not restrictive to assume that F consists of a
% single polynomial f (x) \in  k[x].
% Definition 4.1. Let k be a field, and let f (x) \in  k[x] be a polynomial of
% degree d.
% The splitting field for f (x) over k is an extension F of k such that
% ?
%  d
% f (x) = c
%  (x - \alpha i)
% i=1
% splits in F [x], and further F = k(\alpha 1 , . . . , \alpha d ) is generated
% over k by the roots of f (x)
% in F .
%  ?
% Note that it is clear that a splitting field exists: given f (x) \in  k[x], the
% subfield
% F \subseteq  k generated over k by the roots of f (x) in k satisfies the
% requirements in
% Definition 4.1. But I have written the splitting field, and this is justified by
% the
% uniqueness part of the following basic observation, which also evaluates `how
% big'
% this extension can be.
% Lemma 4.2. Let k be a field, and let f (x) \in  k[x]. Then the splitting field F
% for
% f (x) over k is unique up to isomorphism, and [F : k] \leq  (deg f )!.
% In fact, if \iota  : k\setminus  \to  k is any isomorphism of fields and g(x)
% \in  k \setminus  [x] is such that
% f (x) = \iota (g(x)), then \iota  extends to an isomorphism of any splitting
% field of g(x) over
% \setminus k to any splitting field of f (x) over k.
% Proof. We first construct explicitly a splitting field and obtain the bound on
% the
% degree mentioned in the statement. Then we prove the second part, which implies
% uniqueness up to isomorphism.
% The construction of a splitting field and the given bound on the degree are an
% easy application of our basic simple extension found in Proposition V.5.7.
% Arguing
% inductively, assume the splitting field has been constructed and the bound has
% been proved for all fields and all polynomials of degree (deg f - 1). Let q(t)
% be any
% irreducible factor of f (t) over k; then
% k[t]
% k \subseteq  F \setminus  :=
% (q(t))
% is an extension of degree deg q \leq  deg f , in which q(x) (and hence f (x))
% has a root
% (the coset \alpha  of t) and therefore a linear factor x - \alpha . The
% polynomial h(x) :=
% f (x)/(x - \alpha ) \in  F \setminus  [x] has degree (deg f - 1); therefore a
% splitting field F exists for
% h(x), and
% [F : F \setminus ] \leq  (deg f - 1)!.
% The factors of f (x) over F are (x - \alpha ) and the factors of h(x), which are
% linear; it
% follows that F is a splitting field for f (x) and
% [F : k] = [F : F \setminus ][F \setminus  : k] \leq  (deg f )(deg f - 1)! = (deg
% f )!
% as stated.
% To prove that isomorphisms \iota  : k\setminus  \to  k extend to isomorphisms of
% splitting
% fields as stated, let G be a splitting field for g(x) and consider the
% composition
% k\setminus  \to  k \subseteq  k. Since the extension k \setminus  \subseteq  G
% is algebraic, by Lemma 2.8 there exists
% a morphism of extensions \iota  : G \to  k over k \setminus  . Since G is
% generated over k \setminus  by the
% roots of g(x) in G, we have
% \iota (G) = k(\alpha 1 , . . . , \alpha d ) \subseteq  k,
% where the \alpha i's are the roots of f (x) = \iota (g(x)) in k. Therefore L :=
% \iota (G) is indepen-
% dent of the chosen isomorphism \iota  and splitting field G. Applying this
% observation
% to the given \iota  and to the identity k \to  k proves the statement (as both G
% and F
% are isomorphic to L).
%  ?
% Example 4.3. By definition, Q(i) is the splitting field of x2 + 1 over Q, and C
% is
% the splitting field for the same polynomial, over R.
%  ?
% Example 4.4. The splitting field F of x8 - 1 over Q is generated by \zeta  :=
% e2\pi i/8 :
% indeed, the roots of x8 - 1 are all the 8-th roots of 1, and all of them are
% powers
% of \zeta :
% i
% \zeta 
% \sqrt 
% In fact, \zeta  is a root of the polynomial x4 + 1, which is irreducible over Q;
% therefore
% F = Q(\zeta ) is `already' the splitting field of x4 + 1. The degree of F over Q
% is
% [Q(\zeta ) : Q] = 4,
% way less than the bounds 8!, 4! obtained in Lemma 4.2.
% To understand this splitting field (even) better, note that i = \zeta  2 is in F
% , and
% \sqrt 
%  \sqrt 
%  \sqrt 2 \sqrt 
% so is 2 = \zeta  + \zeta  7 ; thus F contains Q(i, 2). Conversely, \zeta  = 2 (1
% + i) \in  Q(i, 2).
% \sqrt 
% Therefore, the splitting
%  field of x4
%  +1 (a.k.a. the splitting field of x8
%  -1) is Q(i, 2).
% \sqrt Analyzing Q \subseteq  Q(i, 2) (as we did for the extension in Example 1.19) shows that
% its group of automorphisms over Q is Z/2Z \times  Z/2Z.
%  ?
% Example 4.5. Variation on the theme: x4 - 1. The situation changes if instead
% of x4 + 1 we consider x4 - 1: this polynomial factors over Q,
% x4 - 1 = (x - 1)(x + 1)(x2 + 1),
% and it follows that the splitting field is the same as for x2 + 1, that is, just
% Q(i). ?
% Example 4.6. Variation on the theme: x4 +2. The overoptimistic reader may now
% hope that the difference between the splitting fields of x4 +1 vs. x4 -1 over Q
% is just
% due to the fact that the first polynomial is irreducible over Q and the second
% is not.
% This example will nip any such guess in the bud. With notation as in Example
% 4.4,
% the roots of x4 + 2 are
%  \sqrt 
%  \sqrt 
%  3 \sqrt 
%  5 \sqrt 
%  72\zeta  , 2\zeta  , 2\zeta  , 2\zeta  .
% \sqrt 
%  4 \sqrt 
%  4 3 \sqrt 
%  4 5 \sqrt 
%  4 7Therefore, with K = Q( 2\zeta , 2\zeta  , 2\zeta  , 2\zeta  ) the splitting
%  field of x4 + 2,
% \sqrt 
%  \sqrt  \sqrt 
%  \sqrt 
% K \subseteq  Q(\zeta , 2) = Q(i, 2, 2) = Q(i, 2).
% \sqrt 
%  \sqrt 
%  4 \sqrt 
%  4 \sqrt 
%  4 \sqrt 
% On the
%  other hand, 2 = ( 2\zeta )3 /( 2\zeta  3) \in  K; hence i = ( 2\zeta )2 / 2 \in
%  K; hence
% \sqrt 
% 2 \sqrt 
%  4 \sqrt 
%  4 \sqrt 
%  4\zeta  = 2 (1 + i) \in  K; hence 2 = ( 2\zeta )/\zeta  \in  K. Therefore, Q(i,
%  2) \subseteq  K, and the
% \sqrt 
%  4conclusion is that the splitting field of x4 + 2 equals K = Q(i, 2). A simple
%  degree
% computation (Exercise 4.3) shows [K : Q] = 8 and in particular the splitting
% field
% of x4 + 2 is certainly not isomorphic to the splitting field of x4 + 1.
%  ?
% Examples such as these have always left me with the impression that field
% theory must be rather mysterious: innocent-looking variations on the parameters
% of a problem may cause dramatic changes in the corresponding field extensions.
% On
% the plus side, this hints that field theory can indeed be a very precise tool in
% (for
% example) the study of polynomials.
% Splitting fields will play an important role in the rest of the story. They are
% even more special than they may appear to be at first: it turns out that not
% only do
% they split the given polynomial, but they also automatically split any
% irreducible
% polynomial which dares touch them with a root. This makes splitting fields
% normal
% extensions:
% Definition 4.7. A field extension k \subseteq  F is normal if for every
% irreducible poly-
% nomial f (x) \in  k[x], f (x) has a root in F if and only if f (x) splits as a
% product of
% linear factors over F .
%  ?
% Theorem 4.8. A field extension k \subseteq  F is finite and normal if and only
% if F is the
% splitting field of some polynomial f (x) \in  k[x].
% Proof. Assume k \subseteq  F is finite and normal. Then F is finitely generated:
% F =
% k(\alpha 1 , . . . , \alpha r ) with \alpha i algebraic over k. Let pi(t) be the
% minimal polynomial of \alpha i
% over k. As F is normal over k, each pi(t) splits completely over F , and hence
% so
% does f (t) = p1 (t) \dots  pr (t). It follows that F is the splitting field of f
% (x).
% Conversely, assume that F is a splitting field for a polynomial f (x) \in  k[x],
% and
% let p(x) \in  k[x] be an irreducible polynomial, such that F contains a root
% \alpha  of p(x).
% Viewing F as a subfield of the algebraic closure k, let \beta  \in  k be any
% other root
% of p(x); we are going to show that \beta  \in  F . This will prove that F
% contains all roots
% of p(x), implying that k \subseteq  F is normal, and k \subseteq  F is finite by
% Lemma 4.2.
% By Proposition 1.5, there exists an isomorphism \iota  : k(\alpha ) \to  k(\beta
% ) extending the
% identity on k and sending \alpha  ?\to  \beta . We also consider the subfield F
% (\beta ) \subseteq  k, viewed
% as an extension of k(\beta ). Putting everything in one diagram:
% k(\alpha ) ?
%  ?
%  / F ??
%  / k
% k) ? n P u n P n
%  P
%  n
%  P n
%  P n
%  P 7
% '
%  k(\beta ) \sqcup 
%  \iota 
%  ?
%  ?
%  / F (\beta ) ? ?
%  / k
% (Note: If we join the two copies of k on the right by the identity map, the
% corre-
% sponding diagram does not commute: \iota  sends \alpha  \in  k in the top row to
% \beta  \in  k in
% the bottom row.) Now observe that F may be viewed as the splitting field of f
% (x)
% over k(\alpha ): indeed, it contains all the roots of f (x) and is generated
% over k (and
% hence over k(\alpha )) by these roots. By the same token, F (\beta ) is the
% splitting field
% for f (x) over k(\beta ). By Lemma 4.2, \iota  extends to an isomorphism \iota
% \setminus  : F \to  F (\beta ).
% Since \iota  restricts to the identity on k, \iota \setminus  is an isomorphism
% of k-vector spaces; in
% particular, dim k F = dimk F (\beta ) (keep in mind that splitting fields are
% finite exten-
% sions).
% Now consider the different k-linear map of k-vector spaces i : F \to  F (\beta )
% simply
% given by the inclusion within k:
% k \subseteq  F \subseteq  F (\beta ) \subseteq  k.
% Since F and F (\beta ) are k-vector spaces of the same finite dimension, i must
% also be
% an isomorphism. In other words,
% [F (\beta ) : F ] = 1;
% that is, \beta  \in  F as needed.
%  ?
% Remark 4.9. This argument is somewhat delicate. With notation as in the proof,
% the diagram
% k(\alpha ) ?
%  ?
%  / F
% \iota 
% \sqcup  ?
%  ?
% k(\beta )is commutative, while the diagram
% ?
%  ?
% k(\alpha )/\iota ?
% \sqcup 
% F (\beta )
% / F
% \iota 
% \sqcup  ?
%  ?
% k(\beta )/i
% \sqcup 
% F (\beta )
% is not commutative if \beta  = \alpha . Indeed, \iota  sends \alpha  to \beta ,
% while i sends \alpha  to \alpha . The
% key point in the argument is the observation that if there exists one
% isomorphism
% between two finite-dimensional vector spaces V , W , then every injective linear
% map
% V \to  W must be an isomorphism. Finite dimensionality is necessary in order to
% draw this conclusion; cf. Exercise VI.6.5.
%  ?
% Example 4.10. If a complex root of an irreducible polynomial p(x) \in  Q[x] may
% \sqrt 
%  4be expressed as a polynomial in i and 2 with rational
%  \sqrt 
%  coefficients, then all roots
% of p(x) may be expressed
%  \sqrt 
%  likewise
%  in
%  terms
%  of
%  i
%  and
%  2.
%  Indeed, we have checked
% (Example 4.6) that Q(i, 2) is a splitting field over Q; hence it is a normal
% extension
% of Q.
%  ?

\subsection{Separable polynomials}
\label{subsec:fields:separable-polynomials}
% Our intuition may lead us to think that if a poly-

% nomial factors as a product of linear factors and we are not `purposely'
% repeating
% one of the factors (as in (x - 1)2 (x - 2)), then these will be distinct. For
% example,
% surely irreducible polynomials necessarily split as products of distinct factors
% in an
% algebraic closure, right? Wrong.
% Example 4.11. Let p be a prime, and consider the field Fp (t) of rational
% functions
% over Fp . Then the polynomial
% xp - t \in  Fp(t)[x]
% is irreducible: by Eisenstein's criterion it is irreducible in Fp [t][x] (since
% (t) is prime
% in Fp [t]), hence in Fp (t)[x] by Proposition V.4.16. Let u be a root of this
% polynomial
% in an extension L of Fp (t) (for example L could be the algebraic closure of Fp
% (t),
% or more modestly a splitting field for the polynomial). Then (Exercise 4.8)
% xp - t = (x - u)p
% in L[x]; that is, u has multiplicity p as a root of f (x).
% In other words, the minimal polynomial over Fp (t) of u \in  L vanishes p times
% at u, and there is nothing to do about this: no smaller power of (x - u) than
% (x - u)p = xp - t has coefficients in Fp(t) (a smaller power would give a
% nontrivial
% factor of xp - t, and xp - t is irreducible).
%  ?
% I have always found this example difficult to visualize, because of intuition
% developed in characteristic 0, and as we will see, no such pathology can occur
% in
% characteristic 0.
% Definition 4.12. Let k be a field. A polynomial f (x) \in  k[x] is separable if
% it
% has no multiple factors over its splitting field; f (x) is inseparable if it has
% multiple
% factors over its splitting field.
%  ?
% Thus, the polynomial xp - t in Example 4.11 is inseparable. I suppose the
% terminology reflects the fact that we cannot `separate' its roots: they come
% clumped
% together like quarks in a proton, and we cannot take them apart.
% Of course Definition 4.12 does not depend on the chosen splitting field, since
% these are unique up to isomorphism (Lemma 4.2). In fact, to detect separability,
% we can use any field in which the polynomial splits as a product of linear
% factors
% (by Exercise 4.1). The first, somewhat surprising, observation about
% separability is
% that we can in fact detect it without leaving the field of coefficients of f
% (x). This
% fact uses a notion borrowed from calculus: for a polynomial
% f (x) = a0 + a1 x + a2 x2 + \dots  + an xn ,
% \setminus we denote by f (x) the `derivative'
% f \setminus  (x) = a1 + 2a2x + \dots  + nan xn-1.
% Of course this is a purely formal operation; no limiting process is at work over
% arbitrary fields. Still, all expected properties of derivatives hold, as the
% reader
% should check: for example, (f g)\setminus  = f \setminus  g + f g\setminus  as
% usual.
% Lemma 4.13. Let k be a field, and let f (x) \in  k[x]. Then f (x) is separable
% if and
% only if f (x) and f \setminus (x) are relatively prime.
% \setminus By definition, f (x) and f (x) are relatively prime precisely when the greatest
% common divisor of f (x) and f \setminus  (x) is 1. Note that if k \subseteq  F ,
% the gcd of f (x) and f \setminus  (x)
% is the same whether it is considered in k[x] or in F [x]: for example because it
% can
% be computed by applying the Euclidean algorithm (\S{}V.2.4), and this proceeds
% in
% exactly the same way whether it is performed in k[x] or in F [x].
% Proof. First assume that f (x) is not separable. Then f (x) has a multiple root
% in
% a splitting field F ; that is,
% f (x) = (x - \alpha )m g(x)
% for some \alpha  \in  F , g(x) \in  F [x], and m \geq  2. Therefore
% f \setminus  (x) = m(x - \alpha )m-1 g(x) + (x - \alpha )mg \setminus  (x),
% and it follows that (x - \alpha ) is a common factor of f (x) and f \setminus
% (x) in F [x]. Therefore
% \setminus  \setminus  \setminus gcd(f (x), f (x)) = 1 in F [x]; hence gcd(f (x), f (x)) = 1 in k[x]; that is, f (x), f (x)
% are not relatively prime.
% \setminus  \setminus Conversely, assume gcd(f (x), f (x)) = 1, so f (x) and f (x) have a common
% irreducible factor (x - \alpha ) in the algebraic closure k of k. Write f (x) =
% (x - \alpha )h(x);
% we have
% f \setminus  (x) = h(x) + (x - \alpha )h\setminus (x),
% and it follows that x - \alpha  divides h(x) since it divides f \setminus  (x).
% But then
% (x - \alpha ) 2 | f (x);
% hence f (x) is inseparable.
%  ?
% For example, the polynomial xp - t of Example 4.11 may be seen to be insepa-
% rable without invoking splitting fields: the derivative of xp - t equals pxp-1 =
% 0 in
% characteristic p, and gcd(xp - t, 0) = xp - t = 1. This example captures one of
% the
% key features of inseparability:
% Lemma 4.14. Let k be a field, and let f (x) \in  k[x] be an inseparable
% irreducible
% \setminus polynomial. Then f (x) = 0.
% \setminus Proof. Since f (x) is inseparable, f (x) and f (x) have a common irreducible fac-
% tor q(x) by Lemma 4.13; but as f (x) is itself irreducible, q(x) must be an
% associate
% of f (x), and in particular its degree is larger than the degree of f \setminus
% (x) if f \setminus  (x) = 0.
% As q(x) | f \setminus  (x), the only option is that f \setminus  (x) = 0.
%  ?
% This already tells us that irreducible polynomials are necessarily separable in
% characteristic 0: in characteristic 0, the derivative of a nonconstant
% polynomial
% clearly cannot vanish. In fact, Lemma 4.14 gives us a precise picture of what
% inseparable, irreducible polynomials must look like. If
% ?
%  n
% f (x) =
%  ai xi
% i=0
% is irreducible and inseparable, then the characteristic of the field must be a
% positive
% prime p, and by Lemma 4.14 we must have
% ?
%  n
% f \setminus (x) =
%  iai xi-1 = 0;
% i=0
% that is, iai = 0 for all i. Now iai = 0 is automatic if i is a multiple of p,
% and it
% implies ai = 0 for all the indices i which are not multiples of p. Therefore,
% the only
% nonvanishing coefficients in f (x) must be those corresponding to indices which
% are
% multiples of p:
% f (x) = a0 + ap xp + a2px2p + . . . ;
% therefore, f (x) must in fact be a polynomial in xp . (The reader will refine
% this
% picture further by working out Exercise 4.13 and following.)
% This leads us to a precise result linking separability and a flavor of fields
% that
% we have not yet run into.
% Definition 4.15. Let k be a field of characteristic p > 0. The Frobenius homo-
% morphism is the map k \to  k defined by x ?\to  xp .
%  ?
% The Frobenius homomorphism may not look like a homomorphism of rings, but
% it is (Exercise 4.8). It must be injective, as is every nontrivial ring
% homomorphism
% from a field; but it is not necessarily surjective.
% Definition 4.16. A field k is perfect if char k = 0 or if char k > 0 and the
% Frobenius
% homomorphism is surjective.
%  ?
% Proposition 4.17. Let k be a field. Then k is perfect if and only if all
% irreducible
% polynomials in k[x] are separable.
% Proof. I will prove that irreducible polynomials over a perfect field are
% separable,
% leaving the other implication to the reader (Exercise 4.12).
% We have already noted that irreducible polynomials are separable over fields of
% characteristic zero. In positive characteristic p, we have observed that an
% insepara-
% ble irreducible polynomial must be of the form
% ?
%  m
% f (x) =
%  ai \cdot  (xp )i .
% i=0
% Since Frobenius is surjective, there exist bi such that bpi = ai . Thus
% ? ?p
% ?
%  m
%  ?
%  m
%  ?
%  m
% f (x) =
%  bi p (xp ) i
%  =
%  (bi xi )p
%  =
%  bi x
% i
%  = g(x)p ,
% i=0
%  i=0
%  i=0
% ?
%  i
% where g(x) =
%  bi x and keeping in mind that the Frobenius map is a homomor-
% phism. But this contradicts the irreducibility of f (x), so there is no such
% polyno-
% mial.
%  ?
% Corollary 4.18. Finite fields are perfect. Therefore, over finite fields,
% irreducible
% polynomials are separable.
% Proof. The Frobenius map is injective (because it is a homomorphism of fields),
% so
% it is surjective over finite fields, by the pigeon-hole principle. Therefore
% finite fields
% are perfect, and the second part of the statement follows from Proposition 4.17.
% ?
% The reader now sees why I have chosen the somewhat unusual field Fp (t) in
% Example 4.11: we need a field of positive characteristic, but no example of
% irre-
% ducible inseparable polynomial can be found over Fp , by Corollary 4.18; to
% concoct
% an example, we need to enlarge the field so that it is infinite and to throw in
% an
% element (that is, t) for which there is no p-th root, that is, which is not in
% the image
% of Frobenius.

\subsection{Separable extensions and embeddings in algebraic closures}
\label{subsec:fields:separable-extensions-and-embeddings-in-algebraic-closures}
% The ter-

% minology examined in the previous section extends to the language of field
% exten-
% sions. If k \subseteq  F is an extension and \alpha  \in  F is algebraic over k,
% we say that \alpha  is
% separable over k if the minimal polynomial of \alpha  over k is separable;
% \alpha  is inseparable
% otherwise.
% Definition 4.19. An algebraic field extension k \subseteq  F is separable if
% every \alpha  \in  F is
% separable over k.
%  ?
% With this terminology, Proposition 4.17 may be restated in the following more
% impressive form:
% Proposition 4.20. A field k is perfect if and only if every algebraic extension
% of k
% is separable.
% In particular, algebraic extensions of Q (or any field of characteristic zero)
% and
% of every finite field are necessarily separable.
% The separability condition is exceedingly convenient, and we will essentially
% adopt it henceforth: all extensions we will seriously consider will be
% separable. One
% convenient feature of separable extensions is the following alternative
% description
% of the separability condition.
% We have seen (Lemma 2.8) that every algebraic extension k \subseteq  F may be
% embedded in an algebraic closure k \subseteq  k, and I pointed out that this can
% in general
% be done in many different ways. If finite, the number of different homomorphisms
% F \to  k extending the identity on k is denoted by
% [F : k]s .
% Definition 4.21. This is the separable degree of F over k.
%  ?
% It is clear that this number is independent of the chosen algebraic closure k
% \subseteq  k.
% What does it have to do with separability?
% Lemma 4.22. Let k \subseteq  k(\alpha ) be a simple algebraic extension. Then
% [k(\alpha ) : k]s equals
% the number of distinct roots in k of the minimal polynomial of \alpha . In
% particular,
% [k(\alpha ) : k]s \leq  [k(\alpha ) : k], with equality if and only if \alpha
% is separable over k.
% Proof. The proof is essentially (and not by coincidence) a rehash of the proof
% of
% Corollary 1.7. Associate with each \iota  : k(\alpha ) \to  k extending idk the
% image \iota (\alpha ), which
% must be a root of the minimal polynomial of \alpha . This correspondence is
% injective,
% since \iota (\alpha ) determines \iota  (as \iota  extends the identity on k).
% To see it is surjective, let
% \beta  \in  k be any other root, and consider the extension k(\beta ) \subseteq  k; by Proposition 1.5
% there is an isomorphism k(\alpha ) \to  k(\beta ) sending \alpha  to \beta , and
% composing with the
% embedding k(\beta ) \subseteq  k defines the corresponding \iota , as needed.
%  ?
% Thus, the `separable degree' does detect separability, for simple extensions.
% Further, it is multiplicative over successive extensions:
% Lemma 4.23. Let k \subseteq  E \subseteq  F be algebraic extensions. Then [F :
% k]s is finite if
% and only if both [F : E]s, [E : k]s are finite, and in this case
% [F : k]s = [F : E]s [E : k]s .
% Proof. Different embeddings of E into k extend to different embeddings of F into
% k,
% by Lemma 2.8; and embeddings of F into E = k extending the identity on E extend
% a fortiori the identity on k. Therefore, if any of [F : E]s , [E : k]s is
% infinite, then
% so is [F : k]s .
% For the converse implication, it suffices to prove the stated degree formula.
% But
% every embedding F \subseteq  k extending the identity on k may be obtained in
% two steps:
% first extend the identity to an embedding E \subseteq  k, which can be done in
% [E : k]s ways;
% and then extend the chosen embedding E \subseteq  k = E to an embedding F
% \subseteq  E = k,
% which can be done in [F : E]s ways. There are precisely [F : E]s[E : k]s ways to
% do this, as stated.
%  ?
% Lemmas 4.22 and 4.23 allow us to recast separability of finite extensions
% entirely
% in the light of counting embeddings in an algebraic closure:
% Proposition 4.24. Let k \subseteq  F be a finite extension. Then [F : k]s \leq
% [F : k], and
% the following are equivalent:
% (i) F = k(\alpha 1, . . . , \alpha r ), where each \alpha i is separable over k;
% (ii) k \subseteq  F is separable;
% (iii) [F : k]s = [F : k].
% Proof. Since F is finite over k, it is finitely generated. Let F = k(\alpha 1 ,
% . . . , \alpha r ).
% Then using Lemma 4.23, Lemma 4.22, and Proposition 1.10,
% [F : k]s = [k(\alpha 1 , . . . , \alpha r-1)(\alpha r ) : k(\alpha 1 , . . . ,
% \alpha r-1 )]s \dots  [k(\alpha 1 ) : k]s
% \leq  [k(\alpha 1 , . . . , \alpha r-1)(\alpha r ) : k(\alpha 1 , . . . , \alpha r-1 )] \dots  [k(\alpha 1 ) : k]
% = [F : k].
% This proves the stated inequality.
% (i) \implies  (iii): If each \alpha i is separable over k, then it is separable
% over the field
% k(\alpha 1 , . . . , \alpha i-1 ) (Exercise 4.15), so the inequality is an
% equality by Lemma 4.22.
% (iii) \implies  (ii): Assume [F : k]s = [F : k], and let \alpha  \in  F . We
% have k \subseteq  k(\alpha ) \subseteq  F ;
% hence by Lemma 4.23
% [F : k(\alpha )]s [k(\alpha ) : k]s = [F : k]s = [F : k] = [F : k(\alpha
% )][k(\alpha ) : k].
% Since both separable degrees are less than or equal to their plain counterparts,
% the
% equality implies
% [k(\alpha ) : k]s = [k(\alpha ) : k],
% proving that \alpha  is separable, by Lemma 4.22. Therefore the extension k
% \subseteq  F is
% separable according to Definition 4.19.
% (ii) \implies  (i) is immediate from Definition 4.19, since finite extensions
% are finitely
% generated.
%  ?
% For example, if \alpha  is separable over k, then every \beta  \in  k(\alpha )
% is separable. This
% would seem rather mysterious from the definition alone: why should the fact that
% the minimal polynomial of \alpha  has distinct roots in k imply the same for the
% minimal
% polynomial of \beta ? It does, by virtue of Proposition 4.24.
% Remark 4.25. While we have only proved [F : k]s \leq  [F : k], one can in fact
% prove
% that [F : k]s divides [F : k]: in fact, [F : k]s is the degree of a certain
% intermediate
% field Fsep over k. The diligent reader will prove this in Exercise 4.18.

\subsection*{Exercises}

% 4.1. ? Let k be a field, f (x) \in  k[x], and let F be the splitting field for f
% (x) over k.
% Let k \subseteq  K be an extension such that f (x) splits as a product of linear
% factors
% over K. Prove that there is a homomorphism F \to  K extending the identity on k.
% [\S{}4.2]
% 4.2. Describe the splitting field of x6 + x3 + 1 over Q. Do the same for x4 + 4.
% 4.3. ? Find the order of the automorphism group of the splitting field of x4 + 2
% over Q (cf. Example 4.6). [\S{}4.1]
% \sqrt 
%  44.4. Prove that the field Q( 2) is not the splitting field of any polynomial
%  over Q.
% 4.5. ? Let F be a splitting field for a polynomial f (x) \in  k[x], and let g(x)
% \in  k[x] be
% a factor of f (x). Prove that F contains a unique copy of the splitting field of
% g(x).
% [\S{}5.1]
% 4.6. Let k \subseteq  F 1 , k \subseteq  F2 be two finite extensions, viewed as
% embedded in the
% algebraic closure k of k. Assume that F1 and F2 are splitting fields of
% polynomials
% in k[x]. Prove that the intersection F1 \cap  F2 and the composite F1 F2 (the
% small-
% est subfield of k containing both F1 and F2 ) are both also splitting fields
% over k.
% (Theorem 4.8 is likely going to be helpful.)
% 4.7. ? Let k \subseteq  F = k(\alpha ) be a simple algebraic extension. Prove
% that F is normal
% over k if and only if for every algebraic extension F \subseteq  K and every
% \sigma  \in  Aut k (K),
% \sigma (F ) = F . [\S{}6.1]
% 4.8. ? Let p be a prime, and let k be a field of characteristic p. For a, b \in
% K, prove
% that16 (a + b)p = ap + bp . [\S{}4.2, \S{}5.1, \S{}5.2]
% 4.9. Using the notion of `derivative' given in \S{}4.2, prove that (f
% g)\setminus  = f \setminus  g + f g\setminus  for
% all polynomials f , g.
% 4.10. Let k \subseteq  F be a finite extension in characteristic p > 0. Assume
% that p does
% not divide [F : k]. Prove that k \subseteq  F is separable.
% 4.11. ? Let p be a prime integer. Prove that the Frobenius homomorphism on Fp
% is the identity. (Hint: Fermat.) [\S{}5.1]
% 4.12. ? Let k be a field, and assume that k is not perfect. Prove that there are
% inseparable irreducible polynomials in k[x]. (If char k = p and u \in  k, how
% many
% roots does xp - u have in k?) [\S{}4.2]
% 4.13. ? Let k be a field of positive characteristic p, and let f (x) be an
% irreducible
% polynomial. Prove that there exist an integer d and a separable irreducible
% polyno-
% mial fsep (x) such that
% d
% f (x) = fsep (xp ).
% The number pd is called the inseparable degree of f (x). If f (x) is the minimal
% polynomial of an algebraic element \alpha , the inseparable degree of \alpha  is
% defined to
% be the inseparable degree of f (x). Prove that \alpha  is inseparable if and
% only if its
% inseparable degree is \geq  p.
% The picture to keep in mind is as follows: the roots of the minimal polyno-
% mial f (x) of \alpha  are distributed into deg fsep `clumps', each collecting a
% number of
% coincident roots equal to the inseparable degree of \alpha . We say that \alpha
% is `purely
% inseparable' if there is only one clump, that is, if all roots of f (x) coincide
% (see
% Exercise 4.14). [\S{}4.2, 4.14, 4.18]
% 4.14. \lnot  Let k \subseteq  F be an algebraic extension, in positive
% characteristic p. An
% element \alpha  \in  F is purely inseparable over k if \alpha  pd
%  \in  k for some17 d \geq  0. The
% extension is defined to be purely inseparable if every \alpha  \in  F is purely
% inseparable
% over k.
% Prove that \alpha  is purely inseparable if and only if [k(\alpha ) : k]s = 1,
% if and only if
% its degree equals its inseparability degree (Exercise 4.13). [4.13, 4.17]
% 4.15. ? Let k \subseteq  F be an algebraic extension, and let \alpha  \in  F be
% separable over k.
% For every intermediate field k \subseteq  E \subseteq  F , prove that \alpha  is
% separable over E. [\S{}4.3]
% 4.16. \lnot  Let k \subseteq  E \subseteq  F be algebraic field extensions, and
% assume that k \subseteq  E is
% separable. Prove that if \alpha  \in  F is separable over E, then k \subseteq
% E(\alpha ) is a separable
% extension. (Reduce to the case of finite extensions.)
% Deduce that the set of elements of F which are separable over k form an inter-
% mediate field Fsep , such that every element \alpha  \in  F , \alpha  \in  Fsep
% is inseparable over Fsep .
% For F = k, k sep is called the separable closure of k. [4.17, 4.18]
% 16 This is sometimes referred to as the freshman's dream, for painful reasons
% that are likely
% all too familiar to the reader.
% 17 Note the slightly annoying clash of notation: elements of k are not
% inseparable, yet they
% are purely inseparable according to this definition.
% 4.17. \lnot  Let k \subseteq  F be an algebraic extension, in positive
% characteristic. With
% notation as in Exercises 4.14 and 4.16, prove that the extension Fsep \subseteq
% F is purely
% inseparable. Prove that an extension k \subseteq  F is purely inseparable if and
% only if
% Fsep = k. [4.18]
% 4.18. ? Let k \subseteq  F be a finite extension, in positive characteristic.
% Define the
% inseparable degree [F : k]i to be the quotient [F : k]/[F : k]s .
% \bullet  Prove that [k(\alpha ) : k]i equals the inseparable degree of \alpha , as defined in Exer-
% cise 4.13.
% \bullet  Prove that the inseparable degree is multiplicative: if k \subseteq  E \subseteq  F are finite
% extensions, then [F : k]i = [F : E]i [E : k]i .
% \bullet  Prove that a finite extension is purely inseparable if and only if its inseparable
% degree equals its degree.
% \bullet  With notation as in Exercise 4.16, prove that [F : k]s = [Fsep : k] and [F : k]i =
% [F : Fsep ]. (Use Exercise 4.17.)
% In particular, [F : k]s divides [F : k]; the inseparable degree [F : k]i is an
% integer.
% [\S{}4.3]
% 4.19. \lnot  Let k \subseteq  F be a finite separable extension, and let \iota 1
% , . . . , \iota d be the distinct
% embeddings of F in k extending idk . For \alpha  \in  F , prove that the norm
% Nk\subseteq F (\alpha )
% ?d
% (cf. Exercise 1.12) equals i=1 \iota i (\alpha ) and its trace trk\subseteq F
% (\alpha ) (Exercise 1.13) equals
% ?d
% i=1 \iota i (\alpha ). (Hint: Exercises 1.14 and 1.15.) [4.21, 4.22, 6.15]
% 4.20. \lnot  Let k \subseteq  F be a finite separable extension, and let \alpha
% \in  F . Prove that for
% all \sigma  \in  Autk (F ), Nk\subseteq F (\alpha /\sigma (\alpha )) = 1 and
% trk\subseteq F (\alpha  - \sigma (\alpha )) = 0. [6.16, 6.19]
% 4.21. \lnot  Let k \subseteq  E \subseteq  F be finite separable extensions, and
% let \alpha  \in  F . Prove that
% Nk\subseteq F (\alpha ) = Nk\subseteq E (NE\subseteq F (\alpha )) and
%  trk\subseteq F (\alpha ) = trk\subseteq E (trE\subseteq F (\alpha )).
% (Hint: Use Exercise 4.19: if d = [E : k] and e = [F : E], the de embeddings of F
% into k lifting idk must divide into d groups of e each, according to their
% restriction
% to E.)
% This `transitivity' of norm and trace extends the result of Exercise 1.15 to
% separable extensions. The separability restriction is actually unnecessary; cf.
% Exer-
% cise 4.22. [4.22]
% 4.22. Generalize Exercises 4.19---4.21 to all finite extensions k \subseteq  F .
% (For the norm,
% raise to power [F : k]i ; for the trace, multiply by [F : k]i .)

\section{Field extensions, III}
\label{sec:fields:field-extensions-iii}

% The material in \S{}4 provides us with the main tools needed to tackle several
% key
% examples. In this section we study finite and cyclotomic fields, and we return
% to
% the question of when a finite extension is in fact simple (cf. Example 1.19).

\subsection{Finite fields}
\label{subsec:fields:finite-fields}
%. Let F be a finite field, and let p be its characteristic}
% We know

% (\S{}1.1) that F may be viewed as an extension
% Fp \subseteq  F
% of Fp = Z/pZ; let d = [F : Fp ]. Since F has dimension d as a vector space over
% Fp ,
% it is isomorphic to F dp as a vector space, and in particular |F | = pd is a
% power of p.
% The general question we pose is whether there exist fields of cardinality equal
% to every power of a prime, and we aim to classifying all fields of a given
% cardinality.
% The conclusion we will reach is as neat as it can be: for every prime power q
% there
% exists exactly one field F with q elements, up to isomorphism.
% Also recall (Remark III.1.16) that, by a theorem of Wedderburn, every finite
% division ring is in fact a finite field. Thus, relaxing the hypothesis of
% commutativity
% in this subsection would not lead to a different classification.
% Theorem 5.1. Let q = p d be a power of a prime integer p. Then the polynomial
% xq - x is separable over Fp , and the splitting field of the polynomial xq - x
% over Fp
% is a field with precisely q elements. Conversely, let F be a field with exactly
% q
% elements; then F is a splitting field for xq - x over Fp.
% Proof. Let F be the splitting field of xq - x over Fp . Let E be the set of
% roots
% of f (x) = xq - x in F . Since f \setminus (x) = qxq-1 - 1 = -1 (as q = 0 in
% characteristic p),
% \setminus we have (f (x), f (x)) = 1; hence (Lemma 4.13) f (x) is separable, and E consists of
% precisely q elements. I claim that E is a field, and it follows that E = F .
% Indeed, F
% is generated by the roots of f (x); hence the smallest subfield of F containing
% E is
% F itself.
% To see that E is a field, let a, b \in  E. Then aq = a and bq = b; it follows
% that
% (a - b)q = aq + (-1)q bq = a - b
% (using Exercise 4.8; note that (-1)q = -1 if p is odd and (-1)q = +1 = -1 if
% p = 2). If b = 0,
% (ab-1)q = aq (bq )-1 = ab-1 .
% Thus E is closed under subtraction and division by a nonzero element, proving
% that
% E is a field and concluding the proof of the first statement.
% To prove the second statement, let F be a field with exactly q elements. The
% nonzero elements of F form a group under multiplication, consisting of q - 1
% ele-
% ments; therefore, the (multiplicative) order of every nonzero a \in  F divides q
% - 1
% (Example II.8.15). Therefore,
% a = 0 \implies  aq-1 = 1 \implies  aq - a = 0;
% of course 0q - 0 = 0. In other words, the polynomial xq - x has q roots in F
% (that
% is, all elements of F !); it follows that F is a splitting field for xq - x, as
% stated. ?
% Corollary 5.2. For every prime power q there exists one and only one finite
% field
% of order q, up to isomorphism.
% Proof. This follows immediately from Theorem 5.1 and the uniqueness of splitting
% fields (Lemma 4.2).
%  ?
% Since there is exactly one isomorphism class of fields of order q for any given
% prime power q, we can devise a notation for a field of order q; it makes sense
% to
% adopt18 Fq . This is called the Galois field of order q.
% Example 5.3. Let p be a prime integer. Then I claim that the polynomial x4 + 1
% is reducible 19 over Fp (and therefore over every finite field).
% Since x4 + 1 = (x + 1)4 in F2 [x], the statement holds for p = 2. Thus, we may
% assume that p is an odd prime. Then I claim that x4 + 1 divides xp2
%  - x. Indeed,
% the square of every odd number is congruent to 1 mod 8 (Exercise II.2.11); hence
% 8 | (p2 - 1); hence x8 - 1 divides xp2
%  -1 - 1 (Exercise V.2.13); hence
% (x4 + 1) | (x8 - 1) | (xp2
%  -1 - 1) | (xp2
%  - x).
% It follows that x4 + 1 factors completely in the splitting field of xp2
%  - x, that is,
% in Fp2 . If \alpha  is a root of x4 + 1 in Fp2 , we have the extensions
% Fp \subseteq  Fp (\alpha ) \subseteq  Fp2 ;
% therefore (Corollary 1.11) [Fp(\alpha ) : Fp ] divides [Fp2 : Fp ] = 2. That is,
% \alpha  has degree 1
% or 2 over Fp . But then its minimal polynomial is a factor of degree 1 or 2 of
% (x4 +1),
% showing that the latter is reducible.
%  ?
% Theorem 5.1 has many interesting consequences, and we sample a few in the
% rest of this subsection.
% Corollary 5.4. Let p be a prime, and let d \leq  e be positive integers. Then
% there is
% an extension Fpd \subseteq  Fpe if and only if d | e. Further, if d | e, then
% there is exactly
% one such extension, in the sense that Fpe contains a unique copy of Fp d .
% All extensions F pd \subseteq  Fpe are simple.
% Proof. If there is an extension as stated, then Fp \subseteq  Fp d \subseteq
% Fpe ; hence [Fpd : Fp ]
% divides [Fpe : Fp ] by Corollary 1.11. This says precisely that d | e.
% Conversely, assume that d | e. As
% pe - 1 = (pd - 1)((pd ) d e
%  -1 + \dots  + 1),
% we see that pd - 1 divides pe - 1, and consequently xpd
%  -1 - 1 divides xpe
%  -1 - 1
% (Exercise V.2.13). Therefore
% (xpd
%  - x) | (xpe
%  - x).
% By Theorem 5.1, Fpe is a splitting field for the second polynomial. It follows
% that it
% contains a unique copy of the splitting field for the first polynomial (Exercise
% 4.5),
% that is, of Fpd .
% For the last statement, recall that the multiplicative group of nonzero elements
% of a finite field is necessarily cyclic (Theorem IV.6.10). If \alpha  \in  Fpe
% is a generator of
% this group, then \alpha  will generate Fpe over any subfield; if d | e, this
% says Fpe = Fpd (\alpha ),
% so Fp d \subseteq  Fpe is simple.
%  ?
% 18 Keep in mind that Fq
%  is not the ring Z/qZ unless q is prime: if q is composite, then Z/qZ
% is not an integral domain.
% 19 Of course x4 + 1 is irreducible over Z. This example shows that Proposition
% V.5.15 cannot
% be turned into an `if and only if' statement.
% These results can be translated into rather precise information on the structure
% of the polynomial ring over a finite field. For example,
% Corollary 5.5. Let F be a finite field. Then for all integers n \geq  1 there
% exist
% irreducible polynomials of degree n in F [x].
% Proof. We know F = Fpd for some prime p and some d \geq  1. By Corollary 5.4
% there
% is an extension Fpd \subseteq  Fpdn , generated by an element \alpha . Then
% [Fpdn : F pd ] = n,
% and it follows that the minimal polynomial of \alpha  over F = Fpd is an
% irreducible
% polynomial of degree n in F [x] .
%  ?
% In fact, our analysis of extensions of finite fields tells us about explicit
% fac-
% torizations, leading to an inductive algorithm to find all irreducible
% polynomials
% in Fq [x]:
% Corollary 5.6. Let F = Fq be a finite field, and let n be a positive integer.
% Then
% the factorization of xqn
%  - x in F [x] consists of all irreducible monic polynomials
% of degree d, as d ranges over the positive divisors of n. In particular, all
% these
% polynomials factor completely in Fqn .
% Proof. By Theorem 5.1, Fq n is the splitting field of xq n
%  - x over Fp , and hence
% over Fq = F .
% If f (x) is a monic irreducible polynomial of degree d, then F [x]/(f (x)) = F
% (\alpha )
% is an extension of degree d of F , that is, an isomorphic copy of Fqd . By
% Corollary 5.4,
% if d | n, then there is an embedding of Fqd in Fq n . But then \alpha  must be a
% root
% of xqn
%  -x, and hence xqn
%  -x is a multiple of f (x), as this is the minimal polynomial
% of \alpha . This proves that every irreducible polynomial of degree d | n is a
% factor
% of xqn
%  - x.
% Conversely, if f (x) is an irreducible factor of xq n
%  - x, then Fqn contains a
% root \alpha  of f (x); we have the extensions F = Fq \subseteq  Fq (\alpha )
% \subseteq  Fqn , and Fq (\alpha ) = \sim
%  Fq d for
% d = deg \alpha . It follows that d | n, again by Corollary 5.4.
%  ?
% The picture I am trying to convey is the following: the qn roots of xq n
%  - x
% clump into disjoint subsets, with each subset collecting the roots of each and
% every
% irreducible polynomial of degree d | n in F [x].
% Example 5.7. Let's contemplate the case q = 2: F2 = Z/2Z.
% \bullet  n = 1: the polynomial x2 - x factors as the product of x and (x - 1) (which
% we could write as (x + 1) just as well, since we are working over F2). These are
% all
% the irreducible polynomials of degree 1 over F 2 .
% \bullet  n = 2: the polynomial x4 - x must factor as the product of all irreducible
% polynomials of degree 1 and 2; in fact
% x4 - x = x(x - 1)(x2 + x + 1),
% and the conclusion is that there is exactly one irreducible polynomial of degree
% 2
% over F2 , namely x2 + x + 1.
% \bullet  n = 3: the quotient of x8 - x by x(x - 1) is a polynomial of degree 6, which
% must therefore be the product of the two irreducible polynomials of degree 3
% over F2 .
% It takes a moment to find them:
% x3 + x2 + 1,
%  x3 + x + 1.
% It also follows that F8 may be realized in two ways as a quotient of F2 [x]
% modulo
% an irreducible polynomial:
% F2 [x]
%  F2 [x]
% \sim 
%  = .
% (x3
%  + x2
%  + 1)
%  (x3
%  + x + 1)
% We know that these two fields must be isomorphic because of Theorem 5.1; it is
% instructive to find an explicit isomorphism between them (Exercise 5.3).
% \bullet  n = 4 and 5: these cases are left to the enjoyment of the reader (Exercise 5.4).
% \bullet  n = 6: the factorization of x64 - x must include x, x - 1, the one irreducible
% polynomial of degree 2, and the two irreducible polynomials of degree 3 found
% above,
% and that leaves room for 9 polynomials of degree 6, which must be all and only
% the
% irreducible polynomials of degree 6 over F2 . So the 64 elements of F 64 cluster
% as
% follows:
% The two dotted rectangles delimit the (unique) copies of F4 and F8 contained in
% F64 ;
% these intersect in the (unique) copy of F2 .
% Again, F64 may be realized by quotienting F2 [x] by the ideal generated by any
% of the irreducible polynomials of degree 6; this gives 9 `different'
% realizations of this
% field.
%  ?
% The picture drawn above for F64 means next to nothing, but it may help us focus
% on one last element of information we are going to extract regarding finite
% fields.
% Note that the effect of any automorphism of F64 must be to scramble elements in
% each sector represented in the picture, without mixing elements in different
% sectors
% and without interchanging sectors. Indeed, every automorphism of an extension
% sends roots of an irreducible polynomial to roots of the same polynomial.
% Since extensions of finite fields are simple extensions, our previous work
% allows
% us to be much more precise. Restricting our attention to the extensions Fp
% \subseteq  Fpd ,
% for a prime p, we know these can be realized as simple extensions by an element
% with minimal polynomial of degree d. This polynomial is necessarily separable
% (Fp is perfect), so Corollary 1.7 immediately gives us the size of the
% automorphism
% group:
% | AutFp (Fpd )| = d.
% But what is this group?
% Proposition 5.8. AutFp (Fpd ) is cyclic, generated by the Frobenius isomorphism.
% Proof. Let \phi  be the Frobenius homomorphism Fpd \to  Fpd : \phi (x) = xp .
% The
% Frobenius homomorphism is an isomorphism on a finite field (Corollary 4.18) and
% restricts to the identity on Fp (Exercise 4.11), so \phi  \in  AutFp (Fpd ).
% Since we know a
% priori the size of this group, all we have to show is that the order of \phi  is
% d.
% Let e = |\phi |. Then \phi e = id; hence xpe
%  = x for all x \in  Fpd . In other
% words, the (nonzero) polynomial xpe
%  - x has pd roots in Fpd ; this implies pd \leq  pe
% (Lemma V.5.1). Equivalently, d \leq  e, yielding
% | AutFp (Fpd )| \leq  |\phi |.
% But then these two numbers must be equal, since the order of an element always
% divides the order of the group (Example II.8.15).
%  ?
% Two last comments on finite fields are in order:
% ---For n large enough, the stupendously large (but finite) field Fpn! works as
% an
% approximation of the algebraic closure F p : indeed, this field contains roots
% of all
% polynomials of degree \leq  n in Fp[x], by Corollary 5.6.
% This observation can be turned into the explicit construction of an algebraic
% closure of Fp: by Corollary 5.4 there are extensions
% Fp \subseteq  Fp2 \subseteq  Fp6 \subseteq  Fp4! \subseteq  \dots 
% and the union of this chain of fields20 gives a copy of Fp .
% ---Finite fields form a category, which (by Corollary 5.4) should strongly
% remind
% the reader of the `toy' category encountered in Exercise I.5.6. There is an
% impor-
% tant difference, however: in that category every object has exactly one automor-
% phism, while we have just checked that finite fields may have many
% automorphisms.
% Challenge: The reader has encountered a very similar situation in a different
% (but
% related) context. Where?

\subsection{Cyclotomic polynomials and fields}
\label{subsec:fields:cyclotomic-polynomials-and-fields}
% In the last several sections we have

% often encountered roots of 1 in C; the extensions they generate are very
% important.
% Let n be a positive integer. I will denote by \zeta n the complex number e2\pi
% i/n ;
% thus, the n roots of the polynomial xn - 1 in C are the n distinct powers of
% \zeta n ;
% they form a cyclic subgroup of order n of the multiplicative group of C, which I
% will denote \mu n .
% Pictorially, the roots of 1 are placed at the vertices of a regular n-gon
% centered
% at 0, with one vertex at 1.
% 20 A more formal construction requires the notion of direct limit; cf.
% \S{}VIII.1.4.
% \zeta n
% A primitive n-th root of 1 is a generator of \mu n . Thus, \zeta n is primitive;
% from our
% work on cyclic groups, we know (Corollary II.2.5) that \zeta n m is a primitive
% n-th root
% of 1 if and only if m is relatively prime to n. In particular, there are \varphi
% (n) primitive
% roots of 1, where \varphi  is Euler's totient function (cf. Exercise II.6.14).
% Thus,
% ?
%  ?
% \Phi n (x) :=
%  (x - \zeta ) =
%  (x - \zeta n m )
% \zeta  primitive n-th root of 1
%  1\leq m\leq n,(m,n)=1
% is a polynomial of degree \varphi (n).
% Definition 5.9. The polynomial \Phi n (x) is called the n-th cyclotomic
% polynomial. ?
% It is clear that \Phi n (x) is a monic polynomial of degree \varphi (n). It is
% perhaps a
% little less immediate that \Phi n (x) has rational (in fact, integer )
% coefficients and that
% it is irreducible over Q. We will prove these facts in short order.
% Example 5.10. If n = p is prime, then every nonidentity element of \mu p \sim 
%  = Cp is a
% generator: every p-th root of 1 is primitive except 1 itself. Therefore
% xp - 1
% \Phi p (x) =
%  = xp-1 + \dots  + 1
% x - 1
% is the particular case encountered in Example V.5.19, where we proved that \Phi
% p (x)
% is indeed irreducible.
%  ?
% What if n is not prime?
% Lemma 5.11. For all positive integers n,
% ?
% xn - 1 =
%  \Phi d (x).
% 1\leq d|n
% Proof. If n = de, then every d-th root \zeta  of 1 is an n-th root of 1, because
% \zeta  n =
% \zeta  de = (\zeta  d )e = 1. In particular, every primitive d-th root \zeta  of 1 is an n-th root of 1.
% On the other hand, every \zeta  \in  \mu n generates a subgroup H of \mu n, and
% H = \mu d for
% d equal to the order of \zeta , a divisor of n (Proposition II.6.11). Thus,
% every \zeta  \in  \mu n
% is a primitive d-th root of 1 for some d | n.
% Thus the set of n-th roots of 1 equals the union of the sets of primitive d-th
% roots
% of 1, as d ranges over all positive divisors of n. The statement follows
% immediately:
% (
%  )
% ?
%  ?
%  ?
%  ?
% xn - 1 =
%  (x - \zeta ) =
%  (
%  (x - \zeta )) =
%  \Phi d (x),
% \zeta \in \mu n
%  1\leq d|n
%  \zeta  primitive d-th root of 1
%  1\leq d|n
% as claimed.
%  ?
% This argument brings us back to simple group-theoretic considerations. The
% reader can check that Lemma 5.11 implies immediately the result of Exercise
% II.6.14,
% by comparing degrees of both sides of the stated identity.
% Lemma 5.11 yields an inductive computation of cyclotomic polynomials; the
% fact that \Phi n (x) \in  Z[x] follows from this fact. Explicitly,
% Corollary 5.12. The cyclotomic polynomials \Phi n (x) have integer coefficients.
% Proof. Use induction on n. Note that \Phi 1 (x) = x - 1, and assume we have
% ?
%  shown
%  that all \Phi m (x) have integer coefficients for m < n. In particular, f (x)
%  :=
% 1\leq d|n,d<n \Phi d(x) is a monic polynomial with integer coefficients. Since f
% (x) is
% monic, we can divide it into xn - 1 with remainder, within Z[x]: \exists q(x),
% r(x) \in  Z[x]
% such that
% xn - 1 = f (x)q(x) + r(x),
% with r(x) = 0 or deg r(x) < deg f (x). On the other hand, by Lemma 5.11,
% xn - 1 = f (x)\Phi n (x)
% in C[x]. Therefore
% f (x)(\Phi n(x) - q(x)) = r(x)
% in C[x]. But this forces r(x) = 0 (otherwise we would have deg r(x)\geq  deg f
% (x)).
% Therefore \Phi n(x) = q(x) \in  Z[x].
%  ?
% Example 5.13. The reader can spend some quality time computing explicitly the
% cyclotomic polynomials \Phi n (x) for several nonprime numbers n, working
% inductively
% and capitalizing on the fact that we know explicitly \Phi p (x) for prime p.
% For example, x4 - 1 = \Phi 1 (x)\Phi 2 (x)\Phi 4 (x); therefore
% x4 - 1
% \Phi 4 (x) =
%  = x2 + 1.
% x2 - 1
% Since x6 - 1 = \Phi 1 (x)\Phi 2 (x)\Phi 3 (x)\Phi 6 (x),
% x6 - 1
% \Phi 6 (x) =
%  = x2 - x + 1.
% (x2 - 1)(x2 + x + 1)
% Since x12 - 1 = \Phi 1 (x)\Phi 2(x)\Phi 3(x)\Phi 4(x)\Phi 6(x)\Phi 12 (x),
% x12 - 1
% \Phi 12 (x) =
%  = x4 - x2 + 1,
% (x6 - 1)(x2 + 1)
% and so on ad libitum.
%  ?
% The optimistic reader may now start guessing that the coefficients of \Phi n (x)
% are always 0 or \pm 1 (so would I, if I didn't know better). Apparently the
% first
% counterexample is found for n = 105 (Exercise 5.9), and the coefficients are
% known
% to get as large as one pleases for n ? 0.
% The irreducibility of \Phi n (x) is a bit trickier, in particular because
% separabil-
% ity sneaks into the standard argument (which is why I had to wait until now to
% present it).
% Proposition 5.14. For all positive n, \Phi n (x) \in  Z[x] is irreducible over
% Q.
% (Since \Phi n (x) is monic, irreducibility in Q[x] is equivalent to
% irreducibility
% in Z[x]; cf. Corollary V.4.17.)
% Proof. Arguing by contradiction, assume \Phi n (x) is reducible. Then its roots
% \zeta n m ,
% with (m, n) = 1, are divided among the factors; we can choose a root \zeta n m
% of one
% irreducible monic factor f (x), such that another root \zeta n mp (for some
% prime p not
% dividing n) is not a root of f (x). Write
% \Phi n (x) = f (x)g(x);
% since \Phi n (x) \in  Z[x] and \Phi n(x), f (x) are monic, then f (x) and g(x)
% have integer
% coefficients (cf. Exercise V.4.23). By our choice, f (x) is the minimal
% polynomial
% of \zeta n m over Q, and g(\zeta  n mp ) = 0.
% It follows that \zeta n m is a root of g(xp ), and hence f (x) | g(xp ).
% Therefore we can
% write
% g(xp ) = f (x)h(x)
% with h(x) \in  Z[x]. Reading the last equation modulo p, we get (again using
% Exer-
% cise 4.8, and denoting cosets by underlining)
% g(x)p = f (x)h(x) in F p [x];
% in particular, f (x) and g(x) must have a nontrivial common factor (x) in Fp
% [x].
% But then
% (x)2 | f (x)g(x) :
% the reduction of \Phi n(x) modulo p must have a multiple factor.
% This implies that xn - 1 \in  Fp [x] has a multiple factor; that is, it is
% inseparable.
% However, its derivative nxn-1 \in  Fp [x] is nonzero (because p does not divide
% n by
% assumption), and Lemma 4.13 implies that xn - 1 is separable in Fp [x].
% This contradiction shows that our assumption that \Phi n(x) is reducible must be
% nonsense, proving the statement.
%  ?
% Definition 5.15. The splitting field Q(\zeta n) for the polynomial xn - 1 over Q
% is the
% n-th cyclotomic field.
%  ?
% By Proposition 5.14, Q(\zeta n ) is an extension of Q of degree \varphi (n);
% \Phi n (x) is the
% minimal polynomial of \zeta n .
% As usual, we now preoccupy ourselves with the automorphism group of this
% extension, and as in previous examples the situation is very neat.
% Proposition 5.16. AutQ (Q(\zeta n)) is isomorphic to the group of units in Z/nZ.
% Proof. We know that AutQ (Q(\zeta n )) has cardinality \varphi (n) (Corollary
% 1.7; the roots
% are distinct since \Phi n (x) is separable), so all we need to do is exhibit an
% injective
% homomorphism
% j : (Z/nZ)\ast  \to  AutQ (Q(\zeta n )).
% For [m]n \in  (Z/nZ)\ast  , that is, for m an integer relatively prime to n, let
% j([m]n ) be
% the unique automorphism Q(\zeta n ) \to  Q(\zeta n ) sending \zeta n to \zeta n
% m (cf. Proposition 1.5);
% this is evidently independent of the representative m. Then j is clearly
% injective,
% and
% j([m1]n ) \circ  j([m2]n )(\zeta n) = j([m1]n )(\zeta n m2 ) = \zeta n m2 m1
% agrees with the action of j([m1 ]n[m2 ]n ), showing that j is a homomorphism and
% concluding the proof.
%  ?
% Example 5.17. The reader should consider Example 4.4 again: by what we have
% just proved, the automorphism group of the splitting field of x8 - 1 is
% isomorphic
% to the group of units in Z/8Z; this group is immediately seen to be isomorphic
% to (Z/2Z) \times  (Z/2Z), confirming the claim made in Example 4.4.
%  ?
% The constructibility of the regular n-gon by straightedge and compass is of
% course equivalent to the constructibility of the complex number \zeta n . Since
% we now
% know that [Q(\zeta n ) : Q] = \varphi (n), Corollary 3.6 tells us that the
% regular n-gon is con-
% structible by straightedge and compass only if \varphi (n) is a power of 2. This
% condition
% can be chased a little further (Exercise 5.19). In any case, we shall return to
% this
% theme one more time, after we absorb a little Galois theory.

\subsection{Separability and simple extensions}
\label{subsec:fields:separability-and-simple-extensions}
% All the examples of field extensions

% we have encountered so far were simple extensions, although sometimes this may
% not be completely apparent at first (cf. Example 1.19). There is a good reason
% for
% this: most of the examples we have seen were separable extensions, and we can
% now
% prove that a finite separable extension is necessarily simple.
% Here is a nice criterion for simplicity (which does not involve any separability
% condition):
% Proposition 5.18. An algebraic extension k \subseteq  F is simple if and only if
% the
% number of distinct intermediate fields k \subseteq  E \subseteq  F is finite.
% Proof. Assume F = k(\alpha ) is simple and algebraic, and let qk (x) be the
% minimal
% polynomial of \alpha  over k. Embed F in an algebraic closure k of k. If E is an
% intermediate field, then F = E(\alpha ) is also a simple, algebraic extension;
% denote
% by qE (x) the minimal polynomial of \alpha  over E. Since qk (x) \in  E[x] for
% all E and
% qk (\alpha ) = 0, we know that each qE (x) is a factor of qk (x).
% I claim that E is in fact determined by qE (x). Since qk (x) has finitely many
% factors in k, this proves that there are only finitely many intermediate fields,
% that
% is, the `only if' part of the statement.
% To verify my claim, I will show that E is generated (over k) by the coefficients
% \setminus of qE (x). To see this, let E be the subfield of E generated by k and the coefficients
% of qE (x). Then qE (x) \in  E \setminus  [x], and since qE (x) is irreducible
% over E, it must be
% irreducible over E \setminus  . Note that E \setminus  (\alpha ) = F = E(\alpha
% ). We have E \setminus  \subseteq  E \subseteq  F ; therefore
% \setminus  \setminus  \setminus deg(qE (x)) = [F : E ] = [F : E][E : E ] = deg(qE (x))[E : E ],
% \setminus  \setminus from which [E : E ] = 1; that is, E = E as needed. This concludes the proof of
% the `only if' part of the statement.
% To prove the `if' part, assume that there are only finitely many intermediate
% fields k \subseteq  E \subseteq  F . The extension k \subseteq  F must be
% finitely generated (otherwise
% we could construct an infinite sequence of subextensions k \subseteq  k(\alpha 1
% ) \subseteq  k(\alpha 1 , \alpha 2) \subseteq
% \dots  \subseteq  F , and we would have infinitely many intermediate fields), hence finite since
% it is algebraic. If k is a finite field, then so is F , and the extension is
% automatically
% simple by Corollary 5.4; thus we may assume that k is infinite, and we have to
% show
% that every finitely generated algebraic extension F = k(\alpha 1 , . . . ,
% \alpha r ) is simple.
% Arguing inductively, we may assume without loss of generality that F = k(\alpha
% , \beta ).
% For every c \in  k, we have the intermediate field
% k \subseteq  k(c\alpha  + \beta ) \subseteq  k(\alpha , \beta ).
% But there are only finitely many intermediate fields, while k is infinite, so
% for some
% c = c\setminus  in k we must have
% k(c\setminus  \alpha  + \beta ) = k(c\alpha  + \beta ).
% It follows that
% (c\setminus  \alpha  + \beta ) - (c\alpha  + \beta )
% \alpha  =
%  \in  k(c\alpha  + \beta ) and
%  \beta  = (c\alpha  + \beta ) - c\alpha  \in  k(c\alpha  + \beta )
% c\setminus  - c
% and therefore k(\alpha , \beta ) \subseteq  k(c\alpha  + \beta ). The other
% inclusion holds trivially, so k(\alpha , \beta ) =
% k(c\alpha  + \beta ) is simple.
%  ?
% Here is the connection with separability:
% Proposition 5.19. Every finite separable extension is simple.
% Proof. Arguing inductively as in the proof of Proposition 5.18, we may assume
% F = k(\alpha , \beta ), with \alpha  and \beta  separable (and in particular
% algebraic) over k, and we
% may assume k is an infinite field.
% Consider the set I of embeddings \iota  : F \setminus \to  k of F in an
% algebraic closure of k,
% extending the identity of k. If \iota  = \iota \setminus  in I and x is an
% indeterminate, then the
% polynomials
% \iota (\alpha )x + \iota (\beta ), \iota \setminus (\alpha )x + \iota \setminus  (\beta )
% are different: otherwise \iota \setminus  (\alpha ) = \iota (\alpha ) and \iota
% \setminus  (\beta ) = \iota (\beta ), so \iota , \iota \setminus  would act in
% the same
% way on the whole of k(\alpha , \beta ), forcing \iota  = \iota \setminus  .
% Therefore, the polynomial
% ?
% f (x) =
%  ((\iota (\alpha )x + \iota (\beta )) - (\iota \setminus (\alpha )x + \iota
%  \setminus  (\beta ))) \in  k[x]
% \iota ?=\iota ?
% is not identically 0. Since k is infinite, it follows that \exists c \in  k such
% that f (c) = 0;
% that is, distinct \iota  \in  I map
% \gamma  = c\alpha  + \beta 
% to distinct elements
% \iota (\gamma ) = \iota (\alpha )c + \iota (\beta )
% (each \iota  extends idk , so \iota (c) = c). Since the cardinality of I is [F :
% k]s (Defini-
% tion 4.21) and each \iota (\gamma ) is a root of the minimal polynomial of
% \gamma  over k, we have
% [F : k]s \leq  [k(\gamma ) : k] \leq  [F : k].
% By assumption the extension is separable, so [F : k]s = [F : k] by Proposition
% 4.24,
% implying [k(\gamma ) : k] = [F : k] and finally
% F = k(\gamma ),
% concluding the proof.
%  ?
% Proposition 5.19 is a good illustration of why separability is a technically
% desir-
% able condition. To stress the point even further, note that we are now in a
% position
% to extend the convenient inequality found in Corollary 1.7 to all finite
% separable
% extensions:
% Corollary 5.20. Let k \subseteq  F be a finite, separable extension. Then
% | Autk (F )| \leq  [F : k],
% with equality if and only if k \subseteq  F is a normal extension.
% Proof. Since k \subseteq  F is finite and separable, Proposition 5.19 implies it
% is simple:
% F = k(\alpha ) for some \alpha  \in  F . The inequality follows immediately from
% Corollary 1.7,
% and equality holds if and only if the minimal polynomial f (x) of \alpha
% factors into
% distinct linear factors in F . In this case F is the splitting field of f (x),
% so it
% is normal over k by Theorem 4.8. Conversely, if F is normal over k, then f (x)
% splits completely in F and has distinct roots since \alpha  is separable over k;
% therefore
% | Autk (F )| = [F : k], again by Corollary 1.7.
%  ?
% Example 5.21. By Proposition 5.19, if we want to construct a nonsimple finite
% extension, we have to use inseparable elements; to produce an example, we jazz
% up
% Example 4.11 a little. Consider the field of rational functions F = Fp (u, v) in
% two
% variables over Fp and the subfield k = Fp (s, t), where s = up and t = v p .
% Then
% k \subseteq  F is an algebraic extension; the minimal polynomial of u over k is
% xp - s; the
% minimal polynomial of v over k(u) is y p - t; it follows that
% [F : k] = p2.
% As c ranges in k, we obtain intermediate fields
% k \subseteq  k(cu + v) \subseteq  F.
% If any two choices c, c\setminus  led to the same intermediate field, then we
% would deduce
% that k(u, v) = k(cu+v) by arguing as in the proof of Proposition 5.18. In
% particular,
% we would have
% [k(cu + v) : k] = p2 .
% But this is not the case, since (cu + v)p = cp s + t \in  k(s, t) = k; hence the
% minimal
% polynomial of cu + v over k has degree at most p. Since k = Fp (s, t) is
% infinite,
% there are infinitely many intermediate fields, and it follows (by Proposition
% 5.18)
% that F is not a simple extension of k.
%  ?

\subsection*{Exercises}

% 5.1. Let p > 0 be a prime integer. Prove that the additive group of a finite
% field
% with pd elements is isomorphic to (Z/pZ)d .
% 5.2. Prove that every element of a finite field F is a sum of two squares in F .
% (Hint:
% Keep in mind that the multiplicative group of F is cyclic.)
% 5.3. ? Find an explicit isomorphism
% F2[x]
%  \sim 
%  F2 [x]
% /
%  .
% (x3 + x2 + 1)
%  (x3 + x + 1)
% [\S{}5.1]
% 5.4. ? Find all irreducible polynomials of degree 4 over F2 , and count those of
% degree 5. [\S{}5.1]
% 5.5. Find the number of irreducible polynomials of degree 12 over F9 .
% 5.6. Write out explicitly the action of the cyclic group C4 on F16, in terms of
% any
% realization of this field as a quotient of F2 [x].
% 5.7. Let p be a prime integer. View the Frobenius automorphism \phi  : Fpd \to
% Fpd as
% a linear transformation of the Fp -vector space Fpd . Find the rational
% canonical form
% of \phi . (Adapt the proof of Proposition 5.8 to show that the minimal
% polynomial
% of \phi  is xd - 1.)
% 5.8. For a prime p, find the factorization of \Phi  p (x) over Fp .
% 5.9. ? Find all cyclotomic polynomials \Phi n(x) for 1 \leq  n \leq  15. Compute
% \Phi 105 (x).
% [\S{}5.2]
% 5.10. Find the cyclotomic polynomials \Phi 2m (x) for all m \geq  0.
% 5.11. Prove that if n > 1 is odd, then \Phi 2n (x) = \Phi n(-x). (Hint: Draw the
% primitive
% 14-th roots of 1 side-by-side to the primitive 7-th roots of 1; then go back to
% Exercise II.2.15 to justify the fact you observe.)
% 5.12. \lnot  Let a, n be positive integers, with a > 1. Prove that if \Phi n (a)
% divides a - 1,
% then n = 1. (Remember that \Phi n (a) is a product of complex numbers a - \zeta
% , where
% \zeta  is a primitive n-th root of 1. What does this tell you about the size of \Phi n(a)?)
% [5.14]
% 5.13. \lnot  Let a, d, n be positive integers, with d < n and a > 1. Assume that
% ad - 1
% divides an - 1. Prove that \Phi n (a) divides the quotient (an - 1)/(ad - 1).
% (Hint:
% Exercise V.2.13.) [5.14]
% 5.14. ? Let R be a finite division ring.
% \bullet  Prove that the center of R (Exercise III.2.9) is isomorphic to Fq , for q a prime
% power. Prove that |R| = q n for some n.
% \bullet  For every r \in  R, prove that the centralizer of r in the multiplicative group
% (R\ast  , \cdot ) has order q d - 1 for some d \leq  n.
% \bullet  Prove that there are integers d1 , . . . , dr < n such that
% n
%  ?
%  r
%  qn - 1
% (*)
%  q - 1= q - 1+
%  .
% qdi - 1
% i=1
% (Hint: Class equation.)
% \bullet  Deduce that \Phi n (q) divides q-1 and hence n = 1. (Use Exercises 5.12 and 5.13.)
% \bullet  Conclude that R equals its center, showing that R is commutative.
% Thus, every finite division ring is a field: this is Wedderburn's little
% theorem. The
% argument given here is due to Ernst Witt. [\S{}III.1.2]
% 5.15. ? Let a, p, n be integers, with p, n positive and p prime, p ? n.
% \bullet  Show that xn - 1 has no multiple roots modulo p.
% \bullet  Show that if p divides \Phi n (a), then a n \equiv  1 modulo p. (In particular, p ? a, so
% [a]p \in  (Z/pZ)\ast  .)
% \bullet  Show that if p divides \Phi n (a), then ad \equiv  1 modulo p for every d < n.
% \bullet  Deduce that p | \Phi n (a) if and only if the order of [a]p in (Z/pZ)\ast  is n.
% \bullet  Compute \Phi 15 (9), and show it is divisible by 31. Then look back at the first
% part of Exercise II.4.12.
% [\S{}II.4.3, 5.16, 5.17]
% 5.16. If q is a prime that divides an - 1 and does not divide ad - 1 for any d <
% n,
% then q is said to be a primitive prime divisor of an - 1. By Exercise 5.15, if q
% ? n
% and q | \Phi n (a), then q is a primitive prime divisor of an - 1. The
% Birkhoff-Vandiver
% theorem asserts that an - 1 has primitive prime divisors for all but a very
% short
% list of exceptions: n = 1, a = 2; n = 2, a + 1 a power of 2; and n = 6, a = 2.
% Assuming this statement, we can give another proof of Wedderburn's theorem
% (published in 2003 by Nicolas Lichiardopol). Let R be a finite division ring.
% \bullet  Prove that there is a prime p such that |R| = pn for some integer n.
% \bullet  If n = 1, then R \sim 
%  = F p . If n = 2 or if p = 2 and n = 6, then R is commutative:
% the reader has proved this in Exercises III.2.11 and IV.2.17.
% \bullet  Therefore, by Birkhoff-Vandiver we may assume that pn - 1 has a primitive
% prime divisor p\setminus  .
% \bullet  Prove that there is an element a of order p\setminus  in the multiplicative group (R\ast  , \cdot )
% of R.
% \bullet  Prove that R is the only sub-division ring of R containing a.
% \bullet  Prove that the centralizer of a in R (Exercise III.2.10) is R itself; deduce that
% a is in the center of R (Exercise III.2.9).
% \bullet  Prove that the center of R is R: this shows that R is commutative.
% Thus, R is a finite field.
% Note that the Birkhoff-Vandiver theorem gives another rapid way to conclude
% Witt's proof: a primitive prime divisor of qn - 1 would divide q - 1 by (*),
% showing
% that n = 1.
% 5.17. \lnot  Let a, p, n be integers, with p, n positive and p prime, p ? n.
% Assume that
% p divides \Phi n (a). Prove that p \equiv  1 mod n. (Use Exercise 5.15.) [5.18]
% 5.18. ? Let n > 0 be any integer. Prove that there are infinitely many prime
% integers \equiv  1 mod n. (Use Exercise 5.17 together with Exercise V.2.25.)
% The result of this exercise is a particular case of Dirichlet's theorem on
% primes
% in arithmetic progressions; see \S{}7.6. [\S{}7.6]
% 5.19. ? Prove that the regular n-gon can be constructed by straightedge and com-
% pass only if n = 2m p1 . . . pr , where m \geq  0 and the factors pi are
% distinct Fermat
% primes. (Hint: Use Exercise V.6.8.) [\S{}5.2, \S{}7.2]
% 5.20. Recall from Exercise 1.10 that every field extension of degree \leq  n
% over a
% field k is a sub-k-algebra of the ring of matrices Mn (k). Prove that if k is
% finite or
% has characteristic 0, then every extension of k contained in Mn (k) has degree
% \leq  n.
% 5.21. Prove that if k \subseteq  F is the splitting field of a separable
% polynomial, then it is
% the splitting field of an irreducible separable polynomial.
% 5.22. Let k be an infinite field. If F = k(\alpha 1 , . . . , \alpha r ), with
% each \alpha i separable over k,
% prove that there exist c1 , . . . , cr \in  k such that F = k(c1 \alpha 1 +
% \dots  + cr \alpha r ).
% 5.23. ? Let k be a field, and let n > 0 be an integer. Assume that there are
% no irreducible polynomials of degree n in k[x]. Prove that there are no
% separable
% extensions of k of degree n. [\S{}7.1]

\section{A little Galois theory}
\label{sec:fields:a-little-galois-theory}

% Galois theory is a beautiful interplay of field theory and group theory,
% originally
% motivated by the problem of determining `symmetries' among roots of a
% polynomial.
% Galois was interested in concrete relations which must necessarily hold among
% the
% roots of a polynomial, the most trivial example of which being the fact that if
% \alpha , \beta
% are the two roots of x2 + px + q, then \alpha  + \beta  = -p and \alpha \beta  =
% q. Interchanging
% the roots \alpha , \beta  has no effect on the quantities \alpha  + \beta ,
% \alpha \beta . More generally, for a
% higher degree polynomial there may be several quantities which are invariant
% under
% certain permutations of the roots. These quantities and the corresponding groups
% of
% permutations may be viewed as invariants determined by the polynomial, yielding
% a sophisticated tool to study the polynomial.

\subsection{The Galois correspondence and Galois extensions}
\label{subsec:fields:the-galois-correspondence-and-galois-extensions}
% In the language

% of field theory, subsets of the roots of a polynomial f (x) determine
% intermediate
% fields of the splitting field of f (x), and groups of permutations of the roots
% give
% automorphisms of these intermediate fields. With this in mind, it is not
% surprising
% that splitting fields should come to the fore; the fact that one can
% characterize such
% fields in terms of groups of automorphisms (Theorem 6.9 below) will be key to
% the
% whole discussion.
% Before wading into these waters, we should formalize the precise relation be-
% tween groups of automorphisms of an extension and intermediate fields.
% Definition 6.1. Let k \subseteq  F be a field extension, and let G \subseteq
% Autk (F ) be a group
% of automorphisms of the extension. The fixed field of G is the intermediate
% field
% F G := {\alpha  \in  F | \forall g \in  G, g\alpha  = \alpha }.
%  ?
% The fact that F G is indeed a subfield of F containing k is immediate. The
% notion of fixed field allows us to set up a correspondence
% {intermediate fields E: k \subseteq  E \subseteq  F } o
%  / {subgroups of Autk
%  (F )} ,
% sending the intermediate field E to the subgroup AutE (F ) of Autk (F ) and the
% subgroup G \subseteq  Autk (F ) to the fixed field F G .
% Definition 6.2. This is known as the Galois correspondence.
%  ?
% Lemma 6.3. The Galois correspondence is inclusion-reversing.subgroups G of Autk
% (F ) and all intermediate fields k \subseteq  E \subseteq  F :
% \bullet  E \subseteq  F AutE (F ) ;
% \bullet  G \subseteq  AutF G (F ).
% Further, for all
% Further still, denote by E1 E2 the smallest subfield of F containing two
% intermediate
% fields E1 , E2 , and denote by ?G1 , G2 ? the smallest subgroup of Autk (F )
% containing
% two subgroups G 1 , G2 . Then
% \bullet \bullet AutE1E2 (F ) = AutE1 (F ) \cap  AutE2 (F );
% F ?G1 ,G 2 ? = F G 1 \cap  F G2 .
% Proof. Exercise 6.1.
%  ?
% Of course the reader should start wondering whether there are situations guar-
% anteeing that the inclusions appearing in Lemma 6.3 are equalities, making the
% Galois correspondence a bijection. This is precisely where we are heading. The
% first observation is that this is not always the case:
% \sqrt 
%  3Example 6.4. Consider the extension Q \subseteq  Q( 2). Since
% \sqrt 
% [Q( 2) : Q] = 3
% \sqrt 
%  3is prime, the only intermediate
%  fields are Q and Q( 2) (by Corollary
%  1.11). Con-
% \sqrt  3 \sqrt 
%  3 \sqrt 
%  3 \sqrt 
% cerning AutQ (Q( 2)), since 2 \in  R, we have an extension Q( 2) \subseteq  R;
% since
%  2 is
% \sqrt 
% the only cube root
%  \sqrt 
%  3 of
%  in
%  R,
%  we
%  see
%  that
%  the
%  minimal
%  \sqrt 
%  3 polynomial
%  t
%  -
%  of
%  has
%  a
% single root in Q( 2). By Corollary 1.7, AutQ (Q( 2)) consists of a single
% element:
% it is trivial.
% Thus, in this example the Galois correspondence acts between a set with two
% elements and a singleton:
% \sqrt 
%  3 / \sqrt 
% {Q, Q( 2)} o
%  {AutQ
%  (Q( 2))} = {e} .
% In particular, the function associating with each intermediate field the
% correspond-
% ing automorphism group is not injective in general.
%  ?
% This example shows that the inclusion
%  \sqrt 
%  E \subseteq  F AutE (F ) in Lemma 6.3 may be
% proper: Q is properly contained in Q( 2), which is the fixed field of the (only)
% \sqrt  3subgroup of AutQ(Q( 2)). Fortunately, the situation with the other inclusion is
% more constrained:
% Proposition 6.5. Let k \subseteq  F be a finite extension, and let G be a
% subgroup
% of Autk (F ). Then |G| = [F : F G ], and
% G = AutF G (F ).
% In particular, the Galois correspondence (from intermediate fields to
% automorphism
% groups) is surjective for a finite extension.
% Proving this fact leads us to examining very carefully finite extensions of the
% type F G \subseteq  F . Surprisingly, these turn out to satisfy just about every
% property we
% have encountered so far:
% Lemma 6.6. Let k \subseteq  F be a finite extension, and let G be a subgroup of
% Autk (F ).
% Then F G \subseteq  F is a finite, simple, normal, separable extension.
% Remark 6.7. The key to this lemma is the following observation, which is worth
% highlighting.
% If \alpha  \in  F and g \in  G, note that g\alpha  must be a root of the minimal
% polynomial
% of \alpha  over k; there are only finitely many roots, so the G-orbit of \alpha
% consists of finitely
% many elements \alpha  = \alpha 1 , . . . , \alpha n. The group G acts on the
% orbit by permuting its
% elements; therefore, every element of G leaves the polynomial
% q\alpha  (t) = (t - \alpha 1 ) \dots  (t - \alpha n )
% fixed. In other words, the coefficients of this polynomial must be in the fixed
% field F G . Further, q\alpha  (t) is separable since it has distinct roots.
% Finally, note
% that deg q\alpha  (t) \leq  |G|.
%  ?
% With this in mind, we are ready to prove the lemma.
% Proof of Lemma 6.6. The extension F G \subseteq  F is finite because k \subseteq
% F is finite.
% Let \alpha  \in  F ; by the remark following the statement of the lemma, \alpha
% is a root of
% a separable polynomial q\alpha  (t) with coefficients in F G . It follows that
% \alpha  is separable
% over F G ; hence the extension is separable according to Definition 4.12.
% Since F G \subseteq  F is finite and separable, it is simple by Proposition
% 5.19; let \alpha  be
% a generator. The polynomial q\alpha  (t) splits in F , and F is generated over F
% G by the
% roots of q\alpha (t) (indeed, \alpha 1 = \alpha  suffices to generate F over F G
% ); therefore, F is a
% splitting field for q\alpha  (t) over F G according to Definition 4.1. Therefore
% F G \subseteq  F is
% normal by Theorem 4.8, and we are done.
%  ?
% Proposition 6.5 is a consequence of Lemma 6.6:
% Proof of Proposition 6.5. By Lemma 6.3, G is a subgroup of AutF G (F ), and in
% particular
% |G| \leq  | AutF G (F )|.
% In order to verify that G = AutF G , it suffices to prove the converse
% inequality. By
% Lemma 6.6, F = F G(\alpha ) for some \alpha  \in  F ; therefore | AutF G (F )|
% equals the number
% of distinct roots in F of the minimal polynomial of \alpha  over F G (by
% Corollary 1.7).
% With notation as in Remark 6.7, \alpha  is a root of q\alpha  (x) \in  F G [x];
% hence the minimal
% polynomial of \alpha  is a factor of q\alpha  (x). Since by construction the
% number of roots
% of q\alpha  (x) is \leq  |G|, we obtain
% | AutF G (F )| \leq  |G|,
% as needed.
% To verify that [F : F G ] = |G|, just note that [F : F G ] = | AutF G (F )| by
% Corollary 5.20, since F G \subseteq  F is normal and separable by Lemma 6.6.
%  ?
% Remark 6.8. The hypothesis that the extension is finite in Proposition 6.5 is
% necessary: the Galois correspondence is not necessarily surjective in the
% infinite
% case. Not all is lost, though: one can give a suitable topology to the group of
% automorphisms and limit the Galois correspondence to closed subgroups of the
% automorphism group, recovering results such as Proposition 6.5. The reader is
% welcome to explore this further---we will not be able to take this more general
% point of view here.
%  ?
% Theorem 6.9. Let k \subseteq  F be a finite field extension. Then the following
% are
% equivalent:
% (1) F is the splitting field of a separable polynomial f (t) \in  k[t] over k;
% (2) k \subseteq  F is normal and separable;
% (3) | Autk (F )| = [F : k];
% (4) k = F Autk (F ) is the fixed field of Autk (F );
% (5) the Galois correspondence for k \subseteq  F is a bijection;
% (6) k \subseteq  F is separable, and if F \subseteq  K is an algebraic extension
% and \sigma  \in  Autk (K),
% then \sigma (F ) = F .
% Proof. Most of the needed implications have been proven along the way.
% (1) \iff  (2) by Theorem 4.8; (2) \implies  (3) by Corollary 5.20. (3) \iff  (4)
% follows
% from Proposition 6.5, applied to the extension F Autk (F ) \subseteq  F : by
% Proposition 6.5,
% we have
% [F : F Autk (F ) ] = | Autk (F )|;
% since k \subseteq  F Aut k (F ) \subseteq  F , it follows that k = F Autk (F )
% if and only if | Autk (F )| =
% [F : k]. (4) \implies  (2) by Lemma 6.6.
% (2) \iff  (6) holds by Exercise 4.7, since finite separable extensions are
% simple
% (Proposition 5.19).
% To prove (5) \implies  (4), let E = F Autk (F ) ; by Proposition 6.5, AutE (F )
% =
% Autk (F ). Assuming that the Galois correspondence is bijective, it follows that
% k = E = F Autk (F ) , giving (4).
% Finally, we prove that (1) \implies  (5). Since k \subseteq  F is a finite
% extension, we
% already know that the Galois correspondence from intermediate fields to
% subgroups
% of Autk (F ) has a right inverse (Proposition 6.5), so it suffices to show it
% has a left-
% inverse. Therefore, it suffices to verify that every intermediate field E equals
% the
% fixed field of the corresponding subgroup AutE (F ). Then let E be an
% intermediate
% field. By (1), F is the splitting field of a separable polynomial f (t) \in
% k[t] \subseteq  E[t];
% therefore, condition (1) holds for the extension E \subseteq  F . As we have
% already proved
% that (1) \iff  (4), this implies that E = F AutE (F ), and we are done.
%  ?
% Definition 6.10. A finite extension k \subseteq  F is Galois if it satisfies any
% of the
% conditions listed in Theorem 6.9.
%  ?
% Warning: We will not study the Galois condition for infinite extensions (cf. Re-
% mark 6.8). Thus Galois extensions will implicitly be finite in what follows.
% Condition (6) is technically useful and helps with visualizing the Galois condi-
% tion. If a finite separable extension k \subseteq  F is not Galois, then F can
% be embedded
% in some larger extension k \subseteq  K (for example, in the algebraic closure k
% \subseteq  k) in
% many possible ways:
% K
% k
% \sigma (F )
% F
% If k \subseteq  F is Galois, all these images must coincide:
% K
% F = \sigma (F )
% k
% Of course there are possibly still many ways21 to embed F in K, but they all
% have
% the same image F .
% The mysterious comments at the end of Example 1.4 are now hopefully clear.
% Galois extensions of a field k are well-defined subfields of k, preserved by
% auto-
% morphisms
%  \sqrt  3 of k over k. Contrast this situation with the non-Galois extension
% Q \subseteq  Q( 2), which admits three different embeddings
%  in Q. For example, it does
% \sqrt  \sqrt 
%  3not make sense to talk about `the' composite of Q( 2) and Q( 2), as a subfield
% \sqrt 
%  3of (for instance) C: according to the chosen embedding of Q( 2), this may or
%  may
% 21 In fact, precisely [F : k]s
%  = [F : k] = | Autk
%  (F )| if K = k and hence for all K \supseteq  F .
% not be a subfield of R. But 22 it does make sense to talk about the composite of
% two
% Galois extensions k \subseteq  F1 , k \subseteq  F2 , since both F1 and F2 are
% well-defined as subfield
% of k, so their composite F1 F2 is independent of the choice of embeddings F1
% \subseteq  k,
% F2 \subseteq  k.
% Definition 6.11. If k \subseteq  F is a Galois extension, the corresponding
% automorphism
% group Autk (F ) is called the Galois group of the extension.
%  ?
% \sqrt 
%  3To reiterate, the extension Q \subseteq 
%  Q( 2) is not a Galois extension: as we have
% \sqrt 
%  3observed in Example 6.4, AutQ (Q( 2)) is trivial (hence, for example,
%  condition (3)
% of Theorem 6.9 does not hold).
% On the other hand, we have already run into several examples of Galois exten-
% sions: the examples of splitting fields we have seen in \S{}4.1 are all Galois
% extensions.
% The `Galois fields' of \S{}5.1 are (surprise, surprise) Galois extensions of
% their prime
% subfields, by Theorem 5.1. The cyclotomic fields (\S{}5.2) are Galois extensions
% of Q.
% A Galois extension k \subseteq  F is cyclic, abelian, etc., if the corresponding
% group
% of automorphisms Autk (F ) is cyclic, abelian, etc. For example, Q \subseteq
% Q(\zeta 8) is an
% abelian, but not cyclic, Galois extension (cf. Proposition 5.16 and Example
% 5.17).

\subsection{The fundamental theorem of Galois theory, I}
\label{subsec:fields:the-fundamental-theorem-of-galois-theory-i}
% The fundamental theo-

% rem of Galois theory amounts to a more complete description of the Galois corre-
% spondence for Galois extensions.
% At this stage, this is what we know:
% \bullet  If k \subseteq  F is a (finite) Galois extension, then there is an inclusion-reversing
% bijection
% /{intermediate fields E: k \subseteq  E \subseteq  F } o
%  {subgroups of Autk
%  (F )} :
% through this bijection, the intermediate field E corresponds to the subgroup
% AutE (F ) of Autk (F ), and the subgroup G \subseteq  Autk (F ) corresponds to
% the fixed
% field F G .
% \bullet  For every intermediate field E,
% [F : E] = | Aut E (F )|.
% The extension E \subseteq  F is then also a Galois extension (Exercise 6.3), and
% [E : k] = [Autk (F ) : AutE (F )] :
% this follows from Lagrange's theorem (Corollary II.8.14) and its field theory
% counterpart (Proposition 1.10).
% \sqrt 
%  3\bullet  The extension k \subseteq  E is not necessarily Galois. (Once more Q(
%  2) gives a
% counterexample, since it can be viewed as intermediate in the splitting field
% of t3 - 2 over Q.)
% The fact that the Galois correspondence is a bijection may be upgraded as
% follows:
% 22 `Separability' does not play a role in these considerations; thus they best
% convey `Galoisness'
% in, say, characteristic 0.
% Theorem 6.12. Let k \subseteq  F be a Galois extension. The Galois
% correspondence is
% an inclusion-reversing isomorphism of the lattice of intermediate subfields of k
% \subseteq  F
% with the lattice of subgroups of Autk (F ).
% That is (with notation as in Lemma 6.3), if E1 , E2 are intermediate fields
% and G1, G2 are the corresponding subgroups of Autk (F ), then E1 \cap  E2
% corresponds
% to ?G1 , G2? and E1E2 corresponds to G 1 \cap  G2 .
% Proof. This follows immediately from Theorem 6.9 and Lemma 6.3, which gives
% AutE1 E2 (F ) = G1 \cap  G2 ,
%  F ?G1 ,G2 ? = E1 \cap  E2
% as needed.
%  ?
% In this context, it is common to use the `vertical' depiction of field
% extensions,
% with a parallel notation for subgroups:
% F
%  {e}
% E
%  AutE (F )
% k
%  Autk (F )
% The content of Theorem 6.12 is that the lattice of intermediate fields of a
% Galois
% extension k \subseteq  F and the lattice of subgroups of the corresponding
% Galois group are
% identical.
% Example 6.13. For finite fields, this coincidence of lattices was essentially
% proven
% `by hand' in Corollary 5.4 and Proposition 5.8. For example, the extension F2
% \subseteq  F64
% is Galois, with cyclic Galois group C6, generated by the Frobenius automorphism
% \phi .
% The lattices are
% F8 0
%  0 |
%  |
%  || 0 0
%  F64 0
%  }}
%  0 0
%  }
%  }
%  0 0
%  F4
%  ,
%  {e, \phi 3}
%  H H
%  jjj
%  H H
%  H H
%  j
%  H j
%  H
%  j
%  H jjjj
%  H
%  H H
%  H
%  {e} J
%  J J J
%  J iiii
%  J J
%  J J i
%  J
%  i
%  J i
%  J J
%  {e, J J
%  \phi 2 , \phi 4}
% F2
%  {e, \phi , \phi 2
%  , \phi 3
%  , \phi 4
%  , \phi 5
%  }
% ?
% Theorem 6.12 is often used in the `groups to fields' direction: finding the
% lattice
% of subgroups of a finite group is essentially a combinatorial problem, and
% through
% the Galois correspondence it determines the lattice of intermediate fields of a
% cor-
% responding Galois extension.
% \sqrt  \sqrt 
%  \sqrt 
%  \sqrt 
% Example 6.14. The extension Q( 2, 3) = Q( 2 + 3) studied in Example 1.19
% is the splitting field of the polynomial t4 - 10t2 + 1, so it is Galois. We
% found that
% its Galois group is Z/2Z \times  Z/2Z; the lattice of this group has no
% mysteries for us:
% ?(0, 0)?
% ooo
%  o
%  O O O
%  O O O
%  O
% ooo
% ?(1, 0)?
%  ?(1, 1)?
%  ?(0, 1)?
% O O O
%  O O O
%  O
%  ooo
%  o
%  o
%  o
%  o
% Z/2Z \times  Z/2Z
% and therefore the lattice of intermediate fields is just as transparent:
% \sqrt  \sqrt 
% Q( 2, 3)
% \sqrt 
%  r
%  rrr
%  rr
%  \sqrt 
%  L L L
%  L L L
%  \sqrt 
% Q( 2)
%  Q( 6)
%  Q( 3)
% M M M
%  M M M
%  M M rrr
% rrrrr
% Q
% (The intermediate fields are determined by recalling the \sqrt  generators
%  \sqrt  of the corre-
%  \sqrt 
% sponding subgroups, as found in Example 1.19: the flip 3 ?\to  - 3 leaves Q( 2)
% fixed, etc.) Galois theory tells us that there is no other intermediate field.
%  ?

\subsection{The fundamental theorem of Galois theory, II}
\label{subsec:fields:the-fundamental-theorem-of-galois-theory-ii}
% Another popular no-

% tational device is to mark an extension by the corresponding Galois group, if
% the
% extension is Galois:
% F
% AutE (F )
% Autk (F )
%  E
% k
% Example 6.4 shows that the extension k \subseteq  E need not be a Galois
% extension:
% \sqrt  3Q( 2) is an intermediate field for the extension of Q into the splitting field for the
% polynomial t3 - 2 (which is Galois, by part (1) of Theorem 6.9), but as we have
% seen, it is not Galois.
%  \sqrt 
%  3 The splitting field has degree 6 over Q and is a quadratic
% extension of Q( 2). This is in fact the typical situation: every finite
% separable
% extension may be enlarged to a Galois extension (Exercise 6.4).
% In any case, the question of whether the extension of the base field k in an
% intermediate field E of a Galois extension k \subseteq  F is Galois is central
% to the theory.
% The outcome of the discussion will be very neat, and this is possibly the most
% striking part of the fundamental theorem of Galois theory:
% Theorem 6.15. Let k \subseteq  F be a Galois extension, and let E be an
% intermediate
% field. Then k \subseteq  E is Galois if and only if AutE (F ) is normal in Autk
% (F ); in this
% case, there is an isomorphism
% Autk (E) \sim 
%  =
%  Autk (F )
%  .
% AutE (F )
% \sqrt 
%  3The `problem' with Q \subseteq  Q( 2) is that the automorphism group of the
%  splitting
% field of t3 - 2 is S3 (as we will
%  soon be able to verify: cf. Example 7.19). The
% \sqrt  3subgroup corresponding to Q( 2) has index 3; hence (as we know, since we are
% well acquainted with S3 ) it is not normal.
% The fact that the Galois group of k \subseteq  E will turn out to be isomorphic
% to the
% quotient Autk (F )/ AutE (F ) in the Galois case should not be too surprising,
% since
% we know already that in any case [E : k] equals the index of AutE (F ) in Autk
% (F ),
% and the index equals the order of the quotient if the subgroup is normal. In
% general,
% [E : k] equals the number of cosets of AutE (F ) in Autk (F ), and it is worth
% shooting
% for a concrete interpretation of these cosets, even when k \subseteq  E does not
% turn out to
% be a Galois extension.
% We already know that the number of left-cosets of AutE (F ) equals the index
% of AutE (F ) in Autk (F ), hence [E : k], hence [E : k]s (because k \subseteq  E
% is separable,
% as a subextension of a separable extension), and hence there must exist a
% bijection
% between the set of left-cosets g AutE (F ) as g ranges in Autk (F ) and the set
% of
% different embeddings \iota  : E \to  k in an algebraic closure of k, extending
% idk . Not
% surprisingly, there is a reason for this---that is, a `meaningful' bijection
% between
% these two sets.
% This is an excellent moment to review the material on group actions we covered
% in \S{}II.9, especially \S{}II.9.3. Recall in particular that we were able to
% classify all
% transitive actions of a group G on a set S (Proposition II.9.9): these G-sets
% are all
% isomorphic to the set of left-cosets of a subgroup H of G under the natural
% action
% of G; and H may be taken to be the stabilizer of any element of S.
% We are going to define a natural action of Autk (F ) on the set I of embeddings
% of E in k extending the identity, and prove that, as an Autk (F )-set, I is
% isomorphic
% to the set of left-cosets of AutE (F ) in Autk (F ). This will be our
% `meaningful
% bijection', and after this is done Theorem 6.15 will look perfectly reasonable.
% Fix an embedding F \subseteq  k of F in an algebraic closure of k, and let I be
% the
% set of embeddings \iota  : E \setminus \to  k extending the identity on k. It is
% a good idea to
% view k \subseteq  E as an abstract extension, rather than identifying E with a
% specific
% intermediate field in F ; the fact that we can view E as an intermediate field
% of the
% extension k \subseteq  F simply means that \iota (E) \subseteq  F for some \iota
% \in  I.
% \bullet  For all \iota  \in  I, \iota (E) \subseteq  F .
% Indeed, k \subseteq  E is simple (finite and separable \implies  simple,
% Proposition 5.19). If
% E = k(\alpha ), then \iota  is determined by \iota (\alpha ), which is
% necessarily a root of the minimal
% polynomial of \alpha  over k. But k \subseteq  F is normal and contains some
% root \iota (\alpha ) of this
% polynomial; hence F contains all of them (Definition 4.7). That is, \iota
% (\alpha ) \in  F for all
% \iota  \in  I, implying \iota (k(\alpha )) \subseteq  F for all \iota  \in  I.
% \bullet  Let \iota  \in  I, and let g \in  Autk (F ); then g \circ  \iota  \in  I. Therefore, Autk (F ) acts on I.
% Here g \circ  \iota  is interpreted in the following way: \iota (E) \subseteq  F
% , as we have seen; hence
% g restricts to a homomorphism \iota (E) \to  F , and g \circ  \iota  is the
% composition of this
% homomorphism with \iota .
% \bullet  The action of Autk (F ) on I is transitive.
% This follows from the `uniqueness of splitting fields', Lemma 4.2. Indeed, let
% \iota 1 , \iota 2 \in
% I. Both \iota 1 (E), \iota 2 (E) contain k, and F is a splitting field over k
% (part (1) of
% Theorem 6.9); hence F is a splitting field (for the same polynomial) over both
% \iota 1 (E)
% and \iota 2(E). By Lemma 4.2, there exists an automorphism g : F \to  F
% extending the
% isomorphism
% \iota 2 \circ  \iota -1
%  1 : \iota 1 (E) \to  \iota 2 (E).
% Then g \in  Autk (F ) (since \iota 2 \circ  \iota 1 -1
%  restricts to the identity on k), and the fact that
% g restricts to \iota 2 \circ  \iota -1
%  1 means precisely that g \circ  \iota 1 = \iota 2.
% Now choose one \iota  \in  I, and use it to identify E with an intermediate
% field
% of k \subseteq  F . Then AutE (F ) = Aut\iota (E)(F ) is the subgroup of Autk (F
% ) consisting of
% those elements g which restrict to the identity on E = \iota (E), that is, the
% stabilizer
% of \iota . At this point, Proposition II.9.9 is all that is needed to conclude
% the promised
% result:
% \bullet  As an Autk (F )-set, I is isomorphic to the set of left-cosets of AutE (F ) in
% Autk (F ).
% I view this observation as an integral part of the fundamental theorem of Galois
% theory. The conventional statement (Theorem 6.15) is an essentially immediate
% consequence of this fact.
% Proof of Theorem 6.15. As observed above, the Autk (F )-stabilizer of \iota  \in
% I
% equals Aut\iota (E) (F ). By Proposition II.9.12, stabilizers of different \iota
% are conjugates
% of each other; if AutE (F ) is normal, it follows that for all \iota  \in  I
% Aut\iota (E) (F ) = AutE (F ).
% This implies that \iota (E) = E for all \iota  \in  I, since the Galois
% correspondence is bijective
% for the Galois extension k \subseteq  F . It follows that k \subseteq  E
% satisfies condition (6) of
% Theorem 6.9; hence it is Galois.
% Conversely, assume that k \subseteq  E is Galois; by the same condition (6),
% \iota (E) = E
% for all \iota  \in  I. Restriction to E = \iota (E) then defines a homomorphism
% \rho  : Autk (F )
%  / Aut k (E).
% This homomorphism is surjective (because the Autk (F )-action on I is
% transitive),
% and its kernel consists of those g \in  Autk (F ) which restrict to the identity
% on E,
% that is, precisely of AutE (F ). This shows that AutE (F ) is normal and
% establishes
% the stated isomorphism by virtue of the `first isomorphism theorem' for groups,
% Corollary II.8.2.
%  ?
% Remark 6.16. There is a parallel between Galois theory and the theory of
% covering
% spaces in topology. In this analogy, Galois extensions correspond to regular
% covers;
% the Galois group of an extension corresponds to the group of deck
% transformations;
% and Theorem 6.15 corresponds to the fact that the quotient of a regular cover by
% a
% normal subgroup of the group of deck transformations is again a regular cover.
% More general (connected) covers correspond to more general algebraic exten-
% sions. A space is simply connected if and only if it admits no nontrivial
% connected
% covers, so this notion corresponds in field theory to the condition that a field
% K
% admits no nontrivial algebraic extensions, that is, that K is algebraically
% closed.
% Viewing the fundamental group of a space as the group of deck transformations
% of its fundamental cover suggests that we should think of the Galois group of
% the
% algebraic closure23 k \subseteq  k as `the fundamental group' of a field k.
% In algebraic geometry this analogy is carried out to its natural consequences.
% A covering map of algebraic varieties X \to  Y determines a field extension K(Y
% ) \subseteq
% K(X), where K(X), K(Y ) are the `fields of rational functions' (in the affine
% case,
% these are just the fields of fractions of the corresponding coordinate rings;
% cf. Ex-
% ercise 2.16). One can then use Galois theory to transfer to the
% algebro-geometric
% environment notions such as the fundamental group, without appealing to topolog-
% ical notions (such as `continuous maps from S 1 '), which would be problematic
% in,
% e.g., positive characteristic.
%  ?

\subsection{Further remarks and examples}
\label{subsec:fields:further-remarks-and-examples}
% Suppose k \subseteq  F and k \subseteq  K are two finite

% extensions contained in a larger extension of k (for example, we could choose
% specific
% embeddings F \subseteq  k, K \subseteq  k in an algebraic closure of k). Then we
% can consider the
% composite KF (cf. Lemma 6.3) of F , K in the larger extension. It is natural to
% question how the Galois condition behaves under this operation.
% Proposition 6.17. Suppose k \subseteq  F is a Galois extension and k \subseteq
% K is any finite
% extension. Then K \subseteq  KF is a Galois extension, and AutK (KF ) \sim 
%  = AutF \cap K (F ).
% Pictorially,
% sss
%  s
% KF L L L G
%  L
% F K K
%  K
%  rrr
%  K
% G
%  F \cap  K
% k
% Proof. As k \subseteq  F is Galois, it is the splitting field of a separable
% polynomial
% f (x) \in  k[x] \subseteq  K[x]. The roots of f (x) generate F over k, so they
% generate KF
% over K; in other words, KF is the splitting field of f (x) over K, so K
% \subseteq  KF is
% Galois.
% If \sigma  : KF \to  KF is an automorphism extending the identity on K (and
% hence
% on k), then \sigma  restricts to a homomorphism F \to  \sigma (F ) extending the
% identity on k.
% Since k \subseteq  F is Galois, \sigma (F ) = F ((6) in Theorem 6.9). Thus,
% `restriction' defines
% a natural group homomorphism
% \rho  : AutK (KF ) -\to  Aut k (F ).
% I claim that \rho  is injective. Indeed, by definition every \sigma  \in  AutK
% (KF ) is the identity
% on K; if \sigma  restricts to the identity on F , then \sigma  is the identity
% on the composite KF ;
% that is, \sigma  = id in AutK (KF ).
% Therefore, AutK (KF ) may be identified with a subgroup G of Autk (F ). Since
% every \sigma  \in  AutK (KF ) fixes K, the restriction \rho (\sigma ) fixes F
% \cap K; that is, F G \supseteq  F \cap K.
% 23 Of course the extension k \subseteq  k is not finite in general, so one needs
% the infinite version of
% Galois theory; cf. Remark 6.8.
% Conversely, suppose \alpha  \in  F G ; then K(\alpha ) is in the fixed field of
% AutK (KF ), which
% is K itself since K \subseteq  KF is Galois (condition (4) of Theorem 6.9). That
% is, \alpha  \in  K,
% and it follows that F G = F \cap  K. By Proposition 6.5, G = AutF \cap K (F ),
% concluding
% the proof.
%  ?
% I will leave to the reader the pleasure of verifying that if both k \subseteq  F
% and k \subseteq  K
% are Galois, then KF is also Galois over k, and so is F \cap  K (Exercise 6.13).
% Example 6.18. We have studied cyclotomic fields Q(\zeta n) as extensions of Q;
% Q(\zeta n)
% is the splitting field of xn - 1, so these extensions are Galois; we have proved
% (Proposition 5.16) that AutQ (Q(\zeta n)) is isomorphic to the group of units in
% Z/nZ.
% Now let k be any field of characteristic zero. The splitting field of xn - 1
% over k
% is the composite k(\zeta ) of k and Q(\zeta ). By Proposition 6.17 the extension
% k \subseteq  k(\zeta ) is
% Galois, and Autk (k(\zeta )) is isomorphic to a subgroup of the group of units
% of Z/nZ.
% In particular, it is abelian.
%  ?
% The previous example gives an important class of abelian Galois extensions.
% Remarkably, we can classify all cyclic Galois extensions, if we impose some
% restric-
% tion on the base field k.
% Proposition 6.19. Let k \subseteq  F be an extension of degree m. Assume that k
% contains
% a primitive m-th root of 1 and char k does not divide24 m. Then k \subseteq  F
% is Galois
% and cyclic if and only if F = k(\delta ), with \delta m \in  k.
% That is, under the given hypotheses m \sqrt 
%  on k, the cyclic Galois extensions of k
% are precisely those of the form k \subseteq  k( c), with c \in  k. This is part
% of Kummer
% theory, a basic tool in the study of abelian extensions. Proposition 6.19 will
% be a
% key ingredient in our application of Galois theory to the solvability of
% polynomial
% equations.
% Proof. Let \zeta  \in  k be a primitive m-th root of 1.
% First assume that F = k(\delta ), with \delta m = c \in  k.polynomial xm - c,
% \delta  , \zeta \delta  , \zeta  2\delta  , \dots  , \zeta  m-1 \delta ,
% Then all m roots of the
% are in F , and F is generated by them; therefore F is the splitting field of the
% separa-
% ble (since char k does not divide m) polynomial xm - c over k; hence (Theorem
% 6.9,
% part (1)) k \subseteq  F is Galois.
% To see that Autk (F ) is cyclic, note that every \phi  \in  Autk (F ) is
% determined
% by \phi (\delta ), which is necessarily a root of xm - c; therefore \phi (\delta
% ) = \zeta  i \delta  for some
% i, determined up to a multiple of m. This defines an isomorphism of Autk (F )
% with Z/mZ, as is immediately verified.
% Conversely, assume that k \subseteq  F is Galois and Autk (F ) is cyclic,
% generated by \phi .
% The automorphisms \phi  i , i = 0, . . . , m-1, are pairwise distinct; thus they
% are linearly
% independent over F (Exercise 6.14). Therefore, there exists an \alpha  \in  F
% such that25
% \delta  := \alpha  + \zeta  -1 \phi (\alpha ) + \dots  + \zeta  -(m-2)\phi m-2 (\alpha ) + \zeta  -(m-1) \phi m-1 (\alpha ) = 0.
% 24 We will apply this result in \S{}7.4; for convenience, we will take char k =
% 0 in that application.
% 25This is called a Lagrange resolvent.
% Note that
% \phi (\delta ) = \phi (\alpha ) + \zeta  -1 \phi 2 (\alpha ) + \dots  + \zeta  -(m-2)\phi m-1 (\alpha ) + \zeta  -(m-1) \alpha  = \zeta \delta 
% (since \zeta  \in  k, so \phi (\zeta ) = \zeta ). This implies that \delta  is
% not fixed by any nonidentity
% element of Autk (F ); that is, Autk(\delta ) (F ) = {e}; that is, k(\delta ) = F
% . Further,
% \phi (\delta  m ) = (\phi (\delta ))m = \zeta  m\delta m = \delta  m ;
% that is, \delta  m is fixed by \phi , and hence by the whole of Autk (F ). Since
% k \subseteq  F is
% Galois, this implies \delta  m \in  k (Theorem 6.9, part (4)). That is, \delta
% satisfies the stated
% conditions.
%  ?

\subsection*{Exercises}

% 6.1. ? Prove Lemma 6.3. [\S{}6.1]
% 6.2. Prove that quadratic extensions in characteristic= 2 are Galois.
% 6.3. ? Let k \subseteq  F be a Galois extension, and let E be an intermediate
% field. Prove
% that E \subseteq  F is a Galois extension. [\S{}6.2]
% 6.4. ? Let k \subseteq  E be a finite separable extension. Prove that E may be
% identified
% with an intermediate field of a Galois extension k \subseteq  F of k.
% In fact, prove that there is a smallest such extension k \subseteq  F , in the
% sense that
% if k \subseteq  E \subseteq  K, with k \subseteq  K Galois, then there exists an
% embedding of F in K which
% is the identity on E. (The extension k \subseteq  F is the Galois closure of the
% extension
% k \subseteq  E. It is clearly uniquely determined up to isomorphism.) [\S{}6.3,
% 6.5]
% 6.5. We have proved (Proposition 5.19) that all finite separable extensions k
% \subseteq  E
% are simple. Let k \subseteq  F be a Galois closure of k \subseteq  E (Exercise
% 6.4). For \alpha  \in  E,
% prove that E = k(\alpha ) if and only if \alpha  is moved by all \sigma  \in
% Autk (F ) \setminus  AutE (F ).
% *
%  \sqrt 
% 6.6. \bullet  Prove that Q \subseteq  Q( 2 + 2) is Galois, with cyclic Galois
% group. (Cf. Ex-
% ercise 1.25.)
% *
%  \sqrt 
% \bullet  Prove that Q \subseteq  Q( 3 + 5) is Galois and its Galois group is isomorphic to
% (Z/2Z) \times  (Z/2Z).
% *
%  \sqrt 
% \bullet  Prove that Q \subseteq  Q( 1 + 2) is not Galois, and compute its Galois closure
% Q \subseteq  F . Prove that AutQ (F ) \sim 
%  = D8 . (Use Exercise IV.2.16.)
% 6.7. Let p > 0 be prime, and let d | e be positive integers, so that there is an
% extension Fpd \subseteq  Fpe . Prove that AutFpd Fpe is cyclic, and describe a
% generator of
% this group. (This generalizes Proposition 5.8; use Galois theory.)
% 6.8. Let k \subseteq  F be a Galois extension of degree n, and let E be an
% intermediate
% field. Assume that [E : k] is the smallest prime dividing n. Prove k \subseteq
% E is Galois.
% 6.9. Let k \subseteq  F be a Galois extension of degree 75. Prove that there
% exists an
% intermediate field E, with k ? E ? F , such that the extension k \subseteq  E is
% Galois.
% 6.10. Let k \subseteq  F be a Galois extension and E an intermediate field.
% Prove that the
% normalizer of AutE (F ) in Autk (F ) is the set of \sigma  \in  Autk (F ) such
% that \sigma (E) \subseteq  E.
% Use this to give an alternative proof of the fact that E is Galois over k if and
% only if AutE (F ) is normal in Autk (F ).
% 6.11. Let k \subseteq  E and E \subseteq  F be Galois extensions.
% \bullet  Find an example showing that k \subseteq  F is not necessarily Galois.
% \bullet  Prove that if every \sigma  \in  Autk (E) is the restriction of an element of Autk (F ),
% then k \subseteq  F is Galois.
% 6.12. Find two algebraic extensions k \subseteq  F , k \subseteq  K and
% embeddings F \subseteq  k,
% \sigma 1 : K \subseteq  k, \sigma 2 : K \subseteq  k extending k \subseteq  k, such that the composites F \sigma 1 (K),
% F \sigma 2(K) are not isomorphic.
% Prove that no such example exists if F and K are Galois over k.
% 6.13. ? Let k \subseteq  F and k \subseteq  K be Galois extensions, and assume F
% and K are
% subfields of a larger field. Prove that k \subseteq  F K and k \subseteq  F \cap
% K are both Galois
% extensions. [\S{}6.4, \S{}7.4]
% 6.14. ? Let k \subseteq  F be a field extension, and let \phi 1 , . . . , \phi m
% \in  Autk (F ) be pairwise
% distinct automorphisms. Prove that \phi 1, . . . , \phi m are linearly
% independent over F .
% (Recycle/adapt the hint for Exercise VI.6.15.) [\S{}6.4, 6.16, \S{}IX.7.6]
% 6.15. Let k \subseteq  F be a Galois extension, and let \alpha  \in  F . Prove
% that
% ?
%  ?
% Nk\subseteq F (\alpha ) =
%  \sigma (\alpha ), trk\subseteq F (\alpha ) =
%  \sigma (\alpha ).
% \sigma \in Autk (F )
%  \sigma \in Autk (F )
% (Exercise 4.19.)
% 6.16. ? Let k \subseteq  F be a cyclic Galois extension of degree d, and let
% \phi  be a generator
% of Autk (F ). Let \alpha  \in  F be an element such that Nk\subseteq F (\alpha )
% = 1.
% \bullet  Prove that the automorphisms idF , \phi , . . . , \phi d-1 are linearly independent
% over F . (Exercise 6.14.)
% \bullet  Prove that there exists a \gamma  \in  F such that
% \beta  := \gamma  + \alpha \phi (\gamma ) + \alpha \phi (\alpha )\phi 2 (\gamma ) + \dots  + \alpha \phi (\alpha ) \dots  \phi d-2 (\alpha )\phi d-1 (\gamma ) = 0.
% \bullet  Prove that \alpha \phi (\alpha )\phi 2(\alpha ) \dots  \phi d-1 (\alpha )\phi d (\gamma ) = \gamma , and deduce that \alpha  = \beta /\phi (\beta ).
% Together with the result of Exercise 4.20, the conclusion is that an element
% \alpha  of
% a cyclic Galois extension as above has norm 1 if and only if there exists a
% \beta  such
% that \alpha  = \beta /\phi (\beta ).
% This is Hilbert's theorem 90 (the 90-th theorem in Hilbert's Zahlbericht, a
% report
% on the state of number theory at the end of the nineteenth century commissioned
% by the German Mathematical Society). [6.17, \S{}IX.7.6, IX.7.18]
% 6.17. Exercise 6.16 should remind the reader of the proof of Proposition 6.19,
% and
% for good reasons. Assume that k \subseteq  F is a cyclic Galois extension of
% degree m and
% that k contains a primitive m-th root \zeta  of 1. Use Hilbert's theorem 90 to
% prove
% that F = k(\delta ) for a \delta  such that \delta m \in  k, thereby recovering
% one direction of the
% statement of Proposition 6.19. (Hint: What is Nk\subseteq F (\zeta  -1)?)
% The case of Hilbert's theorem 90 needed here is due to Kummer.
% 6.18. Use Hilbert's theorem 90 to find all rational roots a, b of the equation
% a2 +
% b2 d = 1, where d \sqrt 
%  is a positive integer that is not a square. (Hint: a2 + b2d equals
% NQ\subseteq Q(\sqrt 
% -d) (a + b -d); cf. Exercise 1.12.)
% 6.19. \lnot  Let k \subseteq  F be a cyclic Galois extension of degree d, and
% let \phi  be a generator
% of Autk (F ). Let \alpha  \in  F be an element such that trk\subseteq F (\alpha
% ) = 0. Prove an `additive
% version' of Hilbert's theorem 90:
% \bullet  Prove that there exists a \gamma  \in  F such that trk\subseteq F (\gamma ) = 0.
% \bullet  Stare at the expression
% \alpha \phi (\gamma ) + (\alpha  + \phi (\alpha ))\phi 2 (\gamma ) + \dots  + (\alpha  + \phi (\alpha ) + \dots  + \phi d-2 (\alpha ))\phi d-1 (\gamma ).
% \bullet  Prove that there exists a \beta  \in  F such that \alpha  = \beta  - \phi (\beta ).
% Together with Exercise 4.20, this says that an element \alpha  of a cyclic
% Galois extension
% as above has trace 0 if and only if there exists a \beta  such that \alpha  =
% \beta  - \phi (\beta ). For an
% application, see Exercise 7.15. [7.15]

\section{Short march through applications of Galois theory}
\label{sec:fields:short-march-through-applications-of-galois-theory}

% Galois theory is a pervasive tool, with wide-ranging applications. I have
% neither
% the competence nor room in this book for an adequate survey of applications of
% the
% theory, but even the few examples reviewed in this section should suffice to
% convey
% its considerable power.

\subsection{Fundamental theorem of algebra}
\label{subsec:fields:fundamental-theorem-of-algebra}
% Theorem 6.15 and a little group theory

% are essentially all that is needed26 to give an algebraic proof of the
% fundamental
% theorem of algebra, promised back in \S{}V.5.3.
% Theorem 7.1. C is algebraically closed.
% Proof. Let f (x) \in  C[x] be a nonconstant polynomial; we have to prove that f
% (x)
% has roots in C. Note that if f (x) has no roots in C, then neither does f (x)f
% (x) \in
% R[x]; that is, we may assume that f (x) has real coefficients. Let F be a
% splitting
% field for f (x) over R; embed F in an algebraic closure of R (we `don't know
% yet'
% that C is algebraically closed!), and consider the extension
% R \subseteq  F (i).
% This extension is Galois: it is the splitting field of the square-free part27 of
% f (x)(x2 +
% 1). Let G = AutR(F (i)).
% 26 I am cheating a little, actually, since the intermediate value theorem will
% also be used behind
% the scenes. This is not too surprising: the completeness of R must enter the
% proof somewhere.
% 27 That is, take each irreducible factor of the polynomial with a power of
% exactly 1. Over a
% perfect field, square-free polynomials are clearly separable.
% By Sylow I (Theorem IV.2.5), G has a 2-Sylow subgroup H. The index [G : H]
% is odd, so it corresponds to a finite, separable extension R \subseteq  E with
% [E : R] odd.
% However, every polynomial of odd degree over R has a real root (Exercise V.5.17,
% which only uses elementary calculus), and it follows that R \subseteq  E is
% trivial (Exer-
% cise 5.23).
% Therefore [G : H] = 1, proving that G is a 2-group. Since C = R(i) \subseteq  F
% (i),
% the Galois group of the (Galois) extension
% C \subseteq  F (i)
% is a subgroup of G; therefore it is a 2-group: | AutC (F (i))| = 2n for some n
% \geq  0.
% Now recall (Proposition IV.2.6) that every group of order pn , with p prime,
% contains subgroups of order pm for all 0 \leq  m \leq  n. In particular, if n
% \geq  1, then
% AutC (F (i)) contains a subgroup of order 2n-1, which corresponds to a quadratic
% extension of C via the Galois correspondence. However, there are no quadratic
% extensions of C, because there are no irreducible quadratic polynomials in C[x]
% (Exercise 5.23 again).
% This contradiction shows that n = 0; that is, AutC (F (i)) is necessarily
% trivial.
% This proves that F (i) = C and in particular that C contains the roots of f (x).
% ?
% The argument can be simplified further, bypassing the use of Sylow's theorem
% at the price of a slight increase in length. But why should we attempt to bypass
% Sylow's theorem, after the effort put into proving it in Chapter IV?

\subsection{Constructibility of regular n-gons}
\label{subsec:fields:constructibility-of-regular-n-gons}
% We return one last time to the issue

% of constructing regular n-gons. Simpler considerations have given us constraints
% on n that must necessarily be fulfilled for the regular n-gon to be
% constructible: if
% n = p is prime, we have found that p - 1 must be a power of 2 (\S{}3.3); this
% was
% upgraded for any n to the condition that \varphi (n) is a power of 2, in \S{}5.2
% (and an
% even more explicit condition is given in Exercise 5.19). But these are all
% `negative'
% results: we have not as yet proved that if \varphi (n) is a power of 2, then the
% regular
% n-gon is constructible. This is where Galois theory comes into the picture.
% Proposition 7.2. Let k \subseteq  F be a Galois extension, and assume [F : k] =
% pr for
% some prime p and r \geq  0. Then there exist intermediate fields
% k = E0 \subseteq  E 1 \subseteq  E2 \subseteq  \dots  \subseteq  Er = F
% such that [Ei : Ei-1 ] = p for i = 1, . . . , r.
% Proof. As the Galois correspondence is bijective for Galois extensions (Theo-
% rem 6.9, part (5)), this statement follows immediately from the fact that a
% group
% of order pr , with p prime, has a complete series of p-subgroups; cf. for
% example the
% discussion following the statement of Theorem IV.2.8.
%  ?
% Theorem 7.3. The regular n-gon is constructible by straightedge and compass if
% and only if \varphi (n) is a power of 2.
% Proof. As recalled above, we have already established the \implies  direction.
% For the converse, assume \varphi (n) = 2r for some r. The extension Q \subseteq
% Q(\zeta n)
% is Galois (it is the splitting field of \Phi n (x)), of order [Q(\zeta n ) : Q]
% = \varphi (n) = 2r
% (Proposition 5.14). Proposition 7.2 shows that the condition given in Theorem
% 3.4
% is satisfied and therefore that \zeta n is constructible, as needed.
%  ?
% For example, \varphi (17) = 16 = 24; the Galois group of Q(\zeta 17 ) over Q is
% isomorphic
% to (Z/17Z)\ast  (by Proposition 5.16), and therefore (Theorem IV.6.10) it is
% cyclic, of
% order 16. A generator of (Z/17Z)\ast  is [6]17 ; it follows that AutQ (Q(\zeta
% 17 )) is generated
% by the automorphism \sigma  defined by \zeta 17 ?\to  \zeta 17
%  . The sequence of subgroups
% AutQ (Q(\zeta 17 )) = ?\sigma ? \supseteq  ?\sigma  2 ? \supseteq  ?\sigma 4 ?
% \supseteq  ?\sigma  8 ? \supseteq  {e}
% corresponds via the Galois correspondence to a sequence
% Q \subseteq  Q(\delta 1 ) \subseteq  Q(\delta 2) \subseteq  Q(\delta 3 )
% \subseteq  Q(\zeta 17 ).
% It is instructive to apply what we now know to see how one may determine a
% generator \delta 1 for the first extension, that is, the fixed field of \sigma
% 2 . Let \zeta  = \zeta 17 . As
% a Q-vector space, Q(\zeta ) is generated by \zeta , \zeta  2 , . . . , \zeta  16
% . To find the fixed field of \sigma  2 ,
% write out a general element of Q(\zeta ):
% a1 \zeta  + a2 \zeta  2 + a3 \zeta  3 + a4 \zeta  4 + a5\zeta  5 + a6 \zeta  6 +
% a7 \zeta  7 + a8 \zeta  8 + a9 \zeta  9
% + a10 \zeta  10 + a11\zeta  11 + a12 \zeta  12 + a13 \zeta  13 + a14 \zeta  14 +
% a15 \zeta  15 + a 16 \zeta  16
% and apply \sigma 2 : \zeta  ?\to  (\zeta  6 )6 = \zeta  36 = \zeta  2 :
% a1 \zeta  2 + a2 \zeta  4 + a3 \zeta  6 + a4\zeta  8 + a5 \zeta  10 + a6 \zeta
% 12 + a7 \zeta  14 + a8\zeta  16 + a9 \zeta
% + a10\zeta  3 + a11\zeta  5 + a12\zeta  7 + a13 \zeta  9 + a14 \zeta  11 +
% a15\zeta  13 + a16 \zeta  15 .
% The element is fixed if and only if
% a1 = a2 = a4 = a8 = a16 = a15 = a13 = a 9 ,
% a3 = a6 = a12 = a7 = a14 = a11 = a5 = a10 ,
% and it follows that the fixed field of \sigma  2 is generated by
% \delta 1 = \zeta  + \zeta  2 + \zeta  4 + \zeta  8 + \zeta  16 + \zeta  15 + \zeta  13 + \zeta  9
% and by \zeta  3 + \zeta  6 + \zeta  12 + \zeta  7 + \zeta  14 + \zeta  11 +
% \zeta  5 + \zeta  10 = -1 - \delta 1 (remember that \zeta  is a
% root of \Phi 17 (x) = x16 + \dots  + x + 1). Pictorially,
% \zeta 
% The sum of the roots marked by black dots is \delta 1; the white roots add up to
% -1 - \delta 1 .
% The theory tells us that \delta 1 has degree 2 over Q, so \delta 1 2 must be a
% rational combination
% of 1 and \delta 1 . Indeed, pencil and paper and a bit of patience (and the
% relation \zeta  17 = 1)
% give
% \delta 1 2 = (\zeta  + \zeta  2 + \zeta  4 + \zeta  8 + \zeta  16 + \zeta  15 + \zeta  13 + \zeta  9 )2
% = 3(\zeta  + \zeta  2 + \zeta  4 + \zeta  8 + \zeta  16 + \zeta  15 + \zeta  13
% + \zeta  9)
% + 4(\zeta  3 + \zeta  6 + \zeta  12 + \zeta  7 + \zeta  14 + \zeta  11 + \zeta
% 5 + \zeta  10) + 8
% = 3\delta 1 + 4(-1 - \delta 1 ) + 8
% = -\delta 1 + 4.
% Thus, we find that \delta 1 2 + \delta 1 - 4 = 0. Solving for \delta 1 (note
% that \delta 1 > 0, as the black
% dots to the right of 0 outweigh those to the left),
% \sqrt 
% -1 + 17
% \delta 1 =
%  .
% So we could easily construct \delta 1 . The same procedure applied to the other
% extensions
% may be used to produce \delta 2 , \delta 3 , and finally \zeta 17 , whose real
% part is the complicated
% expression given at the end of \S{}3.3.

\subsection{Fundamental theorem on symmetric functions}
\label{subsec:fields:fundamental-theorem-on-symmetric-functions}
% Let t1 , . . . , tn be inde-

% terminates. The polynomial
% Pn(x) := (x - t1 ) \dots  (x - tn ) \in  Z[t1 , . . . , tn ][x]
% is universal in the sense that every polynomial of degree n with coefficients in
% (say)
% an integral domain may be obtained from Pn (x) by suitably specifying t1 , . . .
% , tn ,
% possibly in a larger ring (Exercise 7.2). The coefficients of the expansion
% P n(x) = xn - s1 (t1, . . . , tn)xn-1 + \dots  + (-1)nsn (t1, . . . , tn )
% are (up to sign) the elementary symmetric functions of t1 , . . . , tn . For
% example, for
% n = 3,
% s1(t1 , t2, t3 ) = t1 + t2 + t3 ,
% s2(t1 , t2, t3 ) = t1 t2 + t1 t3 + t2t3 ,
% s3(t1 , t2, t3 ) = t1 t2 t3.
% The functions si (t1, . . . , tn) are symmetric in the sense that they are
% invariant
% under permutations of t1 , . . . , tn (because Pn(x) is invariant); they are
% elementary
% by virtue of the following important result, which has a particularly simple
% proof
% thanks to Galois theory.
% Theorem 7.4 (Fundamental theorem on symmetric functions). Let K be a field,
% and let \phi  \in  K(t1 , . . . , tn ). Then \phi  is symmetric if and only if
% it is a rational
% function (with coefficients in K) of the elementary symmetric functions s1 , . .
% . , sn .
% Here we are viewing the si 's as elements of K[t1 , . . . , tn ], which we can
% do
% unambiguously since Z is initial in Ring.
% Proof. Let F = K(t1, . . . , tn ), and let k = K(s1, . . . , sn ) be the
% subfield generated
% by the elementary symmetric functions over K.
% Then F is a splitting field of the separable polynomial Pn (x) over k. In
% partic-
% ular k \subseteq  F is a Galois extension, and
% | Autk (F )| = [F : k] \leq  n!
% by Lemma 4.2. The symmetric group Sn acts (faithfully) on F by permuting
% t1 , . . . , tn , and this action is the identity on k = K(s1, . . . , sn )
% since each si is
% symmetric. Thus Sn may be identified with a subgroup of Autk (F ); but |Sn | =
% n!,
% so it follows that Autk (F ) = Sn.
% Since k \subseteq  F is Galois, we obtain k = F Sn . This says precisely that if
% \phi  \in
% K(t1, . . . , tn), then \phi  is invariant under the Sn action on t1, . . . , tn
% if and only if
% \phi  \in  k = K(s1, . . . , sn ), which is the statement.
%  ?
% It is not difficult to strengthen the result of Theorem 7.4 and prove that every
% symmetric polynomial is in fact a polynomial in the elementary symmetric
% functions.
% Note that the argument given in the proof amounts to the following result,
% which should be recorded for individual attention:
% Lemma 7.5. Let K be a field, t1 , . . . , tn indeterminates, and let s1, . . . ,
% sn be the
% elementary symmetric functions on t1 , . . . , tn . Then the extension
% K(s1 , . . . , sn) \subseteq  K(t1 , . . . , tn )
% is Galois, with Galois group Sn .
% This result and Cayley's theorem (Theorem II.9.5) immediately yield the fol-
% lowing:
% Corollary 7.6. Let G be a finite group. Then there exists a Galois extension k
% \subseteq  F
% such that Autk (F ) \sim 
%  = G.
% Proof. Exercise 7.4.
%  ?
% It is not known whether k may be chosen to be Q in Corollary 7.6. The proof
% hinted at above realizes k as a rather large field, and it is not at all clear
% how one
% may `descend' from this field to Q. This is the inverse Galois problem, a
% subject
% of current research. The reader can appreciate the difficulty of this problem by
% noting that even realizing Sn as a Galois group over Q (for all n) is not so
% easy,
% although it is true that for every n there are infinitely many polynomials of
% degree n
% in Z[x] whose splitting fields have Galois group Sn over Q (see Example 7.20 and
% Exercise 7.11 for the case n = p prime). If you want a memorable example to
% carry
% with you, Schur showed (1931) that the splitting field of the `truncated
% exponential'
% x2
%  x3
%  xn
% 1+ x +
%  +
%  + \cdot \cdot \cdot  +
% 2!
%  3!
%  n!
% has Galois group Sn over Q for n = 4 (and An for n = 4).
% It is known (Shafarevich, 1954) that every finite solvable group is the Galois
% group of a Galois extension over Q.
% The inverse Galois problem may be rephrased as asking whether every finite
% group is a quotient of the Galois group of Q over Q. I have already mentioned
% that this object (which may be interpreted as the `fundamental group of Q'; cf.
% Re-
% mark 6.16) is mysterious and extremely interesting28 .
% Going back to Lemma 7.5, the alternating group An has index 2 in Sn , so
% it corresponds to a quadratic extension of K(s1, . . . , sn ) via the Galois
% correspon-
% dence. It is easy to determine a generator of this extension: it must be a
% function
% of t1, . . . , t n which is invariant under the action of all even permutations
% and not
% invariant under the action of odd permutations. We have used such a function
% precisely to define even/odd permutations, in \S{}IV.4.3:
% ?
% \Delta =
%  (ti - tj ).
% 1\leq i<j\leq n
% To see that \Delta  has degree 2 over K(s1 , . . . , sn), simply note that the
% discriminant
% ?
% D = \Delta 2 =
%  (ti - tj )2
% 1\leq i<j\leq n
% is invariant under the action of all permutations, and therefore it belongs to
% the
% field K(s1, . . . , sn ) by Theorem 7.4. This proves
% Corollary 7.7. With notation as above, the extension
% \sqrt 
% K(s1 , . . . , s n)( D) \subseteq  K(t1 , . . . , tn )
% is Galois, with Galois group An .
% 28 According to J. S. Milne, the `most interesting object in mathematics' is the
% \'etale fun-
% damental group of the projective line over Q with the three points 0, 1, \infty
% removed. This is an
% extension of the Galois group of Q over Q.

\subsection{Solvability of polynomial equations by radicals}
\label{subsec:fields:solvability-of-polynomial-equations-by-radicals}
% This is possibly the

% most famous elementary application of Galois theory and one that is very close
% to
% Galois' own original motivation.
% Quadratic polynomial equations
% x2 + bx + c = 0
% are famously solved (in29 characteristic = 2) by the quadratic formula:
% \sqrt 
% -b \pm  b2 - 4ac
% x =
%  .
% Analogous, but more complicated, formulas exist for equations of degree 3 and 4
% (Tartaglia/Cardano/Ferrari, around 1540; this generated a priority dispute which
% degenerated into base personal attacks). Equations of degree 5 and higher
% resisted
% the best efforts of mathematicians for centuries. The most ambitious goal, in
% line
% with what occurs for degrees 2, 3, and 4, would be to produce a formula for the
% solutions to the `general' polynomial equation
% xn + an-1xn-2 + \dots  + a0 = 0
% in terms of the coefficients ai and basic operations such as taking roots. This
% would
% be a solution `by radicals' of the equation.
% Well, such a formula simply does not exist:
% Theorem 7.8. The general polynomial equation of degree 5 or higher admits no
% solution by radicals.
% I am being intentionally vague about the ground field in which the problem
% is posed; a `formula' such as the quadratic formula ought to hold over any field
% in which it makes sense; for example, the field should not have characteristic 2
% in the case of the quadratic formula. An easy way to avoid any such snag is
% to work in characteristic zero: this has the main technical advantage of
% ensuring
% automatic separability of polynomials xm - 1, whose roots define the `radicals';
% cf. Proposition 6.19. Therefore, we will assume in what follows that the ground
% field has characteristic 0.
% The `general polynomial' is the `universal' polynomial Pn (x) introduced in
% \S{}7.3,
% taken over a chosen ground field K: the general polynomial
% xn - s1 xn-1 + \dots  + (-1)n sn \in  k[x],
% with k = K(s1, . . . , sn ) and si elementary symmetric functions in
% (`indeterminate')
% roots t1 , . . . , tn . It is in fact no harder to deal with the more general
% case of arbitrary
% irreducible polynomials f (x) \in  k[x], over any field k (of characteristic
% zero); this is
% what we will do.
% A formula by radicals expresses a root ti of f (x) as an element of an extension
% obtained from k in the following fashion:
% k \subseteq  k(\delta 1 ) \subseteq  k(\delta 1, \delta 2 ) \subseteq  \dots
% \subseteq  k(\delta 1 , . . . , \delta r )
% where for all i = 1, . . . , r we have \delta i mi \in  k(\delta 1 , . . . ,
% \delta i-1) for some exponent mi .
% 29 In characteristic 2, one can give a formula in terms of `2-roots', which are
% roots of the
% Artin-Schreier equation, and avoid the problematic division by 2; see Exercise
% 7.16.
% Definition 7.9. Extensions of this type are called radical extensions.
%  ?
% Some facts about radical extensions are immediate from the definition: for ex-
% ample, if k \subseteq  F and k \subseteq  K are two radical extensions and F , K
% are subfields
% of a larger field, then k \subseteq  F K is radical (Exercise 7.5). Further,
% radical exten-
% sions are composite of cyclic ones provided the appropriate roots of 1 are in k,
% by
% Proposition 6.19. Also,
% Lemma 7.10. Every separable radical extension is contained in a Galois radical
% extension.
% Proof. Let k \subseteq  F be a separable radical extension. In particular k
% \subseteq  F is finite
% and separable, so F = k(\alpha ) for some \alpha  \in  F (Proposition 5.19). Let
% p(x) be the
% minimal polynomial of \alpha  over k. The splitting field L of p(x) over k may
% be obtained
% as the composite of \sigma (F ) as \sigma  ranges over the embeddings of F in k,
% extending idk :
% indeed, this composite is generated over k by all \sigma (\alpha ), which range
% over all roots
% of p(x). As a composite of radical extensions, L is radical over k. As a
% splitting
% field of a separable polynomial, L is Galois over k (Theorem 6.9, part (1)), and
% we
% are done.
%  ?
% Galois theory will allow us to relate solvability by radicals of a polynomial f
% (x)
% with a precise statement on the splitting field of f (x). Here is a first
% formalization
% of this connection:
% Lemma 7.11. Let k be a field of characteristic 0, and let f (x) \in  k[x] be an
% irre-
% ducible polynomial. Then there exists a formula solving f (x) by radicals if and
% only
% if the splitting field of f (x) is contained in a Galois radical extension.
% Proof. If the splitting field of f (x) is contained in a radical extension
% (Galois or
% not), then the roots may be written as combinations of field operations and
% radicals,
% as needed.
% For the converse, assume f (x) is solvable by radicals. A formula for a root t1
% can be turned into a formula for any other root ti by applying an element \sigma
% i of
% (for example) Autk (k) sending t1 to ti (such an automorphism exists since f (x)
% is
% irreducible). Thus, if a formula exists for one root, then a formula exists for
% all
% roots. The composite of all the corresponding radical extensions gives one,
% possibly
% larger, radical extension of k containing all roots of f (x). Therefore, the
% splitting
% field of f (x) is contained in a radical extension of k. By Lemma 7.10, this
% extension
% may be assumed to be Galois, as needed.
%  ?
% In a nutshell, we will understand solvability of polynomial equations by
% radicals
% if we understand radical Galois extensions. As demanded by the general
% philosophy
% underlying Galois theory, we should aim to characterize these extensions in
% terms
% of their Galois groups. It turns out that the right condition (modulo technical
% considerations) is that the Galois group should be solvable.
% Definition 7.12. A Galois extension k \subseteq  F is solvable if Autk (F ) is a
% solvable
% group.
%  ?
% The reader may want to revisit \S{}IV.3.3 now, for a reminder on solvability in
% group theory. (We are approaching an explanation for the `solvable'
% terminology!)
% Just as radical extensions decompose as sequence of cyclic extensions, solvable
% groups have cyclic composition factors (Proposition IV.3.11). I will say30 that
% a radical, resp., solvable, extension k \subseteq  F `has enough roots of 1' if
% k contains a
% primitive M -th root of 1 for a common multiple M of the order of the
% corresponding
% cyclic factors.
% Lemma 7.13. Let k \subseteq  F be a Galois extension, with char k = 0. Provided
% it has
% enough roots of 1, k \subseteq  F is radical if and only if it is solvable.
% Proof. Modulo the fundamental theorem of Galois theory, this is a
% straightforward
% generalization of Proposition 6.19.
% Indeed, assume that k \subseteq  F is radical:
% k \subseteq  k(\delta 1 ) \subseteq  \dots  \subseteq  k(\delta 1 , . . . ,
% \delta r ) = F
% with \delta i mi \in  k(\delta 1, . . . , \delta i-1 ). By hypothesis, k
% contains primitive mi-th roots of 1 for
% all i. Consider the corresponding series of subgroups of G = Autk (F (\zeta )):
% G = G0 \supseteq  G 1 \supseteq  \dots  \supseteq  Gr = {e},
% with Gi = Autk(\delta  1 ,...,\delta i ) (F ). Note that F is Galois over each
% intermediate field.
% Each extension
% k(\delta 1 , . . . , \delta i-1 ) \subseteq  k(\delta 1 , . . . , \delta i )
% satisfies the hypothesis of Proposition 6.19, and it follows that it is Galois
% and
% cyclic. By the fundamental theorem (Theorem 6.15) each Gi is normal in the
% previous Gi-1 , with cyclic quotient, and it follows that G is solvable
% (Proposi-
% tion IV.3.11 (ii)).
% The same argument in reverse, using the converse implication in Proposi-
% tion 6.19, proves that solvable extensions with enough roots of 1 are radical.
%  ?
% How do we account for the possible absence of the needed roots of 1? Surpris-
% ingly, this affects the statement in a rather minor way.
% Proposition 7.14. Let k be a field of characteristic 0, and let k \subseteq  F
% be a Galois
% extension. Then k \subseteq  F is solvable if and only if it is contained in a
% Galois radical
% extension.
% The characteristic 0 in Proposition 7.14 is an overshoot; as in Proposition
% 6.19,
% it would suffice to require that char k does not divide the relevant degrees.
% Proof. Assume Autk (F ) is solvable. Let M be a common multiple of the order
% of the cyclic quotients, and let \zeta  be a primitive M -th root of 1 in an
% algebraic
% closure k of k. The splitting field k(\zeta ) of xM - 1 over k is Galois over k
% (xM - 1 is
% separable since k has characteristic 0). By Proposition 6.17, the extension
% k(\zeta ) \subseteq
% F (\zeta ) is also Galois, with Galois group isomorphic to Autk(\zeta )\cap F (F
% ). It follows that
% Autk(\zeta )(F (\zeta )) is solvable, since Autk(\zeta )\cap F (F ) is a
% subgroup of Autk F and the latter
% is solvable by assumption. (See the comment following Corollary IV.3.13.)
% 30Warning! This is nonstandard terminology.
% Since k(\zeta ) \subseteq  F (\zeta ) has enough roots of 1, Lemma 7.13 implies
% that it is rad-
% ical. So is k \subseteq  F (\zeta ), since k \subseteq  k(\zeta ) is itself
% trivially radical. This proves that
% k \subseteq  F is contained in a radical extension, hence in a Galois radical
% extension by
% Lemma 7.10.
% Conversely, assume that k \subseteq  F is Galois and contained in a radical
% Galois
% extension k \subseteq  L. Let M be a common multiple of the order of cyclic
% factors in this
% extension, and let \zeta  be a primitive M -th root of 1. The composite L(\zeta
% ) is Galois
% and radical over k(\zeta ) (Proposition 6.17 and Exercise 7.5); further, it has
% enough
% roots of 1.
% By Lemma 7.13, k(\zeta ) \subseteq  L(\zeta ) is solvable. It follows that k
% \subseteq  L(\zeta ) is solvable.
% Indeed, this extension is Galois (Exercise 6.13); apply the fundamental theorem
% to
% k \subseteq  k(\zeta ) \subseteq  L(\zeta )
% to establish that Autk (k(\zeta )) is isomorphic to the quotient of Autk
% (L(\zeta )) by its nor-
% mal subgroup Autk(\zeta ) (L(\zeta )). Both Autk(\zeta )(L(\zeta )) (as shown
% above) and Autk (k(\zeta ))
% (an abelian group; cf. Example 6.18) are solvable, and it follows that Autk
% (L(\zeta ))
% is solvable as promised, by Corollary IV.3.13.
% Since F is an intermediate field and Galois over k, by the fundamental theorem
% Autk (F ) is a homomorphic image of Autk (L(\zeta )). This shows that Autk (F )
% is a
% homomorphic image of a solvable group; hence it is itself solvable (Corollary
% IV.3.13
% again), and we are done.
%  ?
% We may have lost sight of where we were heading; but we are ready to reconnect
% with the study of solutions to polynomial equations.
% Definition 7.15. The Galois group Galk (f (x)) of a separable polynomial f (x)
% \in
% k[x] is the Galois group of the splitting field of f (x) over k.
%  ?
% Corollary 7.16. Let k be a field of characteristic 0, and let f (x) \in  k[x] be
% an
% irreducible polynomial. Then f (x) is solvable by radicals if and only if its
% Galois
% group is solvable.
% Proof. This is an immediate consequence of Lemma 7.11 and Proposition 7.14.
%  ?
% Corollary 7.16 is called Galois' criterion. Ruffini (1799) and Abel (1824) had
% previously established that general formulas in radicals for the solutions of
% equa-
% tions of degree \geq  5 do not exist (that is, Theorem 7.8); but it was Galois
% who
% identified the precise condition given in Corollary 7.16. Of course, Theorem 7.8
% follows immediately from Galois' criterion:
% Proof of Theorem 7.8. By Lemma 7.5, the Galois group of the general polyno-
% mial of degree n is Sn . The group Sn is not solvable for n \geq  5 (Corollary
% IV.4.21);
% hence the statement follows from Corollary 7.16.
%  ?
% In fact, we could now do more: we know that S3 and S4 are solvable (cf. Exer-
% cise IV.3.16); from a composition series with cyclic quotients we could in
% principle
% decompose explicitly the splitting field of general polynomials of degree 3 and
% 4
% as radical extensions and as a consequence recover the Tartaglia/Cardano/Ferrari
% formulas for their solutions.

\subsection{Galois groups of polynomials}
\label{subsec:fields:galois-groups-of-polynomials}
% Impressive as Galois' criterion is, it does

% not in itself produce a single polynomial of degree (say) 5 over Q that is not
% solvable
% by radicals. Finding such polynomials requires a certain amount of extra work;
% we
% will be able to do this by the end of this subsection.
% Computing Galois groups of polynomials is in fact a popular sport among
% algebraists, excelling at which would require much more information than the
% reader
% will glean here. I will just list a few straightforward observations.
% To begin with, recall that an element of Autk (F ) must send roots of a polyno-
% mial f (x) \in  k[x] to roots of the same polynomial, and if f (x) is
% irreducible and F
% is its splitting field, then there are automorphisms of F sending any root of f
% (x) to
% any other root (Proposition 1.5, Lemma 4.2). This observation may be rephrased
% as follows:
% Lemma 7.17. Let f (x) \in  k[x] be a separable irreducible polynomial of degree
% n.
% Then Galk (f (x)) acts transitively on the set of roots of f (x) in k. In
% particular,
% Galk (f (x)) may be identified with a transitive subgroup of the symmetric group
% Sn .
% Of course a subgroup of Sn is transitive if the corresponding action on {1, . .
% . , n}
% is transitive; the diligent reader has run across this terminology in Exercise
% IV.4.12.
% The reader will easily produce a statement analogous to Lemma 7.17 in case
% f (x) is reducible. Of course the action is not transitive in this case, and if
% the
% factors of f (x) have degrees n1 , . . . , nr , then Gal k (f (x)) is contained
% in a subgroup
% of Sn isomorphic to Sn1 \times  \dots  \times  Snr .
% Lemma 7.17 (or its `reducible' variations) already give some information. For
% example, we see that the Galois group of a separable irreducible cubic can only
% be
% A3 or S3 , since these are the only transitive subgroups of S3 . Similarly, the
% range
% of possibilities for irreducible polynomials of degree 4 is rather restricted:
% S4 , A4 ,
% and isomorphic copies of the dihedral group D8 , of Z/2Z \times  Z/2Z, and of
% Z/4Z
% (Exercise 7.8). One approach to the computation of Galois groups of polynomials
% amounts to defining invariants imposing further restrictions, leading to
% algorithms
% deciding which of such possibilities occurs for a given polynomial.
% We have already encountered the most important of these invariants: if the
% roots of a separable polynomial f (x) \in  k[x] in its splitting field are
% \alpha 1 , . . . , \alpha n , the
% discriminant of f (x) is the element D = \Delta 2 , where
% ?
% \Delta =
%  (\alpha i - \alpha j ).
% 1\leq i<j\leq n
% Every permutation of the roots fixes D, so D must be fixed by the whole Galois
% group; therefore, D \in  k. Odd permutations move \Delta  (if char k = 2), and
% even
% permutations fix it; therefore \Delta  is fixed by the Galois group G (in other
% words,
% \Delta  \in  k) if and only if G \subseteq  An . This proves
% Lemma 7.18. Let k be a field of characteristic = 2, and let f (x) \in  k[x] be a
% sepa-
% rable polynomial, with discriminant D. Then the Galois group of f (x) is
% contained
% in the alternating group An if and only if D is a square in k.
% Example 7.19. Lemma 7.18 and a discriminant computation are all that is needed
% to compute the Galois group of an irreducible cubic polynomial
% f (x) = x3 + ax2 + bx + c.
% It would be futile to try to remember the discriminant
% D = a2 b2 - 4a3 c - 4b3 + 18abc - 27c2 ;
% but one may remember the trick of shifting x by a/3 (in characteristic = 3),
% with
% the effect of killing the coefficient of x2 :
% a
% f (x - ) = x3 + px + q
% for suitable p and q. This does not change D (shifting all roots \alpha i by the
% same
% amount has no effect on the differences \alpha i - \alpha j ), yet
% D = -4p3 - 27q 2
% is a little more memorable.
% By Lemma 7.18 and the preceding considerations, the Galois group is A3 if D
% is a square and S3 otherwise. For example, x3 - 2 has Galois group S3 over Q,
% since
% D = -108 is not a square in Q; x3 - 3x + 1 has discriminant 81 = 92, and
% therefore
% it has Galois group A3 \sim 
%  = Z/3Z.
%  ?
% The reader will have no difficulty locating a discussion of the different pos-
% sibilities for polynomials of degree 4 and detailed information for higher
% degree
% polynomials. I will just highlight the following simple observation.
% Example 7.20. Let f (x) \in  Q[x] be an irreducible polynomial of degree p,
% where
% p is prime. Assume that f (x) has p - 2 real roots and 2 nonreal, complex roots.
% Then the Galois group of f (x) is S p .
% Indeed, complex conjugation induces an automorphism of the splitting field and
% acts by interchanging the two nonreal roots, so the Galois group G, as a
% subgroup
% of Sp , contains a transposition. On the other hand, the degree of the splitting
% field
% (and hence |G|) is divisible by p, because it contains a simple extension of
% order p,
% obtained by adjoining any one root to Q. Since p is prime, G contains an element
% of
% order p by Cauchy's theorem (Theorem IV.2.1); the only elements of order p in Sp
% are p-cycles, so G contains a p-cycle. It follows that G = Sp , by (a simple
% variation
% of) Exercise IV.4.7.
% For example, the Galois group of f (x) = x5 - 5x - 1 over Q is S5, giving a
% concrete example of a quintic that cannot be solved by radicals (Exercise 7.10).
% The reader will use this technique to produce polynomials of every prime de-
% gree p in Z[x] whose Galois group is Sp (Exercise 7.11).
%  ?

\subsection{Abelian groups as Galois groups over Q}
\label{subsec:fields:abelian-groups-as-galois-groups-over-q}
% Having mentioned the inverse

% Galois problem in \S{}7.3, I should point out that the reader is now in the
% position of
% proving that every finite abelian group may be realized as the Galois group of
% an
% extension over Q.
% In fact, the reader can prove a much more precise result: every finite abelian
% group may be realized as the group of some intermediate field of the extension
% Q \subseteq  Q(\zeta n) of Q in a cyclotomic field. This uses a number-theoretic
% fact:
% For every integer N , there are infinitely many primes p such that p \equiv  1
% mod N .
% This is a particular case of Dirichlet's theorem (1837), which states that if a,
% b are
% positive integers and gcd(a, b) = 1, then there are infinitely many primes of
% the form
% a + nb with n > 0. The particular case a = 1, b = N needed here was apparently
% already known to Euler and can be proven by elementary means (in fact, there is
% a
% proof using cyclotomic polynomials: see Exercise 5.18). Assuming this fact,
% argue
% as follows:
% ---By the classification theorem (Theorem IV.6.6), every finite abelian group G
% is isomorphic to a product of cyclic groups
% (Z/n1 Z) \times  \dots  \times  (Z/nr Z).
% ---By Dirichlet's theorem, we can then choose distinct primes pi such that pi
% \equiv  1
% mod ni .
% ---Let n = p1 \dots  pr ; by the Chinese Remainder Theorem (Theorem V.6.1)
% (Z/nZ) \ast  \sim 
%  = (Z/p1Z)\ast  \times  \dots  \times  (Z/pr Z)\ast  ,
% and the i-th factor on the right-hand side is cyclic of order pi - 1, a multiple
% of ni .
% ---It follows that (Z/nZ)\ast  has a subgroup H such that G \sim 
%  = (Z/nZ)\ast  /H.
% ---Since (Z/nZ)\ast  \sim 
%  = AutQ(Q(\zeta n )) (Proposition 5.16) and cyclotomic fields are
% Galois over Q, H corresponds to an intermediate field F , Q \subseteq  F
% \subseteq  Q(\zeta n).
% ---Since AutQ (Q(\zeta n)) is abelian, H is automatically normal; hence, Q
% \subseteq  F is
% Galois, and AutQ (F ) \sim 
%  = G, as needed.
% Thus, the `inverse Galois problem' has an easy solution for abelian groups.
% A much stronger result is true: it can be proven that every finite abelian
% extension of Q is contained in some cyclotomic field. This is the
% Kronecker-Weber
% theorem, dating back to the second half of the nineteenth century. The natural
% context for this result is class field theory and is well beyond the scope of
% this book.
% The reader will be able to verify (Exercise 7.19) that every quadratic extension
% of Q
% is contained in some cyclotomic field.

\subsection*{Exercises}

% 7.1. Find explicitly a generator of a quadratic intermediate field of the
% extension
% Q \subseteq  Q(\zeta 10 ).
% 7.2. ? Let R be an integral domain, and let f (x) \in  R[x] be a polynomial of
% degree n.
% Show how to obtain f (x) by specializing the `universal' polynomial Pn (x)
% defined
% in \S{}7.3. [\S{}7.3]
% 7.3. Prove that the elementary symmetric functions s1, . . . , sn (see \S{}7.3)
% are alge-
% braically independent. (Hint: Use Exercise 1.28.)
% 7.4. ? Prove that every finite group is isomorphic to the group of automorphisms
% of some Galois extension. [\S{}7.3]
% 7.5. ? Let k \subseteq  F be a radical extension, and let k \subseteq  K be any
% extension; assume
% F and K are contained in a larger field. Prove that K \subseteq  F K is radical.
% [\S{}7.4]
% 7.6. Let f (x) \in  Q[x] be an irreducible cubic, and
%  let \rho  be a real root of f (x). Prove
% \sqrt that the splitting field of f (x) over Q is Q(\rho , D), where D is the discriminant
% of f (x).
% 7.7. Let f (x) \in  Q[x] be an irreducible cubic with exactly one real root.
% Prove that
% the discriminant of f (x) is not a square in Q.
% 7.8. ? (Cf. Exercise IV.4.12.) Find (mathematically or by library search) the
% lattice
% of subgroups of S4 , and verify that the transitive subgroups of S4 are
% (isomorphic
% to) S4, A4 , D8 , Z/2Z \times  Z/2Z, and Z/4Z. [\S{}7.5]
% 7.9. Compute the Galois group of the polynomial x4 - 2.
% 7.10. ? Prove that the polynomial x5 - 5x - 1 has exactly 3 real roots (this is
% a
% calculus exercise!) and is irreducible over Q. Prove that its Galois group is S5
% .
% [\S{}7.5]
% 7.11. ? Let n > 0 be an integer. Note that
% fn (x) := (x2 + 2) \cdot  x \cdot  (x - 2) \dots  (x - 2(n - 4)) \cdot  (x - 2(n
% - 3))
% has n - 2 rational roots and 2 nonreal, complex roots. Prove that for infinitely
% many integers q, the polynomial
% qfn (x) + 2 \in  Z[x]
% has (n - 2) real roots, 2 nonreal complex roots, and is irreducible over Q.
% Conclude
% that for each prime p there are infinitely many polynomials of degree p in Z[x]
% whose Galois group is Sp . [\S{}7.3, \S{}7.5]
% 7.12. \lnot  Let f (x) \in  k[x] be a separable irreducible polynomial of prime
% degree p
% over a field k, and let \alpha 1 , . . . , \alpha p be the roots of f (x) in its
% splitting field F . Prove
% that the Galois group of f (x) contains an element \sigma  of order p, `cycling'
% through
% the roots. [7.13]
% 7.13. \lnot  Let f (x) \in  k[x] be a separable irreducible polynomial of prime
% degree p
% over a field k. Let \alpha  be a root of f (x) in k, and suppose you can express
% another root
% of f (x) as a polynomial in \alpha , with coefficients in k. Prove that you can
% express
% all roots as polynomials in \alpha  and that the Galois group of f (x) is Z/pZ.
% (Use
% Exercise 7.12.) [7.14]
% 7.14. \lnot  Let k be a field of characteristic p > 0, and let f (x) = xp - x -
% a \in  k[x].
% Assume that f (x) has no roots in k. Prove that f (x) is irreducible and that
% its
% Galois group is Z/pZ. (Note that if \alpha  is a root of f (x), so is \alpha  +
% 1; use Exercises 1.8
% and 7.13.)
% The splitting field of f (x) is an Artin-Schreier extension of Fp . [7.16]
% 7.15. \lnot  Let k \subseteq  F be a cyclic Galois extension of degree p, where
% char k = p > 0.
% Prove that it is an Artin-Schreier extension. (Hint: What is tr k\subseteq F
% (1)? Use the
% additive version of Hilbert's theorem 90, Exercise 6.19.) [6.19]
% 7.16. ? The Artin-Schreier map on a field k of positive characteristic p is the
% function
% x ?\to  xp - x.
% AS \sqrt 
% Denote this function by AS, and let \cdot  denote its inverse (which is only
% defined
% up to the addition of an element of Fp : see Exercise 7.14).
% Express the
%  solutions of a quadratic equation x2
%  + bx + c = 0 in characteristic 2
% AS \sqrt in terms of \cdot , for b = 0. (For b = 0, the roots must live in an
% Artin-Schreier
% extension of k, since the polynomial is then separable. It is inseparable for b
% = 0;
% solving the equation in this case amounts to extending k by the unique square
% root
% of c.) [\S{}7.4]
% 7.17. \lnot  Let f (x) \in  k[x] be a separable irreducible polynomial of degree
% n over a
% field k, and let F be its splitting field. Assume Autk (F ) = \sim 
%  Sn , and let \alpha  be a root
% of f (x) in F .
% \bullet  Prove that Autk(\alpha ) (F ) \sim 
%  = Sn-1 .
% \bullet  Prove that there are no proper subfields of k(\alpha ) properly containing k. (Hint:
% Exercise IV.4.8.)
% [7.18]
% 7.18. Let f (x) \in  k[x] be a separable irreducible polynomial of degree n over
% a
% field k, with Galois group Sn . Let \alpha  be a root of f (x) in k, and let
% g(x) \in  k[x] be
% any nonconstant polynomial of degree < n. Prove that \alpha  may be expressed as
% a
% polynomial in g(\alpha ), with coefficients in k. (Use Exercise 7.17.)
% 7.19. ? Let d be a positive integer.
% ?d
% \bullet  Prove that the discriminant of a polynomial f (x) =
%  i=1 (x - \alpha i ) equals the
% ?
% product \pm  di=1 f \setminus  (\alpha i ).
% \bullet  Prove that the discriminant of xd - 1 is \pm dd , and note that this is a square
% in Q(\zeta d ).
% \bullet  Prove that Q(\zeta 8d ) contains square roots of d and -d.
% \bullet  Conclude that every quadratic extension of Q may be embedded in a cyclotomic
% field Q(\zeta n ), for a suitable n.
% (This is a very special case of the Kronecker-Weber theorem.) [\S{}7.6]

%%% Local Variables:
%%% mode: LaTeX
%%% TeX-engine: xetex
%%% TeX-master: "../master"
%%% End:
